\begin{thebibliography}{999}
\setlength{\itemsep}{0pt}
\small
\raggedright

% ============ SOLAR ============
\bibitem{owid-learning} Our World in Data, ``Learning curves: solar module prices fell $\sim$20\% per doubling of cumulative capacity; \$106/W (1976) to \$0.38/W (2019).'' \url{https://ourworldindata.org/learning-curve}
\bibitem{pv-shares-2023} A.~J\"ager-Waldau / ISES, ``PV Industry Trends'' (Oct.\ 2024): 2023 China production shares---polysilicon 92\%, wafers 98.1\%, cells 92\%, modules 84.6\%. \url{https://www.ises.org/sites/default/files/webinars/2024/October/Oct242024_PVIndustryTrends_final.pdf}
\bibitem{qcells-sharp} EDN, ``Q-Cells tops Sharp in solar rankings'' (2008). \url{https://www.edn.com/q-cells-tops-sharp-in-solar-rankings/}
\bibitem{german-fit} Wikipedia, ``Feed-in tariffs in Germany'' (EEG 2004: 57.4 ct/kWh rooftop, 20-year guarantee). \url{https://en.wikipedia.org/wiki/Feed-in_tariffs_in_Germany}
\bibitem{suntech} Wikipedia, ``Suntech Power'' (NYSE listing Dec.\ 2005; 2 GW capacity 2011; \$541M default and Wuxi bankruptcy, Mar.\ 2013). \url{https://en.wikipedia.org/wiki/Suntech_Power}
\bibitem{cdb-credit} Eco-Business/Bloomberg, ``Chinese renewable companies slow to tap \$47 billion credit'' (CDB lines, 2010). \url{https://www.eco-business.com/news/chinese-renewable-companies-slow-to-tap-47-billion-credit/}
\bibitem{sei-solar} 12th Five-Year Plan for the Solar Photovoltaic Industry (2011), official translation. \url{https://policy.asiapacificenergy.org/sites/default/files/chinas-five-year-plan-for-solar-translation.pdf}
\bibitem{bernreuter} Bernreuter Research, ``Polysilicon price trend'' (\$475/kg Mar.\ 2008 $\rightarrow$ \$15.9/kg Dec.\ 2012 $\rightarrow$ \$5.6/kg Jul.\ 2024). \url{https://www.bernreuter.com/polysilicon/price-trend/}
\bibitem{solyndra} Congressional Research Service, ``Solyndra'' (R42058); \$535M DOE guarantee, 2011 bankruptcy. \url{https://www.congress.gov/crs-product/R42058}
\bibitem{qcells-insolvency} NS Energy, ``Hanwha completes Q-Cells takeover'' (insolvency Apr.\ 2012). \url{https://www.nsenergybusiness.com/news/newshanwha-completes-q-cells-takeover/}
\bibitem{solarworld} PV Tech, ``SolarWorld entering insolvency proceedings'' (May 2017). \url{https://www.pv-tech.org/solarworld-entering-insolvency-proceedings/}
\bibitem{us-adcvd} US Dept.\ of Commerce, AD/CVD final determinations on PRC solar cells (Oct.\ 10, 2012). \url{https://enforcement.trade.gov/download/factsheets/factsheet_prc-solar-cells-ad-cvd-finals-20121010.pdf}
\bibitem{eu-solar} Clifford Chance, ``EU anti-dumping duties imposed against Chinese solar panel producers'' (2013); Al Jazeera, ``China and EU reach deal on solar panel trade'' (Jul.\ 28, 2013). \url{https://www.aljazeera.com/news/2013/7/28/china-and-eu-reach-deal-on-solar-panel-trade}
\bibitem{s201} Utility Dive, ``Trump issues 30\% tariff on solar panel imports'' (Jan.\ 2018). \url{https://www.utilitydive.com/news/trump-issues-30-tariff-on-solar-panel-imports/515265/}
\bibitem{policy531} Climate Home News, ``China solar industry struggles with sudden subsidy cuts'' (531 policy, 2018). \url{https://www.climatechangenews.com/2018/08/15/china-solar-industry-struggles-sudden-subsidy-cuts/}
\bibitem{toprunner} SolarQuotes, ``China's Top Runner Program''; PV Tech on PERC dominance of Top Runner rounds. \url{https://www.solarquotes.com.au/blog/chinas-top-runner-program/}
\bibitem{china-2024-installs} pv magazine, ``China hits 277.17 GW of new PV installations in 2024.'' \url{https://www.pv-magazine.com/2025/01/21/china-hits-277-17-gw-of-new-pv-installations-in-2024/}
\bibitem{china-2025-installs} TaiyangNews (NEA data), ``China installed over 315 GW AC new solar PV capacity in 2025''; cumulative 1.2 TW AC. \url{https://taiyangnews.info/markets/china-installed-over-315-gw-ac-new-solar-pv-capacity-in-2025}
\bibitem{iea-supply} IEA, \emph{Special Report on Solar PV Global Supply Chains} (2022). \url{https://www.iea.org/reports/solar-pv-global-supply-chains/executive-summary}
\bibitem{poly-2024} PV Tech, ``Chinese firms nine of top ten polysilicon manufacturers 2024'' (93.5\% share; sub-\$4.50/kg prices). \url{https://www.pv-tech.org/chinese-firms-nine-top-ten-polysilicon-manufacturers-2024/}
\bibitem{module-capacity} PV Tech, ``Global solar module manufacturing capacity to reach 1.8 TW in 2025.'' \url{https://www.pv-tech.org/global-solar-module-manufacturing-capacity-to-reach-1-8tw-in-2025-report/}
\bibitem{infolink-2025} InfoLink, ``2025 global module shipment ranking: combined shipments reach 536 GW'' (exports 267.6 GW; TOPCon $\sim$95\%). \url{https://www.infolink-group.com/energy-article/solar-topic-infolink-2025-global-module-shipment-ranking-combined-shipments-reach-536-gw}
\bibitem{ember-cells} Ember, ``China solar cell exports grow 73\% in 2025.'' \url{https://ember-energy.org/latest-insights/china-solar-cell-exports-grow-73-in-2025/}
\bibitem{energytrend-prices} EnergyTrend solar price index (Jul.\ 29, 2026): TOPCon modules RMB 0.67--0.72/W ($\approx$\$0.09--0.10/W). \url{https://www.energytrend.com/solar-price.html}
\bibitem{pv-losses-2025} PV Tech, ``Five Chinese PV giants anticipate combined 2025 losses of US\$4.1--4.7 billion.'' \url{https://www.pv-tech.org/five-chinese-pv-giants-anticipate-combined-2025-losses-of-us4-1-4-7-billion/}
\bibitem{anti-involution} SCMP, ``China's top leadership takes aim at disorderly low-price competition'' (Jul.\ 2025); Xinhua commentary. \url{https://www.scmp.com/economy/china-economy/article/3316616/chinas-top-leadership-takes-aim-disorderly-low-price-competition}
\bibitem{poly-fund} pv magazine, ``China's \$7 billion polysilicon consolidation plan'' (Aug.\ 2025). \url{https://www.pv-magazine.com/2025/08/14/chinas-7-billion-consolidation-plan-may-be-reflected-in-prices-in-short-time/}
\bibitem{poly-platform} pv magazine, ``Beijing group creates official body to guide China's polysilicon capacity'' (Dec.\ 2025). \url{https://www.pv-magazine.com/2025/12/11/beijing-group-creates-official-body-to-guide-chinas-polysilicon-capacity/}
\bibitem{longi-record} pv magazine, ``LONGi sets new world record with 35.5\%-efficient perovskite-silicon tandem cell'' (Jul.\ 2026). \url{https://www.pv-magazine.com/2026/07/15/longi-sets-new-world-record-with-35-5-efficient-perovskite-silicon-tandem-cell/}
\bibitem{trina-patents} SolarQuarter, ``Trina Solar tops global perovskite solar cell patent ranking'' (May 2025). \url{https://solarquarter.com/2025/05/12/trinasolar-tops-global-perovskite-solar-cell-patent-ranking-reinforces-leadership-in-next-gen-pv-technologies/}
\bibitem{ira-build} pv magazine USA, ``U.S. solar manufacturing forges on'' (module capacity 8 $\rightarrow$ 56.5 GW; cells lagging; Aug.\ 2025). \url{https://pv-magazine-usa.com/2025/08/19/u-s-solar-manufacturing-forges-on/}
\bibitem{obbba} pv magazine, ``US solar faces steep cuts as OBBBA phases out key tax credits'' (Sep.\ 2025). \url{https://www.pv-magazine.com/2025/09/16/us-solar-faces-steep-cuts-as-obbba-phases-out-key-tax-credits/}
\bibitem{firstsolar} PV Tech, ``First Solar net sales, gross profit tick upwards 2025'' (16.1 GW produced). \url{https://www.pv-tech.org/first-solar-net-sales-gross-profit-tick-upwards-2025/}

% ============ BATTERIES ============
\bibitem{sony-1991} Battery Design, ``Sony lithium-ion batteries: world's first commercialized LIB (1991).'' \url{https://www.batterydesign.net/sony-lithium-ion-batteries-worlds-first-commercialized-lib-1991/}
\bibitem{byd-history} Wikipedia, ``BYD Company'' (founded 1995; world \#1 NiCd by 2002; Qinchuan acquisition 2003). \url{https://en.wikipedia.org/wiki/BYD_Company}
\bibitem{catl-history} Wikipedia, ``CATL'' (ATL 1999; CATL spin-off 2011; world \#1 from 2017). \url{https://en.wikipedia.org/wiki/CATL}
\bibitem{white-list} CGTN, ``China scraps white list of domestic auto battery suppliers'' (Jun.\ 2019). \url{https://news.cgtn.com/news/2019-06-25/China-scraps-white-list-of-domestic-auto-battery-suppliers-HO26IZLsWI/index.html}
\bibitem{lfp-story} H.~Sanderson, ``The battery bet that changed the world'' (Volt Rush Substack): LFP subsidy politics, Blade/CTP, Tesla-CATL 2020. \url{https://voltrush.substack.com/p/the-battery-bet-that-changed-the}
\bibitem{sne-2025} SNE Research via electrive, ``CATL continues to dominate global battery market'' (FY2025: 1,187 GWh; CATL 39.2\%, BYD 16.4\%). \url{https://www.electrive.com/2026/02/04/sne-research-catl-continues-to-dominate-global-battery-market/}
\bibitem{iea-gevo-2026} IEA, \emph{Global EV Outlook 2026}, ``Electric vehicle batteries'' chapter. \url{https://www.iea.org/reports/global-ev-outlook-2026/electric-vehicle-batteries}
\bibitem{eia-components} US EIA, ``China's battery component production shares'' (anodes 85\%, electrolytes 82\%, separators 74\%, cathodes 70\%). \url{https://www.eia.gov/todayinenergy/detail.php?id=65305}
\bibitem{benchmark-graphite} Benchmark Mineral Intelligence, ``China controls three quarters of graphite anode supply chain.'' \url{https://source.benchmarkminerals.com/article/infographic-china-controls-three-quarters-of-graphite-anode-supply-chain}
\bibitem{nickel} Mining Technology, ``China's control of Indonesia's nickel capacity'' ($\sim$75\% of refining). \url{https://www.mining-technology.com/news/chinas-indonesias-nickel-capacity/}
\bibitem{bnef-2025} BloombergNEF, ``Lithium-ion battery pack prices fall to \$108/kWh'' (2025 survey; China \$84/kWh). \url{https://about.bnef.com/insights/clean-transport/lithium-ion-battery-pack-prices-fall-to-108-per-kilowatt-hour-despite-rising-metal-prices-bloombergnef/}
\bibitem{ford-catl} Ford Authority, ``Ford EV battery plant will begin production this summer'' (CATL-licensed LFP, Marshall MI, 2026). \url{https://fordauthority.com/2026/02/ford-ev-battery-plant-will-begin-production-this-summer/}
\bibitem{naxtra} CarNewsChina, ``CATL confirms 2026 large-scale sodium-ion battery deployment'' (Naxtra: 175 Wh/kg, mass production Dec.\ 2025). \url{https://carnewschina.com/2025/12/28/catl-confirms-2026-large-scale-sodium-ion-battery-deployment-in-multiple-sectors/}
\bibitem{shenxing} CNBC, ``CATL claims to beat BYD's EV battery record with longer range on a 5-minute charge'' (2nd-gen Shenxing, Apr.\ 2025). \url{https://www.cnbc.com/2025/04/22/chinas-catl-claims-to-beat-byds-ev-battery-record-with-longer-range-on-a-5-minute-charge.html}
\bibitem{byd-1mw} BYD, ``Super e-Platform megawatt flash charging'' (Mar.\ 2025). \url{https://www.byd.com/en/news-list/BYD-Unveils-Super-e-Platform-Megawatt-Flash-Charging-Electric-Vehicles-Matching-Refueling-Speeds.html}
\bibitem{mofcom-battery-controls} HSF Kramer, ``China battery technology export restrictions'' (Jul.\ 2025); Global Times on Oct.\ 2025 extension to cells $\geq$300 Wh/kg, anodes, equipment. \url{https://www.hsfkramer.com/notes/mining/2025-posts/china-battery-technology-export-restrictions-cathode-lithium-gallium}
\bibitem{northvolt} Wikipedia, ``Northvolt'' (Ch.\ 11 Nov.\ 2024; Swedish bankruptcy Mar.\ 2025; \$15B raised; 79.8 MWh produced Q1--Q3 2023). \url{https://en.wikipedia.org/wiki/Northvolt}
\bibitem{a123} C\&EN, ``A123 Systems goes bankrupt'' (2012; Wanxiang acquisition). \url{https://cen.acs.org/articles/90/i43/A123-Systems-Goes-Bankrupt.html}
\bibitem{us-cancellations} Energy-Storage.News/CEA, ``US battery gigafactories face delays and cancellations'' (46 GWh cancelled by mid-2025). \url{https://www.energy-storage.news/us-battery-gigafactories-face-delays-and-cancellations-amid-market-uncertainty/}
\bibitem{lg-loss} Investing.com, ``LG Energy Solution flags slowing EV demand, posts first quarterly loss in 3 years''; Reuters on Panasonic forecast cuts (2025). \url{https://investing.com/news/stock-market-news/lg-energy-solution-flags-slowing-ev-demand-posts-first-quarterly-loss-in-3-years-3828755}
\bibitem{bess-2025} Energy-Storage.News, ``China deploys 65 GWh of BESS in December, $\sim$65\% of 2025 global total.'' \url{https://www.energy-storage.news/china-deploys-65gwh-of-bess-in-december-25-of-2025-global-total/}

% ============ EV ============
\bibitem{wan-gang} MIT Technology Review, ``How did China come to dominate the world of electric cars?'' (Feb.\ 2023). \url{https://www.technologyreview.com/2023/02/21/1068880/how-did-china-dominate-electric-cars-policy/}
\bibitem{ten-cities} DieselNet, ``China: New Energy Vehicles'' (policy summary: 863 program, Ten Cities, dual credit, tax exemption). \url{https://dieselnet.com/standards/cn/nev.php}
\bibitem{csis-230b} S.~Kennedy, CSIS, ``The Chinese EV dilemma: subsidized yet striking'' (\$230.9B cumulative support 2009--2023). \url{https://www.csis.org/blogs/trustee-china-hand/chinese-ev-dilemma-subsidized-yet-striking}
\bibitem{dual-credit} ICCT, ``China's NEV mandate (dual-credit) policy update'' (2018). \url{https://theicct.org/sites/default/files/publications/China-NEV-mandate_ICCT-policy-update_20032018_vF-updated.pdf}
\bibitem{chargers} EVCIPA via ChinaEVHome, ``China's EV charging facilities hit 19.32 million by Nov.\ 2025''; US comparison $\sim$250k public ports. \url{https://chinaevhome.com/2025/12/23/chinas-ev-charging-facilities-hits-19-32-million-by-nov-2025-up-52-yoy/}
\bibitem{cpca-2025} CnEVPost, ``China NEV retail Dec.\ 2025 (CPCA)'': FY2025 12.86M, penetration 54.07\%, Dec.\ 60.4\%. \url{https://cnevpost.com/2026/01/07/china-nev-retail-1-387-million-dec-2025-preliminary-cpca/}
\bibitem{caam-2025} CnEVPost, ``China NEV sales Dec.\ 2025 (CAAM)'': FY2025 NEV 16.49M; total 34.4M; exports 7.1M incl.\ 2.62M NEVs. \url{https://cnevpost.com/2026/01/14/china-nev-sales-1-71-million-dec-2025-caam/}
\bibitem{exports-june-2026} CarNewsChina, ``China's monthly vehicle exports exceed 1 million for the first time'' (Jun.\ 2026). \url{https://carnewschina.com/2026/07/10/chinas-monthly-vehicle-exports-exceed-1-million-for-the-first-time-with-nevs-claiming-over-half/}
\bibitem{byd-2025} Electrek, ``BYD crushes Tesla all-electric sales for 2025, secures global BEV crown'' (4.55M total; 2.25M BEV vs.\ Tesla 1.64M). \url{https://electrek.co/2026/01/02/byd-crushes-tesla-all-electric-sales-for-2025-secures-global-bev-crown/}
\bibitem{rhodium-evs} Rhodium Group, ``Why are Chinese EVs so cheap?'' (BYD $\sim$80\% Tier-1 in-house; $\sim$\$4,700/vehicle structural advantage vs.\ Tesla; foreign share 62\%$\rightarrow$35\%). \url{https://rhg.com/research/why-are-chinese-evs-so-cheap/}
\bibitem{byd-ships} CarNewsChina, ``BYD completed its massive fleet, now able to export 1 million cars a year'' (Oct.\ 2025). \url{https://carnewschina.com/2025/10/02/byd-completed-its-massive-fleet-now-able-to-export-1-million-cars-a-year-but-not-just-from-china/}
\bibitem{chery} CarNewsChina, ``One Chery vehicle was exported every 23 seconds in 2025'' (1.34M exports). \url{https://carnewschina.com/2026/01/02/one-chery-vehicle-was-exported-every-23-seconds-in-2025-maintaining-its-lead-position-for-23-years/}
\bibitem{mg-million} CarNewsChina, ``MG is first Chinese auto brand to sell more than one million vehicles in Europe'' (Feb.\ 2026). \url{https://carnewschina.com/2026/02/22/mg-is-first-chinese-auto-brand-to-sell-more-than-one-million-vehicles-in-europe-and-uk/}
\bibitem{xiaomi} CarNewsChina, ``Xiaomi YU7 got over 289,000 orders in 60 minutes'' (Jun.\ 2025); ``Xiaomi hits 600,000 deliveries in 22 months'' (Feb.\ 2026). \url{https://carnewschina.com/2026/02/13/xiaomi-hits-600000-deliveries-in-22-months-sets-550000-unit-2026-target/}
\bibitem{shakeout} EVBoosters, ``400 Chinese EV companies ceased operations between 2018--2025.'' \url{https://evboosters.com/ev-charging-news/400-chinese-ev-companies-ceased-operations-between-2018-2025-only-a-few-will-dominate-towards-2030/}
\bibitem{gm-collapse} Fortune, ``GM in China: from 4.04M (2017) to 1.8M (2024), \$5B+ in charges'' (Jan.\ 2025). \url{https://fortune.com/asia/2025/01/28/gm-china-electric-vehicles-byd-buick-mary-barra/}
\bibitem{vw-decline} CarNewsChina, ``Volkswagen Group China sales declined 8\% y/y in 2025 (2.69M).'' \url{https://carnewschina.com/2026/01/13/volkswagen-groups-sales-declined-by-8-y-o-y-in-china-as-the-top-selling-foreign-auto-brand-in-2025/}
\bibitem{reverse-flow} CarNewsChina (exclusive), ``Volkswagen to licence XPeng's autonomous driving solution for its China EVs in 2026''; VW Group and Stellantis-Leapmotor releases. \url{https://carnewschina.com/2025/10/06/exclusive-volkswagen-to-licence-xpengs-autonomous-driving-solution-for-its-china-evs-in-2026/}
\bibitem{us-tariff} The White House, ``President Biden takes action to protect American workers...'' (100\% EV tariff, May 14, 2024). \url{https://bidenwhitehouse.archives.gov/briefing-room/statements-releases/2024/05/14/fact-sheet-president-biden-takes-action-to-protect-american-workers-and-businesses-from-chinas-unfair-trade-practices/}
\bibitem{eu-cvd} Cleary Trade Watch, ``Definitive duties adopted by the EU on Chinese battery electric vehicles'' (Oct.\ 2024: SAIC 35.3\%, Geely 18.8\%, BYD 17.0\%); electrive on the Jan.\ 2026 minimum-price guidance. \url{https://www.clearytradewatch.com/2024/10/definitive-duties-adopted-by-the-eu-on-chinese-battery-electric-vehicles-to-counteract-subsidies-to-apply-by-october-30/}
\bibitem{lidar} ChinaEVHome, ``Hesai leads global ADAS lidar shipments as China suppliers take 95\% share'' (Yole, May 2026). \url{https://chinaevhome.com/2026/05/06/hesai-leads-global-adas-lidar-shipments-as-china-suppliers-take-95-share/}
\bibitem{gods-eye} paultan.org, ``BYD unveils God's Eye DiPilot ADAS suite for all models'' (Feb.\ 2025). \url{https://paultan.org/2025/02/13/byd-unveils-gods-eye-dipilot-adas-suite-to-roll-out-in-all-byd-denza-yangwang-models-sold-in-china/}
\bibitem{mckinsey-2026} Fortune (citing McKinsey), ``Chinese cars frequently more advanced; 20--24-month development cycles; $\sim$30\% lower BOM'' (Jul.\ 30, 2026). \url{https://fortune.com/2026/07/30/byd-germany-dealerships-detroit-china-ev-threat/}
\bibitem{csis-rare-earth} CSIS, ``Rare earth export restrictions: one year later'' (Apr.\ 2025 controls; production stoppages; Oct.\ 2025 escalation and truce). \url{https://www.csis.org/analysis/rare-earth-export-restrictions-one-year-later}

% ============ MODEL ============
\bibitem{price-war} KrASIA, ``LLM prices hit rock bottom in China as Alibaba Cloud enters the fray'' (May 2024); Fortune Asia on the 97\% cuts. \url{https://kr-asia.com/llm-prices-hit-rock-bottom-in-china-as-alibaba-cloud-enters-the-fray}
\bibitem{deepseek-v3} DeepSeek-V3 Technical Report (arXiv:2412.19437); Stanford FSI, ``Taking stock of the DeepSeek shock.'' \url{https://cyber.fsi.stanford.edu/publication/taking-stock-deepseek-shock}
\bibitem{r1-shock} Yahoo Finance, ``Nvidia stock plummets, loses record \$589 billion as DeepSeek prompts questions over AI spending'' (Jan.\ 27, 2025). \url{https://finance.yahoo.com/news/nvidia-stock-plummets-loses-record-589-billion-as-deepseek-prompts-questions-over-ai-spending-135105824.html}
\bibitem{r1-nature} CNN Business, ``DeepSeek discloses R1 RL training cost of \$294,000 in peer-reviewed Nature paper'' (Sep.\ 2025). \url{https://www.cnn.com/2025/09/19/business/deepseek-ai-training-cost-china-intl}
\bibitem{digitalapplied} Digital Applied, ``Chinese AI models Q2 2026 market share report'' (release/pricing compendium; secondary source---verify against vendor blogs). \url{https://www.digitalapplied.com/blog/chinese-ai-models-q2-2026-market-share-report}
\bibitem{openrouter-2026} OpenRouter, ``The open-weight models that matter, June 2026'' (GLM-5.2 \#1 open on AA index; ranking of open models). \url{https://openrouter.ai/blog/insights/the-open-weight-models-that-matter-june-2026/}
\bibitem{hf-spring-2026} Hugging Face, ``State of Open Source, Spring 2026'' (Chinese models 41\% of downloads; Qwen \textgreater113k derivatives). \url{https://huggingface.co/blog/huggingface/state-of-os-hf-spring-2026}
\bibitem{qwen3} Qwen team, ``Qwen3'' release blog (Apr.\ 2025: 36T tokens, Apache 2.0, hybrid reasoning). \url{https://qwenlm.github.io/blog/qwen3/}
\bibitem{kimi-k2-thinking} S.~Willison, ``Kimi K2 Thinking'' (Nov.\ 2025: 1T/32B, INT4 QAT, HLE 44.9 vs GPT-5 41.7). \url{https://simonwillison.net/2025/Nov/6/kimi-k2-thinking/}
\bibitem{kimi-k3} Tom's Hardware, ``Moonshot releases 2.8-trillion-parameter Kimi K3 as China works around U.S. compute limits'' (Jul.\ 2026); VentureBeat coverage. \url{https://www.tomshardware.com/tech-industry/artificial-intelligence/moonshot-releases-2-8-trillion-parameter-kimi-k3}
\bibitem{deepseek-v4} DeepSeek V4 release documentation and third-party analysis (Apr.\ 2026: 1.6T/49B active, MIT license, DSA). \url{https://www.morphllm.com/deepseek-v4}
\bibitem{glm5} Z.ai, GLM-5/GLM-5.2 releases (MIT license; \#1 open-weight, AA Intelligence Index v4.1). \url{https://www.eigent.ai/blog/glm-5-2}
\bibitem{ernie-open} TechNode, ``Baidu open-sources ERNIE 4.5 series'' (Jun.\ 30, 2025, Apache 2.0). \url{https://technode.com/2025/07/01/baidu-open-sources-ernie-4-5-series-models-including-multimodal-moe-architecture/}
\bibitem{ipos} MIT Technology Review, ``What's next for Chinese open-source AI'' (Z.ai and MiniMax HK IPOs, Jan.\ 2026). \url{https://www.technologyreview.com/2026/02/12/1132811/whats-next-for-chinese-open-source-ai/}
\bibitem{token-econ} Seoul Economic Daily, ``China eyes AI supremacy with tokens at one-twentieth of US'' (May 2026; 140T daily token calls). \url{https://en.sedaily.com/international/2026/05/08/china-eyes-ai-supremacy-with-tokens-at-one-twentieth-of-us}
\bibitem{stanford-aiindex-2026} Stanford HAI, \emph{AI Index Report 2026} (US--China top-model gap 2.7\%; US \$285.9B vs China \$12.4B private investment). \url{https://hai.stanford.edu/ai-index/2026-ai-index-report}
\bibitem{aei-2026} AEI, ``China has caught up in frontier AI'' (Jul.\ 2026: release lag now ``a matter of weeks''; ``frontier AI is quickly becoming a commons''). \url{https://www.aei.org/foreign-and-defense-policy/china-has-caught-up-in-frontier-ai/}
\bibitem{datagravity} Data Gravity, ``China's open-weight takeover'' (OpenRouter token-share series; Llama \textless1\%; provider shares). \url{https://www.datagravity.dev/p/chinas-open-weight-takeover}
\bibitem{casado} M.~Casado (a16z) via The Economist and X (Nov.\ 2025): $\sim$80\% of open-stack startups use Chinese models. \url{https://x.com/martin_casado/status/1990462245541982546}
\bibitem{gpt-oss} OpenAI, ``Introducing gpt-oss'' (Aug.\ 5, 2025). \url{https://openai.com/index/introducing-gpt-oss/}
\bibitem{action-plan} PRC MFA, ``Global AI Governance Action Plan'' (WAIC, Jul.\ 2025), full text. \url{https://www.fmprc.gov.cn/mfa_eng/xw/zyxw/202507/t20250729_11679232.html}
\bibitem{carnegie-pivot} Carnegie Endowment, ``China's pivot on global AI'' (May 2026: WAICO, UN Group of Friends, BRICS AI centre). \url{https://carnegieendowment.org/research/2026/05/chinas-pivot-on-global-ai}
\bibitem{ai-plus} N.~Lambert, ``On China's open-source AI trajectory'' (Interconnects; State Council ``AI+'' plan targets). \url{https://www.interconnects.ai/p/on-chinas-open-source-ai-trajectory}
\bibitem{mittr-open} MIT Technology Review, ``What's next for Chinese open-source AI'' (Liu Zhiyuan quote; academic-credit policy). \url{https://www.technologyreview.com/2026/02/12/1132811/whats-next-for-chinese-open-source-ai/}
\bibitem{raschka} S.~Raschka, ``The technical DeepSeek story'' (MLA, GRPO, DSA analysis). \url{https://magazine.sebastianraschka.com/p/technical-deepseek}
\bibitem{ascend-delay} The Register, ``DeepSeek reverts to Nvidia after failed Ascend training runs'' (Aug.\ 2025); TrendForce coverage. \url{https://www.theregister.com/2025/08/14/dodgy_huawei_deepseek/}
\bibitem{ue8m0} DeepSeek V3.1 release notes (UE8M0 FP8 scale format ``for next-generation domestic chips,'' Aug.\ 2025; widely reported).

% ============ COMPUTE ============
\bibitem{riscv-policy} Reuters (exclusive), ``China to publish policy to boost RISC-V chip use nationwide'' (Mar.\ 3, 2025; eight agencies incl.\ MIIT, CAC, MOST). \url{https://money.usnews.com/investing/news/articles/2025-03-03/exclusive-china-to-publish-policy-to-boost-risc-v-chip-use-nationwide-sources}
\bibitem{fyp15} China Briefing, ``China's 15th Five-Year Plan: key insights'' (``extraordinary measures'' for ICs and basic software). \url{https://www.china-briefing.com/news/chinas-15th-five-year-plan-key-insights-for-foreign-investors/}
\bibitem{xiangshan} B.~Yungang et al., XiangShan project, RISC-V Summit Europe 2025 proceedings; Kunminghu microarchitecture. \url{https://riscv-europe.org/summit/2025/media/proceedings/2025-05-14-RISC-V-Summit-Europe-09h30-BAO-slides.pdf}
\bibitem{cas-riscv} Xinhua, ``CAS unveils RISC-V achievements at Zhongguancun Forum 2026'' (openRuyi, RVA23; 600+/400+ researchers per China Daily). \url{https://english.news.cn/20260327/3b8302ac28f64dedbfa3defe514f6e7e/c.html}
\bibitem{c930} Tom's Hardware, ``Alibaba launches RISC-V based XuanTie C930 server CPU'' (Mar.\ 2025). \url{https://www.tomshardware.com/pc-components/cpus/alibaba-launches-risc-v-based-xuantie-c930-server-cpu-ai-hpc-chip-ships-this-month-more-designs-to-follow}
\bibitem{spacemit} SCMP, ``China's SpacemiT to launch server-class RISC-V processor after raising US\$86 million'' (K3, Jan.\ 2026). \url{https://www.scmp.com/tech/tech-trends/article/3339953/chinas-spacemit-launch-server-class-risc-v-processor-after-raising-us86-million}
\bibitem{csis-riscv} CSIS, ``Sustaining standards leadership: the United States cannot disengage from RISC-V'' (membership statistics; 2B+ chips shipped). \url{https://www.csis.org/analysis/sustaining-standards-leadership-united-states-cannot-disengage-risc-v}
\bibitem{summit-2026} RISC-V International, ``RISC-V Summits'' (Summit China 2026: Shenzhen, Oct.\ 18--20). \url{https://riscv.org/community/risc-v-summits/}
\bibitem{matrix-ext} RISC-V International Tech Hub: IME/AME/VME task-group ratification plans (IME spec freeze Jun.\ 2026, ratification target Aug.\ 2026). \url{https://riscv.atlassian.net/wiki/spaces/IMEX/pages/598867969/IME+Ratification+Plan}
\bibitem{evas} EVAS Intelligence RISC-V AI SuperNode, WAIC 2026 (vendor claims: block-quantized FP8 cloud AI chip, 3.2 TB/s ELink). \url{https://finance.biggo.com/news/8eda6f31-9e18-42f8-9461-c5f142ed8a0d}
\bibitem{nvidia-timeline} Fortune, ``Nvidia and AMD to pay 15\% of China chip sale revenue to US government'' (Aug.\ 2025); export-control timeline. \url{https://fortune.com/2025/08/10/nvidia-amd-chips-h20-mi308-china-sales-revenue-trump-export-license}
\bibitem{cac-ban} TechCrunch, ``China tells its tech companies they can't buy AI chips from Nvidia'' (Sep.\ 17, 2025); Global Times on the CAC summons (Jul.\ 31, 2025). \url{https://techcrunch.com/2025/09/17/china-tells-its-tech-companies-they-cant-buy-ai-chips-from-nvidia/}
\bibitem{datacenter-mandate} Reuters (exclusive), ``China bans foreign AI chips from state-funded data centres'' (Nov.\ 2025); provincial power-tariff discounts for domestic-chip facilities. \url{https://www.aol.com/articles/exclusive-china-bans-foreign-ai-090709029.html}
\bibitem{nvidia-zero} NVIDIA, Q2 FY2026 results (zero H20 sales to China-based customers); TradingKey, ``Nvidia China revenue share 26\% $\rightarrow$ $\sim$5\%.'' \url{https://nvidianews.nvidia.com/news/nvidia-announces-financial-results-for-second-quarter-fiscal-2026}
\bibitem{huang-quotes} Axios, ``Jensen Huang: `China is going to win the AI race''' (Nov.\ 5, 2025, and same-day walk-back); Hill \& Valley remarks (Apr.\ 2025). \url{https://www.axios.com/2025/11/05/ai-nvidia-china-race}
\bibitem{semianalysis-ascend} SemiAnalysis, ``Huawei Ascend production ramp'' (2024: 507k units; 2025: $\sim$805k; die bank 2.9M; HBM stockpiles; SMIC 45k$\rightarrow$80k wpm). \url{https://newsletter.semianalysis.com/p/huawei-ascend-production-ramp}
\bibitem{huawei-roadmap} Huawei Connect 2025 keynote (E.~Xu), ``Ascend 950/960/970 roadmap; Atlas 950/960 SuperPoD; TaiShan 950; UnifiedBus 2.0'' (Sep.\ 18, 2025). \url{https://www.huawei.com/en/news/2025/9/hc-xu-keynote-speech}
\bibitem{cambricon} Tom's Hardware, ``Cambricon targets 500,000 AI chips in 2026''; FY2025 revenue +453\%, first annual profit. \url{https://www.tomshardware.com/tech-industry/semiconductors/cambricon-targets-500000-ai-chips-in-2026-as-china-accelerates-domestic-hardware-push}
\bibitem{ipos-chips} CNBC, ``Moore Threads' \$1.1B Shanghai IPO surges $\sim$400\% on debut'' (Dec.\ 2025); Asia Financial on Biren's HK IPO (Jan.\ 2026). \url{https://www.cnbc.com/2025/12/05/china-nvidia-moore-threads-trading-debut-1-billion-listing-ipo-shanghai-gpu-enflame-biren.html}
\bibitem{kunlun} technology.org, ``Baidu lights up 30,000-strong Kunlun P800 chip cluster'' (Apr.\ 2025). \url{https://www.technology.org/2025/04/25/baidu-lights-up-30000-strong-kunlun-chip-cluster-aims-to-train-deepseek-like-ai-models/}
\bibitem{ppu} TrendForce, ``Alibaba T-Head PPU lands China Unicom deal, reportedly matches Nvidia H20'' (Sep.\ 2025); CNBC on the Zhenwu 10k-chip cluster (Apr.\ 2026). \url{https://www.trendforce.com/news/2025/09/17/news-alibaba-t-head-ppu-lands-china-unicom-deal-reportedly-matches-nvidia-h20-performance/}
\bibitem{anke} Tom's Hardware, ``China certifies nine domestic AI chips for government procurement'' (May 2026). \url{https://www.tomshardware.com/tech-industry/semiconductors/china-certifies-nine-domestic-ai-chips-for-government-procurement}
\bibitem{domestic-share} TradingKey, ``Domestic AI chip shipments $\sim$1.65M units, 41\% of China market in 2025; Huawei $\sim$812k.'' \url{https://www.tradingkey.com/analysis/stocks/us-stocks/261889036-nvidia-china-revenue-airforceone-h200-zero-sales-huang-trump-visit-tradingkey}
\bibitem{miit-eflops} Xinhua (MIIT data), ``China's intelligent computing capacity reaches 2,185 EFLOPS, utilization 71.4\%'' (Jul.\ 2026). \url{https://english.news.cn/20260720/c5ebd0bd866b46dcb221f59a47a38d82/c.html}
\bibitem{miit-interconnect} MIIT, Computing Power Interconnection Action Plan (2025) and ``AI+'' implementation documents on unified AI-chip interconnection ecosystems (Chinese-language primary sources; see also Pandaily summary). \url{https://pandaily.com/miit-optoelectronic-chip-rd-ai-communications-2026}
\bibitem{huawei-ub} CRN Asia, ``Huawei's open source gambit: SuperPoD reference architecture and UnifiedBus 2.0 released openly'' (Sep.\ 2025). \url{https://www.crnasia.com/news/2025/artificial-intelligence/huawei-s-open-source-gambit-chinese-tech-giant-makes-superpo}
\bibitem{huawei-cann} Xinhua, ``Huawei fully open-sources CANN'' (Aug.\ 7, 2025). \url{https://english.news.cn/20250807/35025b0c1366473f976e29e0c5adf88f/c.html}
\bibitem{hbm-wafer} Network World, ``Server memory prices could double by 2026 as AI demand strains supply'' (HBM $\sim$3$\times$ wafer capacity per bit; Grace LPDDR5X shift). \url{https://www.networkworld.com/article/4093752/server-memory-prices-could-double-by-2026-as-ai-demand-strains-supply.html}
\bibitem{riken-study} RIKEN Center for Computational Science, feasibility study of memory-bandwidth balance for 2.8T-parameter sparse-MoE pretraining at 20 PF(FP8) per accelerator: HBM4 capacity wall at \textless256 accelerators; LPDDR6X substitution at 2.3 TB/socket, 5.0 TB/s knee, 67\% vs.\ 89\% MFU ceilings (2026, companion report to this volume).
\bibitem{cxmt-lpddr6} wccftech, ``CXMT's LPDDR6 clears R\&D verification at 12.8 Gbps'' (Aug.\ 2026); Gizmochina on 2H26 mass-production targeting. \url{https://wccftech.com/cxmts-lpddr6-clears-rd-verification-at-12-8-gbps-inching-closer-to-a-full-mass-production-push-report/}
\bibitem{socamm2} SK hynix newsroom, ``Mass production of SOCAMM2 192GB'' (Apr.\ 2026); TechPowerUp on Qualcomm/AMD adoption. \url{https://news.skhynix.com/en/mass-production-socamm2-192gb/}
\bibitem{cloudmatrix} SemiAnalysis (Apr.\ 2025), CloudMatrix 384 analysis: $\sim$2$\times$ GB200 NVL72 BF16 throughput at $\sim$3.9$\times$ power. \url{https://techstartups.com/2025/04/28/huaweis-ascend-910c-system-reportedly-outperforms-nvidias-h100-in-key-metrics/}
\bibitem{ifp-b30a} Institute for Progress, ``The B30A decision'' (2026 output projections: US 6.89M vs China $\sim$167k B300-equivalents). \url{https://ifp.org/the-b30a-decision/}

% ============ SEMICONDUCTOR ============
\bibitem{smic-5nm} Design \& Reuse, ``SMIC achieves 5nm production on N+3 node without EUV tools'' (TechInsights teardown of Kirin 9030, Dec.\ 2025). \url{https://www.design-reuse.com/news/202529830-chinese-smic-achieves-5-nm-production-on-n-3-node-without-euv-tools/}
\bibitem{yields} SupplyICs market intelligence, ``China logic fabs: yield and geopolitical risks 2026'' (N+2 yields est.\ 40--55\%; treat as estimates). \url{https://supplyics.com/insights/market-intelligence/china-logic-fabs-yield-geopolitical-risks-2026/}
\bibitem{smic-2025} TrendForce, ``SMIC posts record \$9.3B in 2025 sales'' (Feb.\ 2026; utilization 93.5\%; capex \$8.1B). \url{https://www.trendforce.com/news/2026/02/11/news-smic-posts-record-9-3b-in-2025-sales-7nm-yields-reportedly-weigh-on-margins/}
\bibitem{shadow-fabs} SemiWiki/Digitimes, ``Huawei's quiet IDM ambition: 11 shadow fabs across China''; DSET, ``Uncovering Huawei's shadow network.'' \url{https://dset.tw/en/research/uncovering-huaweis-shadow-network/}
\bibitem{mature-node} TrendForce, ``China 37\% of global 22--40nm output in 2026, 42\% by 2028 (SEMI China)''; ``Chinese foundries drive $\sim$70\% of 12-inch mature-node expansion'' (May 2026). \url{https://www.trendforce.com/presscenter/news/20260507-13036.html}
\bibitem{auto-chips} TrendForce, ``Chinese carmakers reportedly aim for 100\% domestic chips in vehicles by 2026--27'' (Jun.\ 2025). \url{https://www.trendforce.com/news/2025/06/18/news-chinese-carmakers-reportedly-aim-for-100-domestic-chips-in-vehicles-by-2026/}
\bibitem{cxmt-capacity} wccftech, ``CXMT races to increase capacity to 600,000 wafers per month'' (Jul.\ 2026; current 200--300k wpm). \url{https://wccftech.com/cxmts-64gb-server-ddr5-memory-chip-is-now-pricier-than-samsungs-as-it-races-to-increase-capacity-to-600000-wafers-per-month/}
\bibitem{citrini} Tom's Hardware (Citrini Research), ``CXMT close to matching Micron's memory capacity in 2026'' (350k wpm model; China $\sim$1.4M wpm DRAM by 2030). \url{https://www.tomshardware.com/pc-components/dram/cxmt-close-to-matching-microns-memory-capacity-in-2026-research-claims-would-put-china-on-track-to-become-worlds-second-largest-dram-producer}
\bibitem{cxmt-ddr5} wccftech, ``CXMT debuts DDR5-8000 and LPDDR5X-10667'' (Nov.\ 2025); TrendForce on 80\% DDR5 yield (Dec.\ 2024). \url{https://wccftech.com/cxmt-debuts-domestically-produced-ddr5-memory-8000-mtps-lpddr5-10667-mtps/}
\bibitem{cxmt-ipo} CNBC, ``CXMT closes 466\% higher in Shanghai debut'' (Jul.\ 27, 2026; \$8B raised; largest A-share company; prospectus disclosed no commercial HBM). \url{https://www.cnbc.com/2026/07/27/cxmt-china-market-debut-chipmaker-ipo.html}
\bibitem{chinatalk-hbm} ChinaTalk, ``Mapping China's HBM advancement'' (CXMT HBM roadmap; $\sim$4-year gap to leaders). \url{https://www.chinatalk.media/p/mapping-chinas-hbm-advancement}
\bibitem{packaging} Tom's Hardware, ``Chinese memory maker gets \$2.4 billion to build HBM packaging for AI processors'' (Innotron TSV facility); TrendForce on Tongfu/JCET expansions. \url{https://www.tomshardware.com/pc-components/dram/chinese-memory-maker-gets-dollar24-billion-to-build-hbm-for-ai-processors-shanghai-packaging-facility-to-open-in-2026}
\bibitem{dram-shortage} TrendForce, ``Conventional DRAM up 55--60\% QoQ in 1Q26'' (Jan.\ 2026); Network World on the AI-driven shortage. \url{https://www.trendforce.com/presscenter/news/20260105-12860.html}
\bibitem{cxmt-pricing} wccftech, ``CXMT's 64GB server DDR5 is now pricier than Samsung's'' (Jul.\ 2026). \url{https://wccftech.com/cxmts-64gb-server-ddr5-memory-chip-is-now-pricier-than-samsungs-as-it-races-to-increase-capacity-to-600000-wafers-per-month/}
\bibitem{ymtc} Digitimes, ``YMTC Xtacking 4.0 powers 267-layer NAND''; TechPowerUp/Guru3D on the 294-layer 5th generation. \url{https://www.digitimes.com/news/a20250924PD216/ymtc-nand-nand-flash-production-memory.html}
\bibitem{toolmakers-top20} TrendForce, ``Three Chinese chip toolmakers enter global top 20 as Naura rises to fifth by sales'' (Feb.\ 2026; domestic tool share 20--30\%). \url{https://www.trendforce.com/news/2026/02/02/news-three-chinese-chip-toolmakers-reportedly-enter-global-top-20-as-naura-rises-to-fifth-by-sales/}
\bibitem{naura} BigGo Finance, ``Naura 2025 results: revenue $\sim$RMB 39.4B (+31\%)'' (Apr.\ 2026). \url{https://finance.biggo.com/news/9r28m50BvthpMgHBYd-l}
\bibitem{sicarrier} Yahoo Finance, ``SiCarrier targets \$2.8bn initial raise at $\sim$\$11B valuation''; TrendForce on its SAQP-DUV 5nm-class patent. \url{https://www.trendforce.com/news/2025/11/10/news-decoding-chinas-lithography-push-to-challenge-asml-from-sicarrier-to-alternative-euv-paths/}
\bibitem{litho} TrendForce, ``Decoding China's lithography push to challenge ASML'' (Nov.\ 2025: SMEE 90nm shipping, 28nm immersion in development, Yuliangsheng trials; domestic sub-10nm litho $\sim$2030); eeNews on LDP-EUV reports (speculative). \url{https://www.trendforce.com/news/2025/11/10/news-decoding-chinas-lithography-push-to-challenge-asml-from-sicarrier-to-alternative-euv-paths/}
\bibitem{eda} Tom's Hardware, ``White House lifts chip design software export ban on China'' (Jul.\ 2025); TrendForce on Empyrean's full-process memory EDA platform (Aug.\ 2025). \url{https://www.tomshardware.com/tech-industry/semiconductors/white-house-lifts-chip-design-export-ban-on-china-in-exchange-for-rare-earth-materials-compromise-export-licences-for-eda-software-sales-no-longer-required}
\bibitem{wafers} TechWire Asia, ``China semiconductor self-sufficiency: wafer target 2026'' (300mm wafer share 3\% (2020) $\rightarrow$ $\sim$28--32\% (2026); Eswin pricing). \url{https://techwireasia.com/2026/05/china-semiconductor-self-sufficiency-wafer-target-2026/}
\bibitem{bigfund} Asia Society Policy Institute, ``China's Big Fund 3.0: Xi's boldest gamble yet for chip supremacy'' (RMB 344B, May 2024, horizon 2039). \url{https://asiasociety.org/policy-institute/chinas-big-fund-30-xis-boldest-gamble-yet-chip-supremacy}
\bibitem{semi-wfe} SCMP (SEMI), ``China to remain world's biggest buyer of chipmaking equipment through 2027''; SEMI 2026 equipment forecasts. \url{https://www.scmp.com/tech/article/3336933/china-remain-worlds-biggest-buyer-chipmaking-equipment-through-2027-semi}
\bibitem{self-sufficiency} USCC, ``Made in China 2025: Evaluating China's Performance'' (Nov.\ 2025); TechWire Asia on 2026 self-sufficiency measures. \url{https://www.uscc.gov/sites/default/files/2025-11/Made_in_China_2025--Evaluating_Chinas_Performance.pdf}
\bibitem{bis-dec24} US BIS, ``Commerce strengthens export controls...'' (Dec.\ 2, 2024: 140 entities; HBM and SME controls). \url{https://www.bis.gov/press-release/commerce-strengthens-export-controls-restrict-chinas-capability-produce-advanced-semiconductors-military}
\bibitem{policy-whiplash} Built In, ``Timeline: US AI-chip policy toward China 2025--2026'' (AI Diffusion Rule imposed and rescinded; H20/H200 reversals). \url{https://builtin.com/articles/trump-lifts-ai-chip-ban-china-nvidia}
\bibitem{smuggling} CNBC, ``US tech execs smuggled Nvidia chips to China, prosecutors say'' (\$2.5B scheme, Mar.\ 2026). \url{https://www.cnbc.com/2026/03/19/us-tech-execs-smuggled-nvidia-chips-to-china-prosecutors-say.html}

% ============ MODEL SURVEY / QUANTIZATION ============
\bibitem{qwen-billion} OpenSourceForU, ``Alibaba's Qwen crosses one billion downloads, eclipsing Meta's Llama'' (Jul.\ 2026); Xinhua on the 700M milestone (Jan.\ 2026). \url{https://www.opensourceforu.com/2026/07/alibabas-qwen-crosses-one-billion-downloads-eclipsing-metas-llama/}
\bibitem{sos-ai} State of Open Source AI report v1.0.1 (Jul.\ 2026): DeepSeek 26,000+ enterprise deployments; Chinese models $\sim$45\% of OpenRouter weekly traffic (Apr.\ 2026). \url{https://stateofopensource.ai/}
\bibitem{menlo} Menlo Ventures 2025 enterprise LLM survey: open-weight share of enterprise workloads 19\% (2024) $\rightarrow$ 11\% (2025). \url{https://www.luizneto.ai/enterprise-open-weight-models-2026/}
\bibitem{unsloth-r1} Unsloth AI, ``Run DeepSeek-R1 Dynamic 1.58-bit'' (Jan.\ 2025: 720GB $\rightarrow$ 131GB; dynamic vs naive quantization results). \url{https://unsloth.ai/blog/deepseekr1-dynamic}
\bibitem{unsloth-models} Unsloth documentation: Dynamic 2.0 GGUFs methodology; Kimi K3 (594GB--1.5TB), GLM-5 (176--241GB), K2 Thinking 1-bit (245GB) quantizations with retention figures. \url{https://unsloth.ai/docs/basics/unsloth-dynamic-2.0-ggufs}
\bibitem{unsloth-yc} Y Combinator company profile: Unsloth AI ($\sim$10M monthly model downloads). \url{https://www.ycombinator.com/companies/unsloth-ai}
\bibitem{awq} Lin et al., ``AWQ: Activation-aware Weight Quantization'' (MLSys 2024 Best Paper); Frantar et al., ``GPTQ'' (ICLR 2023). \url{https://github.com/mit-han-lab/llm-awq}
\bibitem{vllm-quant} vLLM documentation, supported quantization methods (AWQ, GPTQ, FP8, NVFP4, GGUF, bitsandbytes). \url{https://docs.vllm.ai/en/latest/features/quantization/}
\bibitem{ktransformers} KTransformers (Tsinghua kvcache-ai), SOSP 2025: DeepSeek-R1/V3 on one RTX 4090 + dual Xeon, 286 tok/s prefill via AMX. \url{https://github.com/kvcache-ai/ktransformers/blob/main/doc/en/DeepseekR1_V3_tutorial.md}
\bibitem{epyc-rig} Digital Spaceport, ``How to run DeepSeek-R1 671B fully locally on a \$2,000 EPYC rig'' (Jan.\ 2025). \url{https://digitalspaceport.com/how-to-run-deepseek-r1-671b-fully-locally-on-2000-epyc-rig/}
\bibitem{mac-frontier} MacRumors on M3 Ultra running R1 at 17--18 tok/s (Mar.\ 2025); A.~Hannun on 2$\times$M3 Ultra running Kimi K2/K2.5; Creative Strategies, ``Running a 1T-parameter model on a \$40k Mac Studio cluster'' (Dec.\ 2025). \url{https://creativestrategies.com/research/running-a-1t-parameter-model-on-a-40k-mac-studio-cluster/}
\bibitem{spark-halo} AIMultiple, ``DGX Spark alternatives'' benchmark roundup (Spark, Strix Halo, Mac Studio, multi-3090 gpt-oss-120b figures); Atlantic Council on fall-2026 local-AI laptops. \url{https://aimultiple.com/dgx-spark-alternatives}
\bibitem{llamacpp-moe} llama.cpp PR \#15077, `--n-cpu-moe` MoE expert CPU-offload (Aug.\ 2025); day-0 native MXFP4 support. \url{https://github.com/ggml-org/llama.cpp/pull/15077}
\bibitem{llmflation} G.~Appenzeller (a16z), ``LLMflation: LLM inference cost is falling 10$\times$ per year'' (Nov.\ 2024). \url{https://a16z.com/llmflation-llm-inference-cost/}
\bibitem{selfhost} OECD/McKinsey self-hosting break-even figures via luizneto.ai enterprise analysis (2026). \url{https://www.luizneto.ai/enterprise-open-weight-models-2026/}
\bibitem{atlantic-council} Atlantic Council, ``The best AI you can own is Chinese---the West needs to close that gap quickly'' (2026; Microsoft's DeepSeek V4 Copilot evaluation; Western app bans vs self-hosted weights). \url{https://www.atlanticcouncil.org/blogs/the-best-ai-you-can-own-is-chinese-the-west-needs-to-close-that-gap-quickly/}

% ============ LINESHINE / LX2 / DHM ============
\bibitem{top500-lineshine} TOP500, ``LineShine debuts at No.\ 1 as TOP500 enters new global exascale era'' (Jun.\ 24, 2026: HPL 2.198 EF, HPCG 22.00 PF, 42.2 MW, 52.07 GF/W). \url{https://top500.org/news/lineshine-debuts-no-1-top500-enters-new-global-exascale-era/}
\bibitem{dongarra-report} J.~Dongarra, ``Report on the Chinese LineShine System,'' Univ.\ of Tennessee Tech Report ICL-UT-26-01 (Jun.\ 23, 2026). \url{https://glennklockwood.com/garden/attachments/dongarra-lineshine-0623.pdf}
\bibitem{dhm-cacm} J.~Dongarra, T.~Hoefler, and S.~Matsuoka, ``Do We Still Need GPUs? Where an Evolved CPU Can Replace Them in AI and Scientific Computing---and Where It Cannot Yet,'' \emph{Communications of the ACM}, in revision (2026).
\bibitem{dhm-long} J.~Dongarra, T.~Hoefler, and S.~Matsuoka, ``Do We Still Need GPUs? Rethinking AI and Scientific Computing on Accelerated CPUs,'' full-length companion paper (2026): LX2 vs GPU-fleet analysis of 1T-MoE serving and training; the design-sequence/allocation argument.
\bibitem{lockwood} G.~K.~Lockwood, ISC'26 recap and LX2 processor notes (microarchitecture, HBM2e speculation, software stack analysis). \url{https://blog.glennklockwood.com/2026/07/isc26-recap.html}
\bibitem{globaltimes-lineshine} Global Times on LineShine (Jun.\ 2026): ``China's first domestically developed HBM'' claim; CAS/analyst commentary. \url{https://www.globaltimes.cn/page/202606/1364287.shtml}
\bibitem{dongarra-scmp} SCMP, ``Return to top: China's LineShine beats US' El Capitan in TOP500 supercomputer rankings'' (Jun.\ 24, 2026; Dongarra quotes on export controls). \url{https://www.scmp.com/news/china/science/article/3358107/return-top-chinas-lineshine-beats-us-el-capitan-top500-supercomputer-rankings}

% ============ WEST-TO-RISC-V ============
\bibitem{nvidia-riscv-cores} RISC-V International, ``How NVIDIA shipped one billion RISC-V cores in 2024'' (Feb.\ 2025; 10--40 cores per chip since 2016). \url{https://riscv.org/blog/how-nvidia-shipped-one-billion-risc-v-cores-in-2024/}
\bibitem{cuda-riscv} Tom's Hardware, ``NVIDIA's CUDA platform now supports RISC-V'' (announced at RISC-V Summit China, Jul.\ 18, 2025). \url{https://www.tomshardware.com/pc-components/gpus/nvidias-cuda-platform-now-supports-risc-v-support-brings-open-source-instruction-set-to-ai-platforms-joining-x86-and-arm}
\bibitem{cuda-riscv-analysis} The Register, ``Nvidia's CUDA-on-RISC-V as a China hedge'' (Jul.\ 21, 2025); SCMP on the boost to China's self-sufficiency drive. \url{https://www.theregister.com/2025/07/21/nvidia_cuda_riscv}
\bibitem{sifive-nvlink} SiFive press release, ``SiFive to integrate NVIDIA NVLink Fusion for next-gen AI datacenter chips'' (Jan.\ 15, 2026). \url{https://www.sifive.com/press/sifive-nvidia-nvlinkfusion-datacenter}
\bibitem{redmond} HPCwire, ``RISC-V CEO Calista Redmond resigns'' (Dec.\ 2024); subsequently NVIDIA VP of global AI initiatives. \url{https://www.hpcwire.com/off-the-wire/risc-v-ceo-calista-redmond-resigns-after-5years-of-progress/}
\bibitem{google-riscv} Data Center Dynamics, ``Google deploys SiFive's Intelligence X280 processor for AI workloads'' (Sep.\ 2022); Google Open Source Blog, ``Android and RISC-V'' (Oct.\ 2023); RISE founding membership (May 2023). \url{https://www.datacenterdynamics.com/en/news/google-deploys-sifives-intelligence-x280-processor-for-ai-workloads/}
\bibitem{mtia} Meta AI blog, MTIA v1/v2 (RISC-V processing elements); Data Center Dynamics, ``Meta unveils next four generations of its MTIA chip'' (Mar.\ 11, 2026). \url{https://www.datacenterdynamics.com/en/news/meta-unveils-next-four-generations-of-its-mtia-chip/}
\bibitem{cobalt} TechRadar, ``Microsoft unveils Cobalt 200'' (Nov.\ 2025) --- Azure silicon remains Arm-based; no public Microsoft RISC-V commitment as of mid-2026. \url{https://www.techradar.com/pro/microsoft-unveils-its-next-generation-arm-based-cpu-cobalt-200-looks-to-unlock-even-more-azure-power}
\bibitem{riscv-restrict} Bloomberg, ``Lawmakers urge export controls on RISC-V'' (Nov.\ 2023); Reuters on Commerce review (Apr.\ 2024); Sen.\ Warner letter to BIS (Sep.\ 30, 2025). No restriction enacted as of mid-2026. \url{https://www.warner.senate.gov/public/_cache/files/a/b/ab558775-b187-434e-b70d-040a63aea6f7/1AD7AF181B29E0603B1023A6FABE01893C9F53BFCA2ACAA82B1535CEBE2FC7AC.9.30.25-risc-v-letter-to-bis.pdf}
\bibitem{shd-2026} SHD Group 2026 forecast (Jun.\ 2026): 35.9B RISC-V SoC units and \$318B revenue by 2031, AI accelerators $\sim$70\% of revenue; RISC-V International 25\%-penetration announcement (Oct.\ 2025). \url{https://riscv.org/blog/shd-forecast-2026/}
\bibitem{omdia} Omdia, ``RISC-V adoption will be accelerated by AI'' (May 2024: $\sim$25\% of processors by 2030). \url{https://omdia.tech.informa.com/pr/2024/may/risc-v-adoption-will-be-accelerated-by-ai-according-to-new-omdia-research}

% ============ CXMT HBM DETAIL ============
\bibitem{cxmt-hbm3} DigiTimes, ``CXMT delivers HBM3 samples to Huawei; mass production slated 2026'' (Oct.\ 2025); TechPowerUp on the 20\%-of-capacity HBM3 line report. \url{https://www.digitimes.com/news/a20251023PD208/cxmt-hbm3-2026-dram-huawei.html}
\bibitem{xmc-hbm} TrendForce, ``Huawei and XMC collaborate to develop HBM'' (Jul.\ 2024); Tom's Hardware on Wuhan Xinxin's $\sim$3,000 wpm HBM line; SemiAnalysis, ``Scaling the Memory Wall.'' \url{https://newsletter.semianalysis.com/p/scaling-the-memory-wall-the-rise-and-roadmap-of-hbm}
\bibitem{hbm-teardown} As of mid-2026 no public teardown confirms Chinese-fabbed DRAM in a shipping accelerator; shipping Ascend 910C teardowns show Samsung/SK hynix stacks (SemiAnalysis; prediction-market tracking). \url{https://newsletter.semianalysis.com/p/huawei-ascend-production-ramp}
\bibitem{substrate} The Substrate, ``Where will China get its compute?'' (domestic HBM stack estimates, 2.5--18.6M range for 2026). \url{https://www.the-substrate.net/p/where-will-china-get-its-compute}
\bibitem{hbm-toolchain} Tom's Hardware, ``Chinese semiconductor industry gears up for domestic HBM3 production by end of 2026'' (Naura, Maxwell, U-Precision toolchain). \url{https://www.tomshardware.com/pc-components/dram/chinese-semiconductor-industry-gears-up-for-domestic-hbm3-production-by-the-end-of-2026-cxmt-to-produce-chips-while-naura-maxwell-and-u-preseason-design-tools-for-assembly}

% ============ THE WRONG SIDE OF THE TRADE ============
\bibitem{openai-valuation} Forbes/CNBC/Bloomberg, ``OpenAI valuation reaches \$852 billion after \$122B round'' (Mar.\ 31, 2026); prior \$500B secondary (Oct.\ 2025). \url{https://www.forbes.com/sites/antoniopequenoiv/2026/03/31/openai-valuation-reaches-852-billion-after-massive-funding-round/}
\bibitem{openai-rev} WSJ via Yahoo Finance, ``OpenAI lifts planned compute spending'' (Jul.\ 22, 2026: 2025 revenue \$13.1B; $\sim$\$2B/month early 2026). \url{https://finance.yahoo.com/technology/ai/articles/openai-lifts-planned-compute-spending-144917731.html}
\bibitem{openai-losses} WSJ/analyst projections: $\sim$\$14B 2026 loss; $\sim$\$115B cumulative losses through 2029; profitability targeted 2029--30. \url{https://valueaddvc.com/blog/openai-revenue-2026-25b-arr-14b-loss-and-why-profitability-waits-until-2030}
\bibitem{openai-commit} Altman's \$1.4T commitments figure (Nov.\ 2025); WSJ: planned compute spend raised to \$750B through 2030, \$1.15T agreements through 2035; Oracle \$300B, AWS $\sim$\$100B, Stargate \$500B/10GW, AMD 6GW + warrants, Broadcom 10GW custom. \url{https://finance.yahoo.com/technology/ai/articles/openai-lifts-planned-compute-spending-144917731.html}
\bibitem{openai-ipo} CNBC/Forbes: confidential S-1 filed Jun.\ 8, 2026, targeting up to $\sim$\$1T; reports of delay to 2027 (Jun.\ 25, 2026). \url{https://www.forbes.com/sites/aliciapark/2026/06/25/openai-considers-delaying-ipo-to-2027-after-spacexs-rocky-debut-report-says/}
\bibitem{altman-bubble} S.~Altman (Aug.\ 2025): AI investors ``overexcited''; ``someone will lose a phenomenal amount of money.'' \url{https://www.tipranks.com/news/sam-altman-warns-ai-stocks-may-be-in-a-dot-com-style-bubble}
\bibitem{altman-enough} TechCrunch, ``Sam Altman says `Enough' to questions about OpenAI's revenue'' (BG2 podcast, Nov.\ 2025). \url{https://techcrunch.com/2025/11/02/sam-altman-says-enough-to-questions-about-openais-revenue/}
\bibitem{anthropic-series-h} Anthropic, ``Series H'' (May 28, 2026: \$65B at \$965B post-money); CNBC coverage. \url{https://www.anthropic.com/news/series-h}
\bibitem{anthropic-econ} Bloomberg (Jan.\ 2026: $>$\$9B run-rate); Digitimes/TechCrunch on 2028 breakeven projection and $\sim$14$\times$ lower burn than OpenAI; Forbes on inference gross margins passing 70\% (Jul.\ 2026); Google 1GW+ TPU deal (Oct.\ 2025). \url{https://www.forbes.com/sites/petercohan/2026/07/28/as-token-costs-plunge-enterprise-ai-providers-face-a-new-margin-squeeze/}
\bibitem{nvidia-fy2026} NVIDIA FY2026 results (Feb.\ 2026: revenue \$215.9B, data center \$193.7B); \$500B Blackwell+Rubin bookings visibility (GTC Washington, Oct.\ 2025). \url{https://nvidianews.nvidia.com/news/nvidia-announces-financial-results-for-fourth-quarter-and-fiscal-2026}
\bibitem{nvidia-5t} CNBC/Axios, ``Nvidia becomes first \$5 trillion company'' (Oct.\ 29, 2025). \url{https://www.axios.com/2025/10/29/nvidia-5-trillion}
\bibitem{nvidia-concentration} CNBC, ``Nvidia's top two mystery customers made up 39\% of Q2 revenue'' (10-Q, Aug.\ 2025). \url{https://www.cnbc.com/2025/08/28/nvidias-top-two-mystery-customers-made-up-39percent-of-its-q2-revenue-.html}
\bibitem{nvidia-backstop} WSJ/CNBC, ``Nvidia and OpenAI in talks for up to \$250 billion AI backstop'' (Jul.\ 26--27, 2026); TechTimes: ``proves debt markets said no''; \$100B LOI restructured to $\sim$\$30B equity (Feb.\ 2026). \url{https://www.cnbc.com/2026/07/27/nvidia-and-openai-in-talks-for-up-to-250-billion-dollar-ai-backstop.html}
\bibitem{softbank-nvidia-sale} Bloomberg/AP, ``SoftBank sells entire Nvidia stake for \$5.83B'' (Nov.\ 11, 2025). \url{https://finance.yahoo.com/news/softbank-dumps-entire-nvidia-portfolio-145532866.html}
\bibitem{softbank-openai} MarketMinute, ``SoftBank shares retreat as Son's \$40B OpenAI gamble strains balance sheet'' (Mar.\ 27, 2026: $-$45\% from peak; Jefferies downgrade; LTV ceiling; $\sim$\$50B refinancing wall). \url{https://markets.financialcontent.com/stocks/article/marketminute-2026-3-27-softbank-shares-retreat-as-masayoshi-sons-40-billion-openai-gamble-strains-balance-sheet}
\bibitem{son-crying} CNBC, ``Masayoshi Son says he was `crying' over forced Nvidia sale; `all in' on OpenAI'' (Dec.\ 2, 2025). \url{https://www.cnbc.com/2025/12/02/softbank-masayoshi-son-crying-nvidia-stake-ai-bet-artificial-intelligence-bubble-vision-fund.html}
\bibitem{capex-2026} Introl, ``Hyperscaler capex \$600B+ in 2026''; Forbes (Jun.\ 2026): \$700--900B range, capital intensity 45--57\% of revenue. \url{https://introl.com/blog/hyperscaler-capex-600b-2026-ai-infrastructure-debt-january-2026}
\bibitem{mckinsey-capex} McKinsey, ``The cost of compute: a \$7 trillion race to scale data centers'' (Apr.\ 2025: \$6.7T by 2030). \url{https://www.mckinsey.com/industries/technology-media-and-telecommunications/our-insights/the-cost-of-compute-a-7-trillion-dollar-race-to-scale-data-centers}
\bibitem{ai-debt} Axios, ``The AI debt machine'' (Jul.\ 27, 2026: \$489B AI-linked issuance YTD; \$302B by five majors); Nikkei study of \$1.65T hidden obligations. \url{https://www.axios.com/2026/07/27/debt-data-center-oracle-goole}
\bibitem{meta-spv} Global Data Center Hub, ``Meta and Blue Owl's \$27B SPV'' (Oct.\ 2025); Meta \$30B bond, largest of 2025. \url{https://www.globaldatacenterhub.com/p/meta-blue-owls-27b-bet-is-this-the}
\bibitem{oracle-cds} Axios/CNBC: Oracle 5-yr CDS from 129bps (Feb.\ 2026) to record 212bps (Jul.\ 2026) on AI debt load. \url{https://www.axios.com/2026/07/27/debt-data-center-oracle-goole}
\bibitem{burry} CNBC, ``Michael Burry accuses AI hyperscalers of artificially boosting earnings'' (Nov.\ 11, 2025: \$176B understated depreciation 2026--28, earnings overstated $\sim$27\% by 2028). \url{https://www.cnbc.com/2025/11/11/big-short-investor-michael-burry-accuses-ai-hyperscalers-of-artificially-boosting-earnings.html}
\bibitem{nadella-power} Tom's Hardware/DCD, Nadella on BG2 (Nov.\ 2025): constraint is ``power, not chips''; GPUs ``sitting in inventory that I can't plug in.'' \url{https://www.datacenterdynamics.com/en/news/microsoft-has-ai-gpus-sitting-in-inventory-because-it-lacks-the-power-necessary-to-install-them/}
\bibitem{coreweave} CNBC / TechTimes: CoreWeave Q1 2026 revenue \$2.08B ($+$112\%); backlog \$99.4B; capex guide \$31--35B; Q1 FCF $-$\$4.7B; $\sim$\$21--25B debt; 5-yr CDS $\sim$855\,bps (29 Jul.\ 2026, $\sim$50\% implied default probability); stock $-$40\% over 12 months; Meta contract \$21B through 2032. \url{https://www.techtimes.com/articles/322222/20260730/coreweave-cds-hits-50-default-odds-bond-market-calls-time-gpu-debt-spiral.htm}
\bibitem{alphabet-fcf} Axios (Jul.\ 27, 2026): Alphabet's first negative-FCF quarter since its 2004 IPO, on AI capex. \url{https://www.axios.com/2026/07/27/debt-data-center-oracle-goole}
\bibitem{gdp-ai} BEA third-estimate analyses: AI/data-center capex $\sim$74\% of US Q1 2026 GDP growth; Furman: ex-data-centers H1 2025 growth $\sim$0.1\%. \url{https://www.tftc.io/ai-capex-gdp-growth-q1-2026-bea-third-estimate/}
\bibitem{mit-nanda} MIT Project NANDA, ``The GenAI Divide'' (Aug.\ 2025): $\sim$95\% of enterprise GenAI pilots produced zero P\&L impact. \url{https://www.forbes.com/sites/jasonsnyder/2025/08/26/mit-finds-95-of-genai-pilots-fail-because-companies-avoid-friction/}
\bibitem{bain-2t} Bain \& Company, 6th Global Technology Report (Sep.\ 2025): \$2T annual AI revenue needed by 2030; $\sim$\$800B shortfall. \url{https://www.bain.com/about/media-center/press-releases/20252/$2-trillion-in-new-revenue-needed-to-fund-ais-scaling-trend---bain--companys-6th-annual-global-technology-report/}
\bibitem{capex-divergence} Forbes, ``The AI capex-to-revenue gap is widening'' (Jun.\ 2, 2026: divergence rate 46\% vs 2001 telecom's 32\%). \url{https://www.forbes.com/sites/jasonkirsch/2026/06/02/the-ai-capex-to-revenue-gap-is-widening---and-markets-are-starting-to-notice/}
\bibitem{imf-boe} CNBC, ``IMF and Bank of England join growing chorus warning of an AI bubble'' (Oct.\ 8--9, 2025). \url{https://www.cnbc.com/2025/10/09/imf-and-bank-of-england-join-growing-chorus-warning-of-an-ai-bubble.html}
\bibitem{circular-deals} Bloomberg graphics, ``The AI circular deals'' (2026, $\sim$\$750B+ interlocking arrangements); Bernstein on circularity concerns. \url{https://www.bloomberg.com/graphics/2026-ai-circular-deals/}
\bibitem{chinese-share-forbes} Forbes, ``As token costs plunge, enterprise AI providers face a new margin squeeze'' (Jul.\ 28, 2026: Chinese suppliers $\sim$46\% global model share; Telnyx \$100k/day $\rightarrow$ \$100/day case). \url{https://www.forbes.com/sites/petercohan/2026/07/28/as-token-costs-plunge-enterprise-ai-providers-face-a-new-margin-squeeze/}
\bibitem{semis-selloff} Intellectia/CNN: \$1.3T semiconductor selloff (Jun.\ 4--9, 2026) on Broadcom AI-networking miss. \url{https://intellectia.ai/blog/semiconductor-stocks-selloff-june-2026}
\bibitem{memory-rout} CNBC/NBC (Jul.\ 28--29, 2026): memory-led chip rout, SK hynix $-$15\%, $>$\$1T lost, Nasdaq-100 in correction; Samsung/SK warnings of shortages to 2027+. \url{https://www.cnbc.com/2026/07/29/chip-selloff-sk-hynix-samsung-softbank.html}
\bibitem{stargate-reality} Epoch AI, ``OpenAI Stargate: where the US sites stand'' (Apr.\ 2026: $\sim$0.2--0.3GW operational of 10GW pledged); partner-dispute reporting. \url{https://epoch.ai/publications/openai-stargate-where-the-us-sites-stand}

% ============ ANALOGY LIMITS ============
\bibitem{red-ink} G.~DiPippo et al., ``Red Ink: Estimating Chinese Industrial Policy Spending in Comparative Perspective,'' CSIS (2022): $\sim$\$248B/yr, 1.7\% of GDP. \url{https://csis-website-prod.s3.amazonaws.com/s3fs-public/publication/220523_DiPippo_Red_Ink.pdf}
\bibitem{comac} Leeham News, ``COMAC struggled in 2025'' (Jan.\ 2026: $\sim$13 deliveries vs 75 plan); 36Kr on $>$40\% foreign core content and the 2025 LEAP license suspension/restoration. \url{https://leehamnews.com/2026/01/15/comac-struggled-in-2025-2026-wont-be-much-better/}
\bibitem{cj1000a} 36Kr, ``CJ-1000A: verification flight testing 2025, certification targeted end-2025''; IBA assessment: ``years from readiness.'' \url{https://www.iba.aero/resources/articles/comac-aircraft-programmes-status-outlook/}
\bibitem{easa-c919} International Federation of Airworthiness / EASA Director Guillermet: C919 European certification 3--6 years out. \url{https://ifairworthy.com/comacs-c919-will-not-be-certified-in-europe-for-3-to-6-years-according-to-easa-director/}
\bibitem{litho-localization} Tianxia Gongchang Research, ``China semiconductor equipment 2026'' (localization by segment: litho $<$5\%, etch 30--35\%, ArF resist $<$10\%, valves $<$10\%); TrendForce on SMEE 28nm retraction. \url{https://faxiangongchang.com/en/reports/china-semiconductor-equipment-2026}
\bibitem{instruments} 36Kr, ``High-end scientific instruments: domestic share $\sim$17\%, imports \$17B'' (electron microscopes $>$90\% imported). \url{https://eu.36kr.com/en/p/3363026278434564}
\bibitem{machine-tools} USCC MIC2025 evaluation: high-end CNC $\sim$8\% domestic vs 80\% target; robots 52\% vs 70\%; US Commerce China machine-tools market intelligence. \url{https://www.uscc.gov/sites/default/files/2025-11/Made_in_China_2025--Evaluating_Chinas_Performance.pdf}
\bibitem{harmonic} Jamestown Foundation, ``New gains in PRC robotics'' (Harmonic Drive $\sim$85\% global share 2023; Nabtesco $\sim$60\% RV; Leaderdrive 15\% global; Chinese reducers in Tesla Optimus/Unitree). \url{https://jamestown.org/new-gains-in-prc-robotics-software-hardware/}
\bibitem{harmonyos} Counterpoint via Silicon UK: HarmonyOS 17\% of China smartphone OS, ahead of iOS six straight quarters; Caixin on 8M developers; Kingsoft WPS $\sim$600M MAU. \url{https://www.silicon.co.uk/mobility/mobile-os/huawei-apple-china-2-626789}
\bibitem{biotech} Pharmasource, ``China biopharma out-licensing surges to record \$137.7B in 2025''; Caixin on first-in-class pipeline ($\sim$24\%); The Cancer Letter on the ODAC 14--1 sintilimab vote (2022). \url{https://pharmasource.global/content/china-biopharma-out-licensing-surges-to-record-137-7b-in-2025-2026-on-pace-to-break-it-again/}
\bibitem{csis-aviation} CSIS (Tintor\'e Vicent), ``Industrial Policy Without Illusions: What China's EV Boom and Aviation Struggles Teach'' (2025). \url{https://www.csis.org/blogs/trustee-china-hand/industrial-policy-without-illusions-what-chinas-ev-boom-and-aviation}
\bibitem{kyle-chan} K.~Chan, ``China's Technology Long Game'' (High Capacity); USCC testimony (Feb.\ 2025) on recurring five-year-plan sectors and the moving-frontier problem. \url{https://www.highcapacity.org/p/chinas-tech-long-game}
\bibitem{talent} MacroPolo Global AI Talent Tracker 2.0 (Mar.\ 2024; NeurIPS 2019 vs 2022 authors, top-tier $\approx$ top 20\%): Chinese-undergrad share of top-tier researchers 29\%$\rightarrow$47\%; share working in the US 59\%$\rightarrow$42\%, in China 11\%$\rightarrow$28\% (MIT Technology Review, 27 Mar.\ 2024); MacroPolo ceased operations in March 2025 and later ``Tracker 3.0'' figures circulating online are not the Institute's; Carnegie Endowment (Dec.\ 2025): 87 of 100 tracked 2019-cohort Chinese-origin researchers remained in the US; RIETI on the returnee wave and the K-visa (Oct.\ 2025). \url{https://macropolo.org/interactive/digital-projects/the-global-ai-talent-tracker/}

% ============ MEASUREMENT AND ACCOUNTING SOURCES ============
\bibitem{r1-replication} C.~Zhang et al., ``100 Days After DeepSeek-R1: A Survey on Replication Studies...'' arXiv:2505.00551 (2025). \url{https://arxiv.org/abs/2505.00551}
\bibitem{epoch-output} Epoch AI, ``LLM responses to benchmark questions are getting longer over time'' (Apr.\ 2025: reasoning models $\sim$5$\times$/yr; non-reasoning 2.2$\times$/yr). \url{https://epoch.ai/data-insights/output-length}
\bibitem{epoch-prices} Epoch AI, ``LLM inference prices have fallen rapidly but unequally across tasks'' (Mar.\ 2025: 9$\times$--900$\times$/yr, median $\sim$50$\times$). \url{https://epoch.ai/data-insights/llm-inference-price-trends}
\bibitem{epoch-capacity} Epoch AI, ``Global AI computing capacity is doubling every 7 months'' (Jan.\ 2026); ``Acquisition costs of leading AI supercomputers doubled every 13 months'' (Jun.\ 2025). \url{https://epoch.ai/data-insights/ai-chip-production}
\bibitem{vera-cpu} NVIDIA, ``NVIDIA launches Vera CPU'' (Mar.\ 16, 2026: 88 Olympus cores, LPDDR5X, up to 1.2 TB/s, up to 1.5 TB). \url{https://nvidianews.nvidia.com/news/nvidia-launches-vera-cpu-purpose-built-for-agentic-ai}
\bibitem{techinsights-n3} TechInsights, ``SMIC N+3 confirmed: Kirin 9030 analysis reveals how close SMIC is to 5nm'' (Dec.\ 11, 2025; ``significantly less scaled than leading commercial 5nm''). \url{https://www.techinsights.com/blog/smic-n3-confirmed-kirin-9030-analysis-reveals-how-close-smic-5nm}
\bibitem{cxmt-reuters} Reuters (Ngui/Shen), ``Chipmaker CXMT vaults to top of China's valuation with 466\% surge'' (Jul.\ 27, 2026: \$8.6B raised; 6.73\% free float). \url{https://www.cnbc.com/2026/07/27/cxmt-china-market-debut-chipmaker-ipo.html}
\bibitem{trendforce-hbm} TrendForce HBM market bulletin (Jan.\ 2026): Chinese HBM in early verification stages, limited near-term global impact. \url{https://www.trendforce.com/research/download/RP260120JO3}
\bibitem{msft-q3} Microsoft FY2026 Q3 earnings (Apr.\ 29, 2026): AI business $>$\$37B annualized run-rate, +123\% y/y; capex \$31.9B, guided $>$\$40B; capacity constrained through 2026. \url{https://www.microsoft.com/en-us/investor/earnings/fy-2026-q3/press-release-webcast}
\bibitem{boe} Bank of England Financial Stability Report (Jul.\ 2026): AI-focused companies $\approx$ half of S\&P 500 value (from $\sim$a quarter in 2022); Oct.\ 2025 FPC record and Dec.\ 2025 FSR on stretched valuations. \url{https://www.bankofengland.co.uk/financial-stability-report/2026/july-2026}

% ============ ENERGY ============
\bibitem{china-capacity} NEA data via carboncredits.com: China +540 GW in 2025 (total 3.89 TW; $\sim$\$500B invested); 4.04 TW and 1,551 GW record peak by mid-2026 (Global Times); falling utilization hours. \url{https://carboncredits.com/china-adds-power-7x-more-than-the-us-in-2025-with-500b-energy-build-out-in-a-single-year/}
\bibitem{us-additions} US EIA: 2025 utility-scale additions 53 GW (largest since 2002); 2026 planned record 86 GW. \url{https://www.eia.gov/todayinenergy/detail.php?id=67205}
\bibitem{ember-2026} Ember, Global Electricity Review 2026: China 10,573 TWh (+503 TWh y/y) vs US 4,430 TWh; China fossil generation fell 0.9\% in 2025; solar +40\%. \url{https://ember-energy.org/latest-insights/global-electricity-review-2026/major-countries-and-regions/}
\bibitem{nuclear-china} World Nuclear News / POWER: 10 reactors approved Apr.\ 2025 (\$27.4B), 8 more Jul.\ 2026; 66 GW installed, targets 110 GW by 2030. \url{https://www.world-nuclear-news.org/articles/ten-new-reactors-approved-in-china}
\bibitem{coal-2025} Carbon Brief/CREA-GEM: record $\sim$95 GW coal commissioned 2025; permits at four-year low; coal repositioned as flexibility. \url{https://www.carbonbrief.org/rush-for-new-coal-in-china-hits-record-high-in-2025-as-climate-deadline-looms}
\bibitem{uhv} Global Energy Monitor / OIES: 45 UHV lines in operation (Dec.\ 2025); four new 8-GW UHVDC corridors in 2025. \url{https://www.gem.wiki/Ultra-High-Voltage_(UHV)_Power_Transmission_System_in_China}
\bibitem{us-queue} LBNL ``Queued Up'' 2026: $\sim$2,061 GW active queue; median $>$3 yrs to agreement; 13\% historical completion. \url{https://www.publicpower.org/periodical/article/backlog-power-plants-seeking-transmission-grid-connection-eased-somewhat-2025-lbnl}
\bibitem{sp-dc} S\&P Global 451 Research (Oct.\ 2025): US DC grid demand 61.8 GW (2025) $\rightarrow$ 134.4 GW (2030); LBNL 2024 report: DCs 176 TWh (4.4\%) in 2023 $\rightarrow$ 6.7--12\% by 2028. \url{https://www.spglobal.com/energy/en/news-research/latest-news/electric-power/101425-data-center-grid-power-demand-to-rise-22-in-2025-nearly-triple-by-2030}
\bibitem{grid-strategies} Grid Strategies, 2025 National Load Growth Report: 166 GW five-year peak-load growth forecast, $\sim$90 GW from data centers. \url{https://gridstrategiesllc.com/wp-content/uploads/Grid-Strategies-National-Load-Growth-Report-2025.pdf}
\bibitem{pjm} Canary Media: PJM record \$16.4B capacity auction (2027/28); $\sim$94\% of load growth from DCs; \$70/mo household impact estimates; Good Jobs First: \$4.3B ratepayer-funded DC interconnection (2024). \url{https://www.canarymedia.com/articles/data-centers/pjm-record-capacity-costs-rising-bills}
\bibitem{moratoria} Good Jobs First (Mar.\ 2026): data-center moratorium bills in $\geq$12 states in 2026 sessions. \url{https://goodjobsfirst.org/data-center-moratorium-bills-are-spreading-in-2026/}
\bibitem{gas-turbines} Utility Dive: GE Vernova gas-turbine backlog 116 GW (Q2 2026), reservations for 2031 delivery. \url{https://www.utilitydive.com/news/ge-vernova-gas-turbine-backlog-climbs-to-116-gw/826039/}
\bibitem{transformers} Smartland Energy (2026, citing Wood Mackenzie): large-transformer lead times $>$2 yrs, prices +70\% since 2019. \url{https://smartlandenergy.com/transformer-supply-in-2026-a-constraint-thats-reshaping-grid-development/}
\bibitem{tmi-palisades} WITF (Jul.\ 2026): Three Mile Island/Crane restart on track for 2027; Michigan Public/ANS on Palisades slippage. \url{https://www.witf.org/2026/07/28/three-mile-island-reactor-restart-progresses-through-environmental-review-public-comments/}
\bibitem{openai-100gw} OpenAI letter to OSTP (Oct.\ 27, 2025): US needs 100 GW/yr of new capacity; the ``electron gap'' vs China's 429 GW. \url{https://www.datacenterdynamics.com/en/news/openai-tells-white-house-that-us-needs-100gw-of-new-energy-per-year-to-support-ai-boom/}
\bibitem{edwc} National Data Bureau via TechRadar: East-Data-West-Computing eight-hub program; PUE 1.04 best-in-class; 80\%-renewable requirement for new hub DCs. \url{https://www.techradar.com/pro/website-hosting/china-invested-dollar61-billion-in-eight-national-datacenter-hubs-according-to-director-of-national-data-bureau}
\bibitem{china-dc-power} Semafor/techchannel (Nov.\ 2025): provinces cut DC electricity up to 50\% (to $\sim$RMB 0.4/kWh) exclusively for approved-domestic-chip facilities; offsets 30--50\% chip-efficiency gap. \url{https://www.semafor.com/article/11/05/2025/china-slashes-big-data-centers-electric-bills}
\bibitem{oies} Oxford Institute for Energy Studies, ``The China data centre advantage'' (Feb.\ 2026): TCO decomposition; spot-price overlap; headroom argument. \url{https://www.oxfordenergy.org/wpcms/wp-content/uploads/2026/02/Comment-The-China-data-centre-advantage.pdf}
\bibitem{iea-energy-ai} IEA, ``Energy and AI'' (2025): global DC consumption 415 TWh (2024) $\rightarrow$ $\sim$945 TWh (2030). \url{https://www.iea.org/reports/energy-and-ai/executive-summary}

% ============ TRUST & SECURITY ============
\bibitem{pickle} JFrog (2024: $\sim$100 malicious HF models); ReversingLabs ``nullifAI'' (Feb.\ 2025: broken-pickle evasion); Sonatype Picklescan CVEs (2025). \url{https://www.reversinglabs.com/blog/rl-identifies-malware-ml-model-hosted-on-hugging-face}
\bibitem{namespace} Palo Alto Unit 42, ``Model namespace reuse'' (Sep.\ 2025): trusted-name takeover propagating into Vertex AI and Azure AI Foundry. \url{https://unit42.paloaltonetworks.com/model-namespace-reuse/}
\bibitem{hf-incident} Hugging Face, ``Security incident disclosure --- July 2026'' (Jul.\ 16, 2026: dataset-processing RCE, credential harvest, lateral movement; agentic attacker; no public-artifact tampering found). \url{https://huggingface.co/blog/security-incident-july-2026}
\bibitem{poison} Anthropic / UK AISI / Turing Institute (Oct.\ 2025): $\sim$250 poisoned documents backdoor models regardless of scale. \url{https://www.anthropic.com/research/small-samples-poison}
\bibitem{sleeper} Hubinger et al., ``Sleeper Agents'' (arXiv:2401.05566): backdoors survive SFT/RLHF/adversarial training; agent-supply-chain extensions (CAIS 2026). \url{https://arxiv.org/abs/2401.05566}
\bibitem{virustotal} Hugging Face and VirusTotal collaboration (Oct.\ 22, 2025): continuous scanning of 2.2M+ public repositories. \url{https://huggingface.co/blog/virustotal}
\bibitem{deepseek-censor} Qiu, Zhou, Ferrara, ``Information Suppression in Large Language Models: ...Censorship in DeepSeek'' (arXiv:2506.12349): semantic-level suppression across 646 sensitive prompts. \url{https://arxiv.org/abs/2506.12349}
\bibitem{r1dacted} Naseh et al., ``R1dacted: Investigating Local Censorship in DeepSeek's R1'' (arXiv:2505.12625): topic/language-dependent censorship, transfers to distilled models. \url{https://arxiv.org/abs/2505.12625}
\bibitem{anthropic-distill} Anthropic, ``Detecting and preventing distillation attacks'' (Feb.\ 23, 2026): $>$16M exchanges via $\sim$24,000 accounts (MiniMax $>$13M, Moonshot $>$3.4M, DeepSeek $>$150K). Vendor allegation. \url{https://www.anthropic.com/news/detecting-and-preventing-distillation-attacks}

% ============ AGENTIC SECURITY ============
\bibitem{order66} S.~Matsuoka, ``Latent Model Backdoors and Prompt-Borne Worms: A Mini-Survey of Evidence, the `Order 66' Compound Scenario, and Defenses,'' v6.0, 31 July 2026. Companion survey; supplies the evidence hierarchy, threat taxonomy, and control architecture used in Chapter~\ref{ch:order66} and Appendix~\ref{app:controls}.
\bibitem{poison-constant} A.~Souly, J.~Rando, E.~Chapman et al., ``Poisoning Attacks on LLMs Require a Near-Constant Number of Poison Samples,'' arXiv:2510.07192 (2025). \url{https://arxiv.org/abs/2510.07192}
\bibitem{badagent} Y.~Wang, D.~Xue, S.~Zhang, and S.~Qian, ``BadAgent: Inserting and Activating Backdoor Attacks in LLM Agents,'' ACL 2024, pp.~9811--9827. \url{https://doi.org/10.18653/v1/2024.acl-long.530}
\bibitem{backreveal} W.~Zhang and S.~Pei, ``Your LLM Agent Can Leak Your Data: Data Exfiltration via Backdoored Tool Use,'' Findings of ACL 2026, pp.~25105--25129. \url{https://aclanthology.org/2026.findings-acl.1257.pdf}
\bibitem{trojanstego} D.~Meier, J.~P.~Wahle, P.~R\"ottger, T.~Ruas, and B.~Gipp, ``TrojanStego: Your Language Model Can Secretly Be a Steganographic Privacy Leaking Agent,'' EMNLP 2025, pp.~27244--27261. \url{https://doi.org/10.18653/v1/2025.emnlp-main.1386}
\bibitem{morris2} S.~Cohen, R.~Bitton, and B.~Nassi, ``Here Comes The AI Worm: Unleashing Zero-click Worms that Target GenAI-Powered Applications,'' arXiv:2403.02817 (rev.\ 2025). \url{https://arxiv.org/abs/2403.02817}
\bibitem{agentworm} Y.~Zhang, Z.~Wei, X.~Luan et al., ``AgentWorm: Self-Propagating Attacks Across LLM Agent Ecosystems,'' arXiv:2603.15727 v3 (Jul.\ 2026). \url{https://arxiv.org/abs/2603.15727}
\bibitem{adaptive-worm} J.~Guan, T.~Blanchard, H.~Foerster, H.~Jia, G.~Huang, and N.~Papernot, ``AI Agents Enable Adaptive Computer Worms,'' arXiv:2606.03811 (Jun.\ 2026). \url{https://arxiv.org/abs/2606.03811}
\bibitem{malicious-skills} Y.~Liu, Z.~Chen, Y.~Zhang et al., ``Malicious Agent Skills in the Wild: A Large-Scale Security Empirical Study,'' arXiv:2602.06547 (2026). \url{https://arxiv.org/abs/2602.06547}
\bibitem{haystack} B.~Bullwinkel, G.~Severi, K.~Hines, A.~Minnich, R.~S.~S.~Kumar, and Y.~Zunger, ``The Trigger in the Haystack: Extracting and Reconstructing LLM Backdoor Triggers,'' arXiv:2602.03085 (2026). \url{https://aka.ms/airt-backdoor-detection}
\bibitem{caisi} Center for AI Standards and Innovation, NIST, \emph{Evaluation of DeepSeek AI Models} (Sept.\ 2025): agent-configuration susceptibility to indirect prompt injection, measured as attempted malicious actions. \url{https://www.nist.gov/system/files/documents/2025/09/30/CAISI_Evaluation_of_DeepSeek_AI_Models.pdf}
\bibitem{openai-hf} OpenAI, ``OpenAI and Hugging Face Partner to Address Security Incident During Model Evaluation'' (21 Jul.\ 2026). \url{https://openai.com/index/hugging-face-model-evaluation-security-incident/}
\bibitem{hf-timeline} Hugging Face Security Team, ``Anatomy of a Frontier Lab Agent Intrusion: A Technical Timeline of the July 2026 Incident'' (27 Jul.\ 2026). \url{https://huggingface.co/blog/agent-intrusion-technical-timeline}
\bibitem{anthropic-incidents} Anthropic, ``Investigating Three Real-World Incidents in Our Cybersecurity Evaluations'' (30 Jul.\ 2026). \url{https://www.anthropic.com/news/investigating-incidents-cybersecurity-evals}

% ============ INTERNATIONAL CONSORTIUM ============
\bibitem{ai4s-framework} S.~Matsuoka, ``How Much Compute Does a Nation's Science Need? A Portable, Measurement-Anchored Framework for Sizing Sovereign AI-for-Science Infrastructure, with Instantiations for Japan, Singapore, and the United States,'' RIKEN R-CCS Technical Report, July 2026. Supplies the BGPU unit, the two-layer demand model, the parameter taxonomy and transfer rules, the measurement protocol, and the three national instantiations used in Chapter~\ref{ch:sizing}.
\bibitem{consortium-record} S.~Matsuoka and T.~Schulthess, ``AI for Science and an International Consortium,'' edited discussion record v1.0, 31 July 2026 (RIKEN R-CCS / CSCS). Unpublished working document; numerical values are discussion estimates unless otherwise noted.

% ============ AI4S: THE LARGEST PRIZE ============
\bibitem{ecmwf-aifs} ECMWF, ``ECMWF's AI forecasts become operational'' (25 Feb.\ 2025) and ``Ensemble AI forecasts become operational'' (1 Jul.\ 2025). \url{https://www.ecmwf.int/en/about/media-centre/news/2025/ecmwfs-ai-forecasts-become-operational}
\bibitem{noaa-ai} NOAA, ``NOAA deploys new generation of AI-driven global weather models'' (17 Dec.\ 2025: AIGFS at 0.3\% of GFS compute; AIGEFS 31 members at 9\%). \url{https://www.noaa.gov/news-release/noaa-deploys-new-generation-of-ai-driven-global-weather-models}
\bibitem{aurora} C.~Bodnar et al., ``A foundation model for the Earth system,'' \emph{Nature} 641 (21 May 2025). \url{https://www.nature.com/articles/s41586-025-09005-y}
\bibitem{neuralgcm} D.~Kochkov et al., ``Neural general circulation models for weather and climate,'' \emph{Nature} (22 Jul.\ 2024). \url{https://research.google/blog/fast-accurate-climate-modeling-with-neuralgcm/}
\bibitem{cma-ai} China Meteorological Administration, operational AI models FengQing, FengLei, FengShun (22 Dec.\ 2025). \url{https://www.cma.gov.cn/en/research/news/202512/t20251222_7506498.html}
\bibitem{fuxi-cnops} FuXi-CNOPs typhoon-track system, IAP-CAS and Fudan University, \emph{Advances in Atmospheric Sciences} (2026): up to 32.3\% track-error reduction. \url{https://www.globaltimes.cn/page/202607/1367074.shtml}
\bibitem{alphafold-impact} Google DeepMind, ``AlphaFold: five years of impact'' (2025). \url{https://deepmind.google/blog/alphafold-five-years-of-impact/}
\bibitem{nobel-2024} Nobel Prize in Chemistry 2024 (Hassabis, Jumper, Baker), 9 Oct.\ 2024. \url{https://www.nobelprize.org/prizes/chemistry/2024/press-release/}
\bibitem{alphafold-db} EMBL-EBI, first complexes in the AlphaFold Database (Mar.--May 2026; $\sim$17M GPU-hours). \url{https://www.ebi.ac.uk/about/news/technology-and-innovation/first-complexes-alphafold-database/}
\bibitem{esm3} T.~Hayes et al., ``Simulating 500 million years of evolution with a language model,'' \emph{Science} (Jan.\ 2025). \url{https://www.science.org/doi/10.1126/science.ads0018}
\bibitem{rfdiffusion3} Baker Lab, RFdiffusion3 (bioRxiv, Dec.\ 2025); designed metallohydrolases with catalytic efficiency orders of magnitude above prior designs. \url{https://www.genengnews.com/topics/artificial-intelligence/rfdiffusion3-now-open-source-designs-dna-binders-and-advanced-enzymes/}
\bibitem{imo-gold} Google DeepMind, ``Gemini Deep Think achieves gold-medal standard at IMO 2025'' (21 Jul.\ 2025: 35/42, 5 of 6 problems). \url{https://deepmind.google/blog/advanced-version-of-gemini-with-deep-think-officially-achieves-gold-medal-standard-at-the-international-mathematical-olympiad/}
\bibitem{alphaevolve} Google DeepMind, ``AlphaEvolve'' (May 2025): 48-multiplication 4$\times$4 complex matmul; 0.7\% of worldwide compute recovered; 23\% kernel speedup $\rightarrow$ 1\% Gemini training-time reduction. \url{https://deepmind.google/blog/alphaevolve-a-gemini-powered-coding-agent-for-designing-advanced-algorithms/}
\bibitem{rentosertib} Insilico Medicine, rentosertib Phase IIa, \emph{Nature Medicine} (3 Jun.\ 2025). \url{https://insilico.com/news/tnrecuxsc1-insilico-announces-nature-medicine-publi}
\bibitem{rentosertib-p3} Insilico Medicine, Phase III initiation (7 Jul.\ 2026; 320 patients, 47 centres). \url{https://insilico.com/news/xmjsn4l091-insilico-initiates-phase-iii-clinical-tr}
\bibitem{mit-repudiation} MIT Department of Economics, statement on ``Artificial Intelligence, Scientific Discovery, and Product Innovation'' (16 May 2025): no confidence in the data; withdrawn. \url{https://economics.mit.edu/news/assuring-accurate-research-record}
\bibitem{alab-correction} Author Correction to the A-Lab autonomous-synthesis paper, \emph{Nature} (19 Jan.\ 2026): successes 40 $\rightarrow$ 36. \url{https://www.nature.com/articles/s41586-025-09992-y}
\bibitem{alab-critics} Chemistry World, ``New analysis raises doubts over autonomous lab's materials discoveries.'' \url{https://www.chemistryworld.com/news/new-analysis-raises-doubts-over-autonomous-labs-materials-discoveries/4018791.article}
\bibitem{gnome-critique} A.~K.~Cheetham and R.~Seshadri, critique of AI-generated inorganic materials; coverage in The Register (11 Apr.\ 2024). \url{https://www.theregister.com/2024/04/11/google_deepmind_material_study/}
\bibitem{mattergen} C.~Zeni et al., ``A generative model for inorganic materials design,'' \emph{Nature} 639 (16 Jan.\ 2025). \url{https://www.nature.com/articles/s41586-025-08628-5}
\bibitem{ecmwf-farewell} ECMWF AIFS blog, ``Farewell to external AI models'' (May 2026): all four external models withdrawn from real-time running over analysis sensitivity. \url{https://www.ecmwf.int/en/about/media-centre/aifs-blog/2026/farewell-external-ai-models}
\bibitem{aletheia} Feng, Luong et al., adjudication of an autonomous agent on 700 open Erd\H{o}s problems, arXiv:2601.22401 (2026): 68.5\% of adjudicated candidates fundamentally flawed. \url{https://arxiv.org/html/2601.22401v3}
\bibitem{tao} T.~Tao on AI solutions to Erd\H{o}s problems (Jan.\ 2026): 1--2\% tractable; unreported failures skew perceived rates. \url{https://the-decoder.com/terence-tao-says-gpt-5-2-pro-cracked-an-erdos-problem-but-warns-the-win-says-more-about-speed-than-difficulty/}
\bibitem{agents4science} Agents4Science 2025 (Stanford, Oct.\ 2025): 314 submissions, 48 accepted; reviewer assessments. \url{https://www.sciencenews.org/article/science-conference-test-ai-agents}
\bibitem{openai-science} MIT Technology Review, ``Inside OpenAI's big play for science'' (26 Jan.\ 2026), incl.\ working scientists' assessments. \url{https://www.technologyreview.com/2026/01/26/1131728/inside-openais-big-play-for-science/}
\bibitem{ai-slop} \emph{Organization Science} pipeline study (Apr.\ 2026): submissions $+$42\%, $>$30\% of reviews showing AI use, 3.2\% vs 11.9\% R\&R rates. \url{https://www.forbes.com/sites/johndrake/2026/04/30/ai-slop-is-flooding-academic-journals-a-top-journal-measured-it/}
\bibitem{nsf-indicators} NSF NCSES, \emph{Science and Engineering Indicators 2025}: global R\&D \$3.1T (PPP, 2022); US \$923B, China \$812B. \url{https://ncses.nsf.gov/pubs/nsb20257/global-r-d-and-international-comparisons-2}
\bibitem{gartner} Gartner, worldwide IT spending forecast (27 Jul.\ 2026): software \$1.468T in 2026. \url{https://www.gartner.com/en/newsroom/press-releases/2026-07-27-gartner-forecasts-worldwide-it-spending-to-grow-14-point-2-percent-in-2026-totaling-6-point-37-trillion}
\bibitem{lab-revenue} Forbes (21 May 2026) on frontier-laboratory revenue scale and composition. \url{https://www.forbes.com/sites/paulocarvao/2026/05/21/anthropic-openai-enterprise-ai-profitability/}
\bibitem{mckinsey-genai} McKinsey Global Institute, ``The economic potential of generative AI'' (2023): R\&D the largest affected function ($\sim$24\%); estimates based on the 2022 economy and excluding new product categories. \url{https://www.mckinsey.com/capabilities/tech-and-ai/our-insights/the-economic-potential-of-generative-ai-the-next-productivity-frontier}
\bibitem{jones-summers} B.~F.~Jones and L.~H.~Summers, ``A Calculation of the Social Returns to Innovation,'' NBER w27863: \$13.3 per \$1 at a 5\% discount rate. \url{https://www.nber.org/system/files/working_papers/w27863/w27863.pdf}
\bibitem{ideas-harder} N.~Bloom, C.~I.~Jones, J.~Van Reenen and M.~Webb, ``Are Ideas Getting Harder to Find?'' \emph{AER} 110(4), 2020. \url{https://web.stanford.edu/~chadj/IdeaPF.pdf}
\bibitem{acemoglu} D.~Acemoglu, ``The Simple Macroeconomics of AI,'' NBER w32487 (2024): TFP $+$0.66\% over ten years, \emph{explicitly excluding} protein-folding-class scientific advances as unlikely within the window. \url{https://economics.mit.edu/sites/default/files/2024-05/The\%20Simple\%20Macroeconomics\%20of\%20AI.pdf}
\bibitem{davidson} T.~Davidson, D.~Halperin, N.~Houlden and A.~Korinek, ``When Does Automating AI Research Produce Explosive Growth?'' NBER w35155 (Apr.\ 2026). \url{https://www.nber.org/system/files/working_papers/w35155/w35155.pdf}
\bibitem{quadrillion} AI Frontiers, ``The quadrillion-dollar disagreement on AI and the economy'' (11 May 2026): credible estimates span 0.1\%--30\% annual growth contribution. \url{https://ai-frontiers.org/articles/the-quadrillion-dollar-disagreement-on-ai-and-the-economy}
\bibitem{nature-index-2026} Nature Index 2026 Research Leaders (10 Jun.\ 2026): China first, output $+$22.4\% y/y; CAS the leading institution. \url{https://www.natureasia.com/en/info/press-releases/detail/9350}
\bibitem{nature-index-2025} Nature Index 2025 Research Leaders (11 Jun.\ 2025): China Share 32,122 ($+$17\%), eight of the top ten institutions. \url{https://group.springernature.com/gp/group/media/press-releases/nature-index-research-leaders-/27786652}
\bibitem{china-ai4s} UC Institute on Global Conflict and Cooperation, ``Inside China's AI for Science strategy'' (MOST/NSFC initiative 2023; AI and fusion in the 15th Five-Year Plan). \url{https://ucigcc.org/blog/inside-chinas-ai-for-science-strategy/}
\bibitem{scienceone} Chinese Academy of Sciences, ScienceOne launch (28 Jul.\ 2025): 12 institutes, 170M papers, 300+ tools. \url{https://english.cas.cn/newsroom/cas_media/202507/t20250728_1048587.shtml}
\bibitem{scienceone2} Chinese Academy of Sciences, ScienceOne Omni 2.0 (20 Jul.\ 2026). \url{https://english.cas.cn/newsroom/cas-in-media/202607/t20260720_1178699.shtml}
\bibitem{scnet} Reporting on an autonomous AI research system connected to China's national supercomputing network SCNet (Dec.\ 2025); secondary source, primary announcement not located. \url{https://tech.yahoo.com/ai/articles/china-rolls-super-ai-science-081459774.html}
\bibitem{dp-tech} DP Technology Series C (24 Dec.\ 2025); Bohrium and related platforms serving 1,000+ research organizations and 150 corporate clients. \url{https://siliconangle.com/2025/12/24/dp-technology-raises-114m-accelerate-chinas-ai-science-industry/}
\bibitem{intern-s1} Shanghai AI Laboratory, Intern-S1-Pro (2 Apr.\ 2026): trillion-parameter open-weight scientific multimodal model; SciReasoner 55.5 vs 14.7/13.6 for leading closed models. arXiv:2603.25040. \url{https://arxiv.org/html/2603.25040v2}
\bibitem{genesis} US DOE, ``Energy Department launches Genesis Mission'' (Executive Order, 24 Nov.\ 2025): goal to double the productivity of American science within a decade. \url{https://www.energy.gov/articles/energy-department-launches-genesis-mission-transform-american-science-and-innovation}
\bibitem{genesis-scale} NVIDIA/Oracle/DOE, Solstice (100,000 Blackwell GPUs) and Equinox (10,000), 28 Oct.\ 2025. \url{https://nvidianews.nvidia.com/news/nvidia-oracle-us-department-of-energy-ai-supercomputer-scientific-discovery}
\bibitem{genesis-projects} The White House, Genesis Mission commitments (22 Jul.\ 2026): $>$\$5B, 278 projects, 15+ agencies. \url{https://www.whitehouse.gov/releases/2026/07/45502/}
\bibitem{eu-factories} European Commission, AI Factories (19 operational plus 13 antennas; €10B via EuroHPC JU). \url{https://digital-strategy.ec.europa.eu/en/policies/ai-factories}
\bibitem{eu-gigafactories} European Commission, InvestAI / AI Gigafactories (€20B; 77 proposals across 16 member states). \url{https://commission.europa.eu/topics/competitiveness/competitiveness-coordination-tool-projects/ai-gigafactories_en}
\bibitem{isambard} University of Bristol, Isambard-AI launch (17 Jul.\ 2025): £225M, 5,448 GH200 superchips. \url{https://www.bristol.ac.uk/news/2025/july/isambard-launch.html}
\bibitem{fugakunext} RIKEN, FugakuNEXT announcement (22 Aug.\ 2025): $>$600 EF FP8 target within $\sim$40\,MW; explicit AI-for-Science mandate. \url{https://www.riken.jp/en/news_pubs/news/2025/20250822_1/index.html}
\bibitem{nea-2025} China National Energy Administration, 2025 installed-capacity and renewable-integration statistics. \url{https://www.nea.gov.cn/20260129/6874f211acd0417eab7ac10c3061a7c2/c.html}

% ============ THEORY ============
\bibitem{bresnahan1995} T.~F.~Bresnahan and M.~Trajtenberg, ``General Purpose Technologies: Engines of Growth?'' \emph{Journal of Econometrics} 65(1):83--108, 1995.
\bibitem{helpman1998} E.~Helpman (ed.), \emph{General Purpose Technologies and Economic Growth}, MIT Press, 1998.
\bibitem{brynjolfsson2021} E.~Brynjolfsson, D.~Rock and C.~Syverson, ``The Productivity J-Curve: How Intangibles Complement General Purpose Technologies,'' \emph{AEJ: Macroeconomics} 13(1):333--372, 2021.
\bibitem{schumpeter1942} J.~A.~Schumpeter, \emph{Capitalism, Socialism and Democracy}, Harper \& Brothers, 1942.
\bibitem{porter1980} M.~E.~Porter, \emph{Competitive Strategy}, Free Press, 1980.
\bibitem{christensen1997} C.~M.~Christensen, \emph{The Innovator's Dilemma}, Harvard Business School Press, 1997.
\bibitem{teece1997} D.~J.~Teece, G.~Pisano and A.~Shuen, ``Dynamic Capabilities and Strategic Management,'' \emph{Strategic Management Journal} 18(7):509--533, 1997.
\bibitem{baldwin2000} C.~Y.~Baldwin and K.~B.~Clark, \emph{Design Rules, Volume 1: The Power of Modularity}, MIT Press, 2000.
\bibitem{jacobides2018} M.~G.~Jacobides, C.~Cennamo and A.~Gawer, ``Towards a Theory of Ecosystems,'' \emph{Strategic Management Journal} 39(8):2255--2276, 2018.
\bibitem{roofline2009} S.~Williams, A.~Waterman and D.~Patterson, ``Roofline: An Insightful Visual Performance Model for Multicore Architectures,'' \emph{Communications of the ACM} 52(4):65--76, 2009.

% ============ ADDITIONS ============
\bibitem{archer2-eval} UK Research and Innovation, \emph{ARCHER2 High Performance Computer Evaluation} (Nov.\ 2025 / Mar.\ 2026): independent evaluation with explicit counterfactual scenarios and benefit-cost accounting against public investment. \url{https://www.ukri.org/publications/archer2-high-performance-computer-evaluation-nov-2025/}
\bibitem{anthropic-margin} PitchBook via Yahoo Finance, ``Anthropic's gross margin is the most important number in AI'' (10 Jun.\ 2026): \$0.71 compute cost per revenue dollar in Q1 2026, projected \$0.56 in Q2, trajectory toward $\sim$44\%. \url{https://finance.yahoo.com/markets/stocks/articles/anthropics-gross-margin-most-important-070500338.html}
\bibitem{lbl-queued} Lawrence Berkeley National Laboratory, \emph{Queued Up: 2025 Edition} (data through 2024): $\sim$2,290\,GW of active generation and storage requests; median request-to-operation time above four years; only 13\% of 2000--2019 requested capacity ever reached operation. \url{https://emp.lbl.gov/publications/queued-2025-edition-characteristics}
\bibitem{lbl-speed} Lawrence Berkeley National Laboratory, \emph{Speed to Power: Solutions for Accelerating Large-Load Connections} (Jun.\ 2026): no comprehensive national median exists for large-load connection times; utility examples show 6--18-month load studies and 18--60-month construction with line extension; identifies provisional interconnection and curtailable ``non-firm'' load service among thirteen solution areas. \url{https://eta-publications.lbl.gov/sites/default/files/2026-06/lbnl_large_loads_speed_to_power_final_1.pdf}
\bibitem{dominion-wait} Bloomberg / Construction Owners 2026 market guide: Northern Virginia transmission-level data-center hookups quoted at up to $\sim$7 years; Dominion substation-service queue $>$36 months against 18--24-month shell construction. \url{https://www.constructionowners.com/insights/how-long-it-actually-takes-to-power-a-data-center-in-2026-a-u-s-market-by-market-reality-check}
\bibitem{pjm-2829} PJM Interconnection, 2028/29 Base Residual Auction results (14 Jul.\ 2026): \$325/MW-day, 138,318\,MW procured, $\sim$\$16.4B, 6,831\,MW short of the reliability requirement. \url{https://www.pjm.com/-/media/DotCom/about-pjm/newsroom/2026-releases/20260714-pjm-capacity-auction-procures-138318-mw-of-generation-resources.pdf}
\bibitem{ercot-queue} ERCOT large-load interconnection data and reporting: large-load queue above 230\,GW by Dec.\ 2025 ($\sim$4$\times$ a year earlier), $\sim$7.5\,GW connected; Texas SB6 (Jun.\ 2025) imposes study fees, duplicate-request disclosure, security deposits, and remote-curtailment capability on $\ge$75\,MW loads. \url{https://www.utilitydive.com/news/ercots-large-load-queue-jumped-almost-300-last-year-official/808820/}
\bibitem{doc136} Carbon Brief, analyses of NDRC Document 136 renewable-pricing reform and 2025--26 curtailment: 2025 national wind/solar curtailment 5.4\% (breaching the 5\% target); Q1 2026 capacity-factor declines concentrated in Inner Mongolia, Xinjiang, Liaoning; only 18 provinces had finalized Doc-136 rules by Oct.\ 2025. \url{https://www.carbonbrief.org/analysis-chinas-co2-climbs-2-in-early-2026-due-to-wasted-wind-and-solar}
\bibitem{oies-china-dc} Oxford Institute for Energy Studies, ``The China data centre advantage'' (Feb.\ 2026): interior spot power 300--400 RMB/MWh (\$43--58); coastal DC power \$0.08--0.11/kWh; 39 UHV lines in service, $\sim$28 more planned by 2030; China's risk characterized as generation overcapacity, not shortage. \url{https://www.oxfordenergy.org/wpcms/wp-content/uploads/2026/02/Comment-The-China-data-centre-advantage.pdf}
\bibitem{hub-tariffs} People's Daily / China Daily / Global Times reporting on western hub economics: Qingyang (Gansu) on-site power stabilized at 0.398 RMB/kWh with a dedicated 2\,GW green plant; Ulanqab/Hohhot (Inner Mongolia) data-center band 0.32--0.38 RMB/kWh with $>$80\% green supply and 200+ free-cooling days; Gui'an (Guizhou) Q1 2025 data-center electricity use $+$517\% y/y with dedicated substation buildout. \url{https://en.people.cn/n3/2026/0708/c98649-20475664.html}

\bibitem{mlperf-v6} MLCommons, MLPerf Training v6.0 results (16 Jun.\ 2026): adds DeepSeek-V3 (671B/37B MoE) pretraining and GPT-OSS 20B benchmarks; 95 systems from 24 organizations on 13 accelerator types; no Chinese-designed accelerator submissions. \url{https://mlcommons.org/2026/06/mlperf-training-v6-0-results/}
\bibitem{openai-reset} CNBC, ``OpenAI resets spend expectations, targets around \$600 billion by 2030'' (20 Feb.\ 2026): compute-spending target walked back from headline \$1.4T commitments; $\sim$\$280B revenue projected by 2030. \url{https://www.cnbc.com/2026/02/20/openai-resets-spend-expectations-targets-around-600-billion-by-2030.html}
\bibitem{openai-msft} DCD / Forbes reporting: for-profit conversion completed Oct.\ 2025 (Microsoft $\sim$27\%, \$250B Azure commitment, cloud right of first refusal surrendered); exclusivity and revenue sharing reported ended Apr.\ 2026. \url{https://www.datacenterdynamics.com/en/news/openai-completes-for-profit-move-microsoft-given-27-stake-and-250bn-azure-contract-but-no-longer-has-cloud-right-of-first-refusal/}
\bibitem{anthropic-msft-nvda} Microsoft, ``Microsoft, NVIDIA and Anthropic announce strategic partnerships'' (18 Nov.\ 2025): Microsoft up to \$5B and NVIDIA up to \$10B into Anthropic; Anthropic commits \$30B of Azure purchases. \url{https://blogs.microsoft.com/blog/2025/11/18/microsoft-nvidia-and-anthropic-announce-strategic-partnerships/}
\bibitem{nvidia-gtc} TechCrunch/CNBC, GTC 2026 (16 Mar.\ 2026): Blackwell+Rubin order visibility raised to $\sim$\$1T through 2027; Vera Rubin production ramp H2 2026. \url{https://techcrunch.com/2026/03/16/jensen-just-put-nvidias-blackwell-and-vera-rubin-sales-projections-into-the-1-trillion-stratosphere/}
\bibitem{nvidia-openai-stall} WSJ via Yahoo Finance, ``Stalled: the NVIDIA--OpenAI megadeal'' (31 Jan.\ 2026): the \$100B progressive, non-binding LOI reported stalled amid internal doubts over size and structure. \url{https://finance.yahoo.com/news/stalled-nvidia-openai-megadeal-ai-131959187.html}
\bibitem{softbank-fy26} CNBC, SoftBank FY ended Mar.\ 2026: Vision Fund gains $\sim$\$46B ($>$\$45B from OpenAI); ¥5T annual profit; S\&P outlook negative; $>$\$60B toward $\sim$13\% of OpenAI; \$6.5B Ampere acquisition completed Nov.\ 2025. \url{https://www.cnbc.com/2026/05/13/softbank-earnings-fy-2025-vision-fund-openai-stake.html}
\bibitem{son-asi} CNBC / Fortune (Jun.--Jul.\ 2026): Son declares AI-designing-AI a sign of superintelligence; claims \$5T/yr of global AI investment needed; dismisses bubble concerns; ¥1,000T NAV target. \url{https://fortune.com/2026/07/14/softbank-ceo-5-trillion-ai-demand-critics-bubble/}
\bibitem{arm-fy26} Arm, Q4 FYE26 results: record FY revenue \$4.92B; record royalties \$2.61B; data-center royalties more than doubled; claimed $\sim$50\% CPU share among top hyperscalers; ``AGI'' datacenter CPU with $>$\$2B claimed demand. \url{https://newsroom.arm.com/news/arm-q4-fye26-results}
\bibitem{hyperscaler-capex} Silicon Analysts / trackers: combined 2026 hyperscaler capex guidance $\sim$\$700--745B (AMZN $\sim$\$200B, MSFT $\sim$\$190B, GOOGL \$175--185B, META raised twice toward \$145B); trailing-4Q D\&A $\sim$\$149B; useful-life history incl.\ Amazon's partial reversal to 5 years with $\sim$\$1.3B impact; Alphabet 100-year sterling note; Google Cloud backlog $\sim$\$460B. \url{https://siliconanalysts.com/analysis/hyperscaler-ai-capex-depreciation-wall-2026}
\bibitem{private-credit} Forbes, ``Bond investors push back as AI debt heads toward \$570 billion'' (17 Jul.\ 2026): $\sim$\$236B priced by May 31 at 4$\times$ prior-year pace; $\sim$\$800B private-credit data-center opportunity through 2028; new-issue demand coverage fell from $\sim$5$\times$ (Feb.) to $<$2$\times$ (Jul.). \url{https://www.forbes.com/sites/robertszczerba/2026/07/17/bond-investors-push-back-as-ai-debt-heads-toward-570-billion/}
\bibitem{oracle-rpo} Forbes / Motley Fool analyses (Jan.--Jul.\ 2026): Oracle RPO \$455B$\to$\$523B$\to$\$638B; OpenAI contract $>$\$300B over five years from 2027, est.\ 40--60\% of backlog; OpenAI reassessment right after $\sim$5 yrs vs 15--19-yr Oracle leases; $\sim$\$380B debt-like obligations; negative FCF; CDS $>$130\,bps; both agencies negative outlook. \url{https://www.forbes.com/sites/roomykhan/2026/01/15/oracles-ai-bet-equity-wins-and-credit-waits/}

% ---- author position / institutional record ----
\bibitem{genesis-japan} U.S.\ Department of Energy, MEXT and METI, Statement of Intent establishing Japan as the first international partner under the Genesis Mission (4 Jun.\ 2026): \$1B over five years, \$500M from each nation, binding twelve DOE national laboratories to twelve Japanese research institutions---RIKEN, the University of Tokyo, NIMS, KEK and J-PARC among them---across quantum information science, fusion, biotechnology, advanced materials, particle physics, and AI-and-robotics autonomous laboratories, with access to DOE systems and Fugaku. \url{https://www.energy.gov/articles/united-states-and-japan-announce-historic-1-billion-partnership-under-president-trumps}
\bibitem{riken-genesis} RIKEN, ``Message from President Gonokami regarding the new Japan--US strategic partnership on `AI for Science''' (5 Jun.\ 2026): RIKEN as a central implementing institution, with Fugaku, SPring-8 and the BioResource Center named, and the TRIP (2022) and TRIP-AGIS (2023) platforms as the standing AI-for-Science vehicles; collaborating US laboratories named as Argonne, Oak Ridge and Brookhaven. \url{https://www.riken.jp/en/news_pubs/news/2026/20260605_1/index.html}
\bibitem{anl-riken} Argonne National Laboratory with RIKEN, Fujitsu and NVIDIA, partnership on AI for science and next-generation computing (28 Jan.\ 2026). \url{https://www.anl.gov/article/argonne-partners-with-riken-fujitsu-and-nvidia-to-advance-ai-for-science-and-next-generation-computing}

% ---- primary sources on the Qwen3.8-Max release ----
\bibitem{qwen-announcement} Qwen (@Alibaba\_Qwen), official release announcement, 3 Aug.\ 2026 (pinned): ``Next week, the open weights of Qwen3.8-Max will be released, and Qwen3.8-27B is also going open-weights to meet you all''; 2.4T parameters; autonomous coding described as ``10+ days of self-evolving development, from empty folder to production without hand-holding, complete project trace in the GitHub''; ``system-level autonomous planning with closed-loop adaptive learning, driving 500+ turns of chip design optimization and 365 days of e-commerce strategy''; native multimodal intelligence as ``a continuous feedback loop for planning, execution, and self-correction''; pricing \$2.0/\$6.0 per Mtok with \$0.25 implicit caching. \url{https://x.com/Alibaba_Qwen}
\bibitem{qwen38-blog} Qwen Team, Qwen3.8 release blog with the full cross-model benchmark table, Aug.\ 2026. \url{https://qwen.ai/blog?id=qwen3.8}
\bibitem{berman2026} M.~Berman, ``Open-source is WINNING,'' Matthew Berman (YouTube), Aug.\ 2026. Review of the Qwen3.8-Max release: identifies research reproduction (PaperBench 93.0) as ``the required skill for automated AI research,'' noting the model ``started from nothing but the paper and a set of GPUs---no starter code, no ready-made pipeline'' and ``invented and tested 18 improvement ideas of its own across four rounds''; reports Alibaba's claim that the model ``has independently achieved the autonomous execution of the entire silicon design flow,'' with the observation that ``one of the things that China is behind on is chip design and chip manufacturing---but if they had a model that could help them design these chips, that would no longer be their weakness''; per-task cost data (Qwen3.7-Max \$1.28 vs GPT-5.6 Sol \$1.23, with Anthropic models most expensive); the model-size gap (US closed frontier rumoured $\sim$7T+ against 2.4--2.8T Chinese); the enterprise-adoption argument (``95\% of the capability at a fraction of the price''); the prediction that ``the model will be co-designed with the chip\ldots{} with extreme co-design between model and hardware it might actually put us, as we're becoming dependent on Chinese open source, highly dependent on Chinese chips as well''; and the recursive-self-improvement counter-thesis, including OpenAI's attribution of its 80\% price cut to efficiency gains produced by its own frontier model. \url{https://youtu.be/CVlKp9Ld-Zg}

% ---- the eighteen-day barrage ----
\bibitem{qwen38-max} MarkTechPost / SiliconANGLE / Forbes (3 Aug.\ 2026): Alibaba releases Qwen3.8-Max---2.4T-parameter sparse MoE, 1M context (991K in / 131K out / 262K reasoning budget), multimodal input, \$2/\$6 per Mtok flat across the full window; \emph{benchmark table published by the vendor; licence terms and activated-parameter count not stated}; a $\sim$95B active figure is press consensus of unclear provenance. \url{https://www.marktechpost.com/2026/08/03/alibaba-qwen-releases-qwen3-8-max/}
\bibitem{qwen38-license} South China Morning Post (3 Aug.\ 2026), ``Alibaba's Qwen3.8-Max made widely accessible ahead of open-weights release'': corroborates the open-weight commitment and the Qwen3.8-27B dense companion, with the repository pending at publication. \url{https://www.scmp.com/tech/article/3362738/alibabas-ai-model-qwen38-max-made-widely-accessible-ahead-open-weights-release}
\bibitem{qwen38-bench} Vendor-reported benchmark table for Qwen3.8-Max as compiled by apidog and OfficeChai (3 Aug.\ 2026): Terminal-Bench 2.1 86.6; SWE-bench Pro 67.7; PaperBench 93.0; GPQA Diamond 92.6; IFBench 82.8; HLE 43.6; CoWorkBench 74.8; OSWorld-Verified 86.1. No SWE-bench Verified, no AIME, no Kimi K3 or DeepSeek V4 comparators, and---unlike Qwen3.7-Max---no hallucination-rate disclosure. All figures are Alibaba's own. \url{https://apidog.com/blog/qwen-3-8-benchmarks/}
\bibitem{arena-qwen38} Arena.AI / LMArena standings reported by CGTN and SiliconANGLE (Aug.\ 2026): Qwen3.8-Max \#5 overall in text (highest-ranked Chinese model), \#2 in vision, \#4 on Frontend Code Arena at 1,668 Elo---8 points behind Kimi K3, $\sim$37 behind the leading Claude configuration. Artificial Analysis had not scored the model. \url{https://news.cgtn.com/news/2026-08-03/Alibaba-unveils-Qwen3-8-Max-its-most-capable-AI-model-to-date-1Pj3AAwmjPa/p.html}
\bibitem{infoworld-qwen} InfoWorld (3 Aug.\ 2026), ``Alibaba takes aim at OpenAI and Anthropic with Qwen3.8-Max launch'': Forrester and Gartner assessments; the formulation that ``publishing weights is a separate act from opening an API endpoint---until there is a repository, a licence and a model card, open-weight describes an intention''; and the unaudited-autonomy critique (``sixteen days of what? how many times did a human step in?''). \url{https://www.infoworld.com/article/4204415/alibaba-takes-aim-at-openai-and-anthropic-with-qwen3-8-max-launch.html}
\bibitem{v4flash} Caixin Global (1 Aug.\ 2026) and technical write-ups of DeepSeek V4-Flash-0731 (official release 31 Jul.\ 2026; preview from 24 Apr.): $\sim$280B total / 13B active MoE, MIT licence with weights published, 1M context, 384K max output, \$0.14/\$0.28 per Mtok; generation gains attributable to post-training alone (Terminal-Bench 2.1 61.8$\to$82.7 on unchanged architecture); Artificial Analysis Intelligence Index 50; announced 2$\times$ peak-hour surge pricing across seven daily hours; the preview was the most-used model on OpenRouter for seven consecutive weeks. \url{https://www.caixinglobal.com/2026-08-01/deepseek-releases-official-v4-flash-model-as-chinas-ai-race-intensifies-102470292.html}
\bibitem{pareto-frontier} ZeroHedge (Aug.\ 2026), ``China's AI knife fight'': the eighteen-day chronology (16 Jul.\ Kimi K3 $\to$ 27 Jul.\ K3 weights $\to$ 31 Jul.\ V4-Flash $\to$ 3 Aug.\ Qwen3.8-Max); four of the five models on the cost-performance Pareto frontier Chinese (Kimi K3, Qwen3.8-Max, GLM-5.2, DeepSeek V4-Flash) with Claude Opus 5 the lone American entry defended by $\sim$37 Elo; benchmark-suite run costs of $\sim$\$0.03 (V4-Flash) against \$3.15 (Fable 5) and \$1.86 (GPT-5.6 Sol). \url{https://www.zerohedge.com/ai/chinas-ai-knife-fight-deepseeks-new-model-runs-100x-cheaper-anthropics-flagship}
\bibitem{openai-cut} Forbes / G.~Wang (31 Jul.\ 2026), ``Why OpenAI's 80\% price cut could trigger a race to the bottom in AI'': GPT-5.6 Luna input \$1.00$\to$\$0.20 and output \$6.00$\to$\$1.20 (80\% cut), Terra down 20\%, Sol unchanged; OpenAI attributed the cut to system efficiency while the analysis attributes competitive pressure to Kimi K3 and DeepSeek V4. Note the sequencing: the cut preceded Qwen3.8-Max's general availability by three days. \url{https://www.forbes.com/sites/geruiwang/2026/07/31/why-openais-80-price-cut-could-trigger-a-race-to-the-bottom-in-ai/}
\bibitem{qwen-offshore} Financial Times reporting relayed by Tom's Hardware (2026): leading Chinese AI firms, Alibaba among them, train models on NVIDIA GPUs leased from operators in Singapore and Malaysia---legal under current export rules because no controlled hardware enters China---with the resulting weights served domestically. No public source connects Qwen3.8-Max training to domestic accelerators. \url{https://www.tomshardware.com/tech-industry/semiconductors/chinas-top-ai-firms-shift-model-training-overseas-to-access-nvidia-gpus}
\bibitem{alibaba-moonshot} TechTimes (2 Aug.\ 2026): Alibaba, Moonshot's largest investor, supplied roughly 20,000 NVIDIA Hopper-class accelerators under a compute agreement (Bloomberg and Reuters corroborated; Alibaba disputes the H200 characterization but not the arrangement); Kimi K3 subsequently outscored Qwen models on several evaluations run on that compute. \url{https://www.techtimes.com/articles/322699/20260802/alibaba-bankrolled-kimi-k3-20000-nvidia-chips-now-qwen-losing-its-own-compute.htm}
\bibitem{qwen-downloads} South China Morning Post, reporting an Interconnects AI analysis of Hugging Face data (10 Apr.\ 2026): Qwen at $\sim$1 billion cumulative downloads by March 2026, 153.6M in February alone---more than double the combined total of the next eight families---and over 50\% of all global open-model downloads, with $>$200,000 derivative variants. \url{https://www.scmp.com/tech/big-tech/article/3349552/alibabas-qwen-family-captures-over-50-global-open-source-downloads-report-finds}
\bibitem{alibaba-capex} BNN Bloomberg, Alibaba FY results (May 2026): three-year AI capex plan of RMB\,380B ($\sim$\$56B) to be exceeded without a new stated target; Cloud Intelligence revenue $+$38\% y/y with AI products $\sim$30\% of external cloud revenue; ex-items net income down 99.7\% and adjusted EBITDA down 84\%; CEO statement that ``margin is still secondary''; target of $>$\$100B combined external AI and cloud revenue over five years. \url{https://www.bnnbloomberg.ca/business/2026/05/13/alibaba-reports-continued-profit-squeeze-from-spending-on-ai-and-instant-retail/}

% ---- countries and China self-critique ----
\bibitem{japan-ai-act} White \& Case / FPF: Japan's AI Promotion Act (passed 28 May 2025)---first AI-dedicated law; innovation-first, no penalties, Cabinet-level AI Strategy Headquarters and national AI Basic Plan. \url{https://www.whitecase.com/insight-alert/japans-first-ai-legislation-becomes-law-focus-promoting-research-and-development-no}
\bibitem{japan-budget} Japan Times / Bloomberg (Dec.\ 2025): government moves to roughly quadruple annual chips/AI budget support, deepening state backing for Rapidus; follows the ¥10T-through-2030 package. \url{https://www.japantimes.co.jp/business/2025/12/26/economy/ai-budget-support/}
\bibitem{eu-omnibus} Gibson Dunn / European Parliament: Digital Omnibus provisional agreement (May--Jun.\ 2026)---AI Act Annex III high-risk obligations postponed to 2 Dec.\ 2027, embedded high-risk to Aug.\ 2028; AI Office GPAI powers strengthened. \url{https://www.gibsondunn.com/eu-ai-act-omnibus-agreement-postponed-high-risk-deadlines-and-other-key-changes/}
\bibitem{eu-giga-tender} EuroHPC JU (30 Jul.\ 2026): formal call for up to seven AI gigafactories (each 100k+ advanced chips), expected to trigger $>$€30B; 76 expressions of interest across 16 member states. \url{https://www.eurohpc-ju.europa.eu/eurohpc-joint-undertaking-launches-ai-gigafactories-call-2026-07-30_en}
\bibitem{mistral-asml} Mistral AI / CNBC (9 Sept.\ 2025): €1.7B Series C at €11.7B ($\sim$\$14B) post-money, led by ASML (€1.3B, $\sim$11\%). \url{https://www.cnbc.com/2025/09/09/ai-firm-mistral-valued-at-14-billion-as-asml-takes-major-stake.html}
\bibitem{indiaai} IndiaAI Mission GPU tenders: $\sim$38,000 GPUs empanelled across three rounds at subsidized rates reportedly under \$1/GPU-hr; reported utilization shortfall (``can't find enough users''). \url{https://www.communicationstoday.co.in/38000-gpus-one-big-problem-indias-ai-mission-cant-find-enough-users/}
\bibitem{korea-hbm} Counterpoint / TrendForce: SK hynix $\sim$62\% HBM share; reportedly $\sim$two-thirds of leading HBM4 volume; Samsung/SK hynix scaling into the 2026 memory cycle. \url{https://www.astutegroup.com/news/general/sk-hynix-holds-62-of-hbm-micron-overtakes-samsung-2026-battle-pivots-to-hbm4/}
\bibitem{korea-sovereign} Korea.net / Korea Herald / NVIDIA: national AI computing center; $>$260,000-GPU APEC deal (Oct.\ 2025) across government, Samsung, SK, Hyundai, Naver; sovereign foundation-model tournament---five teams selected Aug.\ 2025, public elimination rounds 2026, 700B-class open models released. \url{https://nvidianews.nvidia.com/news/south-korea-ai-infrastructure}
\bibitem{sea-lion} AI Singapore, SEA-LION v4 (Aug.\ 2025): open multimodal 27B model for 11+ Southeast Asian languages. \url{https://sea-lion.ai/blog/sea-lion-v4-multimodal/}
\bibitem{gulf-ai} NVIDIA/HUMAIN partnership (500\,MW-class ``AI factories,'' several hundred thousand GPUs); AMD--HUMAIN up to \$10B/6\,GW; Stargate UAE 1\,GW within a 5\,GW campus; US Commerce authorization of up to 70,000 advanced chips to G42/HUMAIN (Nov.\ 2025). \url{https://www.commerce.gov/news/press-releases/2025/11/statement-uae-and-saudi-chip-exports}
\bibitem{taiwan-energy} Tom's Hardware / CAPRI: Taiwan's $\sim$NT\$100B ``ten major AI infrastructure'' plan; post-nuclear grid strain as the binding constraint on AI and foundry expansion. \url{https://www.tomshardware.com/tech-industry/semiconductors/taiwan-to-spend-billions-on-ai-island-push-despite-grid-risks}
\bibitem{cassava} Cassava Technologies (Mar.\ 2026): Africa's first NVIDIA-powered AI factory (South Africa); Kenya, Nigeria, Egypt, Morocco expansions planned. \url{https://www.cassavatechnologies.com/cassava-scales-african-ai-infrastructure-with-nvidia-powered-ai-factories-to-accelerate-sovereign-data-capabilities/}
\bibitem{china-involution} Caixin / Fortune / State Council reporting: anti-``involution'' campaign; pricing-law amendments (Aug.\ 2025); State Council warning against ``disorderly competition'' in AI; Xi Jinping publicly questioning whether every province needs AI/compute/EV projects. \url{https://fortune.com/asia/2025/08/29/china-warns-against-disorderly-competition-ai-race}
\bibitem{china-monetization} Recode China AI / Asia Tech Review / Sinocities: Zhipu and MiniMax post-IPO disclosures show thin revenue and heavy losses; documented ``low-trust open-source paradox'' limiting domestic enterprise adoption of domestic open models. \url{https://www.recodechinaai.com/p/zhipu-ai-and-minimax-just-went-public}
\bibitem{china-idle-dc} MIT Technology Review (Mar.\ 2025) / Light Reading / ASPI: hundreds of state-subsidized Chinese AI data centers with reported utilization far below plan (up to $\sim$80\% of some new capacity unused); rental-price collapse ending the buildout binge. \url{https://www.technologyreview.com/2025/03/26/1113802/china-ai-data-centers-unused/}
\bibitem{china-hbm-gap} ChinaTalk / CSIS / trade reporting: HBM and advanced packaging as binding constraints on the Ascend fleet; CXMT assessed roughly three years behind the Korean leaders. \url{https://www.chinatalk.media/p/mapping-chinas-hbm-advancement}

% ---- embodied AI ----
\bibitem{cosmos} NVIDIA, Cosmos world-model platform; Cosmos 3 omnimodal world models for physical AI (Jun.\ 2026). \url{https://developer.nvidia.com/blog/develop-physical-ai-reasoning-world-and-action-models-with-nvidia-cosmos-3/}
\bibitem{genie3} Google DeepMind, ``Genie 3: a new frontier for world models'' (Aug.\ 2025). \url{https://deepmind.google/blog/genie-3-a-new-frontier-for-world-models/}
\bibitem{helix} Figure AI Helix VLA (S2 deliberative 7--9\,Hz + S1 visuomotor 200\,Hz; trained on $\sim$500\,h teleoperation); Helix-02 completed a livestreamed 8-hour autonomous package-sorting shift at human-comparable rates (13 May 2026); 350+ Figure 03s delivered. \url{https://www.techtimes.com/articles/316632/20260514/figure-ais-helix-02-robots-complete-full-8-hour-autonomous-shifts-humanoid-race-intensifies.htm}
\bibitem{pi-models} Physical Intelligence $\pi$-series VLA models; funding \$400M@\$2.4B (2024) $\to$ $\sim$\$5.6B, reported talks at $>$\$11B (Mar.\ 2026); Skild AI reported $\sim$\$14B. \url{https://mlq.ai/news/physical-intelligence-in-talks-to-raise-1-billion-at-11-billion-valuation/}
\bibitem{data-farms} Rest of World, ``Inside China's robot training centers'' (2026): 40+ state-backed data-collection centers announced ($\sim$24 operational); teleoperators with VR/exoskeletons; UBTech RMB\,566M of sales to three provincial centers; China Mobile RMB\,124M orders to Unitree/AgiBot. \url{https://restofworld.org/2026/china-robots-training-centers-workers/}
\bibitem{kylin} China Daily / Global Times: Shanghai heterogeneous humanoid training facility (Jan.\ 2025), 100+ robots simultaneously, target 1,000 by 2027 collecting 10M+ entries. \url{https://www.chinadaily.com.cn/a/202501/23/WS67919d43a310a2ab06ea8c2e.html}
\bibitem{agibot-world} AgiBot World Colosseo open embodied dataset: 1M+ real-robot trajectories, 100 robots, 100+ scenarios; Genie Sim 3.0 open-source simulation (CES 2026). \url{https://opendrivelab.com/AgiBot-World/}
\bibitem{open-x} Open X-Embodiment collaboration (Google DeepMind + 33 institutions): pooled cross-embodiment dataset and RT-X models. \url{https://robotics-transformer-x.github.io/}
\bibitem{agility-spac} GeekWire, Agility Robotics SPAC filings (2026): \$2.5B valuation; $>$\$300M committed multi-year Digit orders incl.\ a 1,000-robot contract; 65,000+ hours of real-world operation; 2025 opex \$111M; GXO multi-year RaaS deployment. \url{https://www.geekwire.com/2026/digit-maker-agility-robotics-to-go-public-in-2-5b-deal-heres-what-the-filings-say-about-its-finances/}
\bibitem{optimus-zero} TechTimes/Electrek (Jul.\ 2026): Tesla Optimus at zero verified production units; Fremont line start slated late Jul./Aug.\ 2026; Gen-3 reveal deferred twice; earnings-call statement that existing units are ``primarily for learning, not productive tasks.'' \url{https://www.techtimes.com/articles/321012/20260720/tesla-optimus-production-count-remains-zero-q2-earnings-call-looms-wednesday.htm}
\bibitem{reducers} Humanoid.guide / Pandaily / Forbes: Leaderdrive as China's largest harmonic-reducer maker lifted by humanoid demand; Zhongda Leader founders became billionaires on joint demand (Apr.\ 2026); Unitree price war squeezing the reducer chain. \url{https://www.forbes.com/sites/zinnialee/2026/04/23/humanoid-mania-turns-chinese-brothers-behind-robotic-joint-maker-into-billionaires/}
\bibitem{roller-screws} Fast Company, ``This tiny screw is powering the humanoid robot revolution'': planetary roller screws $\sim$33\% of humanoid BoM (J.P. Morgan est.); \$1,350--2,700/unit; China-dominant production; Shanghai Beite's \$260M plant. \url{https://www.fastcompany.com/91314612/this-tiny-screw-is-powering-the-humanoid-robot-revolution}
\bibitem{tesla-suppliers} 36Kr / Futu: reported $\sim$\$685M Optimus actuator-assembly order to Sanhua; Ningbo Tuopu building robot drive-system capacity. \url{https://eu.36kr.com/en/p/3510288514980998}
\bibitem{lidar-robotics} Yicai Global (Q2 2025): Hesai robotics lidar shipments 49,000 units ($+$744\% y/y); RoboSense 34,400 ($+$632\%); both pivoting from automotive to robotics. \url{https://www.yicaiglobal.com/news/chinas-robosense-hesai-see-robot-lidar-sales-surge-in-second-quarter}
\bibitem{catl-galbot} CnEVPost / SCMP (24 Jun.\ 2026): CATL--Galbot agreement; Galbot S1 deployed at scale on CATL battery module/pack lines; Galbot cumulative funding $\sim$\$800M at $\sim$\$3B. \url{https://cnevpost.com/2026/06/24/catl-galbot-scale-humanoid-robots-battery-factories/}
\bibitem{jetson-thor} NVIDIA, Jetson Thor launch (25 Aug.\ 2025): Blackwell-based, 2,070 FP4 TFLOPS, 128\,GB, dev kit \$3,499; Isaac GR00T open reference humanoid for academic research (Jun.\ 2026). \url{https://nvidianews.nvidia.com/news/nvidia-blackwell-powered-jetson-thor-now-available-accelerating-the-age-of-general-robotics}
\bibitem{iso-humanoid} ISO 25785-1 (industrial mobile robots with actively controlled stability) drafting toward $\sim$2027 publication; ISO 10218:2025 integrating ISO/TS 15066 force limits; UL 3300 for consumer robots. \url{https://humanoidliability.com/resources/iso-25785-humanoid-safety-standards/}
\bibitem{goldman-humanoid} Goldman Sachs Research: humanoid TAM \$38B by 2035; unit manufacturing cost fell \$50--250k $\to$ \$30--150k in one year ($\sim$40\% vs 15--20\% expected). \url{https://www.goldmansachs.com/insights/articles/the-global-market-for-robots-could-reach-38-billion-by-2035}
\bibitem{ms-humanoid} Morgan Stanley (Jun.\ 2026): China humanoid shipments 12,000 (2025) $\to$ forecast 50,000 (2026) $\to$ 446,000/yr by 2030; localized Chinese chains $\sim$20\% cheaper; average price $\sim$\$46k trending to $\sim$\$21k; ``Humanoid 100'' value-chain map. \url{https://www.cnbc.com/2026/06/24/morgan-stanley-china-humanoid-robot-market-forecast.html}
\bibitem{unitree-ipo} Gasgoo / Fortune (Jun.\ 2026): Unitree STAR-market IPO approved; 5,500 humanoids sold 2025; humanoid ASP RMB\,593k (2023) $\to$ RMB\,166k (2025); gross margin 87.7\% $\to$ 63.2\%; $\sim$\$250M revenue with $\sim$\$41M profit---among the only profitable humanoid makers. \url{https://fortune.com/2026/06/06/chinese-humanoid-robots-global-market-sales-performative-functional/}
\bibitem{optimus-china} BigGo / Tech Digest / 36Kr supply-chain reporting [A]: $\sim$70\% of Optimus supply chain China-based; excluding Chinese parts estimated to triple cost; 2025 rare-earth licensing touched actuator magnets. \url{https://www.techdigest.tv/2026/02/excluding-chinese-parts-could-triple-tesla-optimus-costs.html}
\bibitem{fortune-humanoid} Fortune (6 Jun.\ 2026): 2025 global humanoid shipments $>$13,000, $\sim$85\% Chinese; US firms ``a few hundred or less''; most units performative rather than functional---demos, exhibitions, education. \url{https://fortune.com/2026/06/06/chinese-humanoid-robots-global-market-sales-performative-functional/}
\bibitem{agibot-10k} The Robot Report (31 Mar.\ 2026): AgiBot ships 10,000th humanoid; second five thousand in $\sim$3 months; Omdia/IDC \#1 by 2025 shipments. \url{https://www.therobotreport.com/agibot-rolls-out-10000th-humanoid-robot/}
\bibitem{ubtech} PRNewswire / Interesting Engineering: UBTech Walker S2 mass production from Nov.\ 2025, 1,000th unit built; cumulative Walker orders $>$RMB\,800M; multi-robot deployments at Zeekr, BYD, Foxconn, Audi-FAW. \url{https://www.prnewswire.com/news-releases/ubtech-humanoid-robot-walker-s2-begins-mass-production-and-delivery-with-orders-exceeding-800-million-yuan-302616924.html}
\bibitem{byd-robots} TechTimes / ChinaBizInsider (Jun.\ 2026): BYD confirms four-year internal humanoid program; reported target $\sim$20,000 robots in its own factories in 2026. \url{https://www.techtimes.com/articles/317949/20260607/byd-confirms-four-year-secret-humanoid-robot-program-could-sell-through-car-dealerships.htm}
\bibitem{xpeng-iron} Gasgoo / Humanoids Daily: XPeng IRON mass production targeted end-2026; stated goal 1M units by 2030. \url{https://www.humanoidsdaily.com/news/xpeng-sets-2026-target-for-mass-produced-iron-robot-eyes-1-million-units-by-2030}
\bibitem{miit-humanoid} MIIT, ``Guiding Opinions on Innovation and Development of Humanoid Robots'' (Nov.\ 2023): initial mass production by 2025; humanoids as ``an important new engine of economic growth'' by 2027; embodied AI in the 2025 Government Work Report and 15th Five-Year Plan. \url{https://www.therobotreport.com/china-plans-to-mass-produce-humanoids-by-2025/}
\bibitem{embodied-policy} Carnegie Endowment / Jamestown Foundation (2025--26): \$8.2B national AI-fund allocation; Beijing RMB\,100B robotics fund; Shenzhen/Guangdong/Beijing action plans with enterprise and revenue targets; 150+ humanoid companies; H1 2025 humanoid financing $>$RMB\,10B. \url{https://carnegieendowment.org/research/2025/11/embodied-ai-china-smart-robots}
\bibitem{figure-c} Figure AI Series C (Sept.\ 2025): $\sim$\$1B at \$39B post-money. \url{https://sacra.com/c/figure-ai/}
\bibitem{atlas} Forbes / Hyundai (CES 2026): electric Atlas 2026 production run fully committed; Hyundai plant designed for 30,000/yr; deployment of a committed 25,000 units gated by union negotiations. \url{https://www.forbes.com/sites/johnkoetsier/2026/01/06/atlas-humanoid-robots-production-fully-committed-for-2026-factory-will-build-30000-per-year/}
\bibitem{teleop} Futurism / Electrek / Humanoids Daily: Tesla ``We, Robot'' bartenders remotely operated (analyst-confirmed); Dec.\ 2025 Optimus demo teleoperation controversy; 1X NEO ships with explicit human-teleoperation fallback. \url{https://www.humanoidsdaily.com/news/1x-neo-launch-sparks-debate-on-autonomy-and-teleoperation}
\bibitem{ifr} International Federation of Robotics, \emph{World Robotics 2025}: 542,000 industrial robots installed 2024; China 295,000 (54\%); operational stock 4.66M (China $>$2M); Chinese makers 57\% domestic share vs $\sim$28\% a decade earlier. \url{https://ifr.org/downloads/press_docs/2025-09-25-IFR_press_release_China_in_English.pdf}

% ============ ADDITIONS ============
\bibitem{agrawal2026} A.~K.~Agrawal, J.~McHale and A.~Oettl, ``AI in Science,'' NBER Working Paper 34953, 2026. Models AI in science as augmentation through search, with large variation across workflow stages and fields. \url{https://doi.org/10.3386/w34953}
\bibitem{fort2026} T.~C.~Fort, N.~Goldschlag, J.~Liang, P.~K.~Schott and N.~Zolas, ``Growth is Getting Harder to Find, Not Ideas,'' NBER Working Paper 35182, 2026. \url{https://doi.org/10.3386/w35182}
\bibitem{hyperion-roi} Hyperion Research Holdings, LLC, \emph{Analysis of the Ripple Effects of the Next-Generation Computing Platform}, supplementary final survey report, revised ed., commissioned by RIKEN, 5 Mar.\ 2025 (\#HR4.0523.02.05.2025). Fugaku COVID-19 projects estimated at \$50--75B returned to the Japanese economy; successor-machine lifetime return scenarios \$250B/\$400B+/\$600B+; 65.7\% of surveyed Fugaku researchers reporting realized results on SDG-class projects; innovation-class mapping of 650+ worldwide HPC projects; international five-year AI investment estimates.
\bibitem{hill2026} R.~R.~Hill and C.~Stein, ``How Artificial Intelligence Shapes Science: Evidence from AlphaFold,'' NBER Working Paper 35143, 2026. \url{https://doi.org/10.3386/w35143}
\bibitem{liu2026novelty} M.~Liu and C.~Zhai, ``An Axiomatic Benchmark for Evaluation of Scientific Novelty Metrics,'' arXiv:2604.15145, 2026. \url{https://arxiv.org/abs/2604.15145}
\bibitem{xiong2026} L.~Xiong et al., ``AutoResearchBench: Benchmarking AI Agents on Complex Scientific Literature Discovery,'' arXiv:2604.25256, 2026: frontier agents below 10\% on demanding literature-discovery tasks. \url{https://arxiv.org/abs/2604.25256}
\bibitem{kimiK3-report} Kimi Team et al., ``Kimi K3: Open Frontier Intelligence,'' arXiv:2607.24653, 2026: 2.8T total / 104B activated parameters, 1M-token context, native vision; the report itself acknowledges trailing the strongest proprietary systems overall. \url{https://arxiv.org/abs/2607.24653}
\bibitem{kimiLicense} Moonshot AI, ``Kimi K3 License'' (custom commercial license: use, modification and redistribution permitted; separate agreement required for qualifying model-as-a-service businesses above a revenue threshold; attribution conditions on very large deployments). \url{https://huggingface.co/moonshotai/Kimi-K3/blob/main/LICENSE}
\bibitem{qwen36-repo} Qwen Team, Qwen3.6 official repository: open releases are the 27B dense and 35B-A3B variants, Apache 2.0. \url{https://github.com/QwenLM/Qwen3.6}
\bibitem{reutersQwen38} Reuters, ``Alibaba unveils its largest AI model yet; DeepSeek's latest model is ultra-low cost'' (3 Aug.\ 2026). \emph{Superseded} by the primary-source entries above: the report's pairing of the two releases obscured that DeepSeek V4-Flash shipped officially on 31 July, and its $\sim$95B activated-parameter figure for Qwen3.8-Max is press consensus rather than a vendor disclosure. \url{https://www.reuters.com/business/retail-consumer/alibaba-unveils-its-most-capable-ai-model-date-not-far-behind-moonshots-size-2026-08-03/}
\bibitem{microsoft2026} Microsoft, FY26 Q3 earnings release (29 Apr.\ 2026): AI annual revenue run rate above \$37B, $+$123\% y/y. \url{https://www.microsoft.com/en-us/investor/earnings/fy-2026-q3/press-release-webcast}
\bibitem{nvidia2026} NVIDIA, Q1 FY2027 financial results (20 May 2026): revenue \$81.6B, data center \$75.2B, non-GAAP gross margin $\sim$75\%. \url{https://investor.nvidia.com/news/press-release-details/2026/NVIDIA-Announces-Financial-Results-for-First-Quarter-Fiscal-2027/default.aspx}
\bibitem{imePlan} RISC-V International, IME task-group ratification plan (target Aug.\ 2026). \url{https://riscv.atlassian.net/wiki/spaces/IMEX/pages/598867969/IME+Ratification+Plan}
\bibitem{amePlan} RISC-V International, AME task-group ratification plan (board-ratification target Apr.\ 2027; simulators, ABI, GCC/LLVM support and basic kernels as staged milestones). \url{https://riscv.atlassian.net/wiki/spaces/AMEX/pages/55083420/Ratification+Plan}
\bibitem{skhynixSocamm} SK hynix, ``SK hynix Begins Mass Production of 192GB SOCAMM2'' (20 Apr.\ 2026). \url{https://news.skhynix.com/mass-production-socamm2-192gb/}
\bibitem{nvidiaVera} NVIDIA, Vera data-center CPU: up to 1.5\,TB LPDDR5X capacity and $\sim$1.2\,TB/s memory bandwidth per socket. \url{https://www.nvidia.com/en-us/data-center/vera-cpu/}
\bibitem{reutersCXMT} Reuters, ``CXMT plans second chip plant in Beijing and is in talks on its funding, sources say'' (3 Aug.\ 2026): three fabs of $\sim$100k wafers/month each, potential expansion beyond 600k wpm; CXMT remains far smaller than the three global leaders. \url{https://www.reuters.com/world/asia-pacific/cxmt-plans-second-chip-plant-beijing-is-talks-its-funding-sources-say-2026-08-03/}

% ============ AUGUST 2026 SECOND-WEEK CYCLE ============
\bibitem{muse-glimmer} Meta AI Research, ``Introducing Muse Glimmer: An Open Agentic Model That Runs on Your Device'' (10 Aug.\ 2026): 30B dense multimodal, Apache 2.0, weights on Hugging Face/Ollama; distilled from Muse Spark. \url{https://research.meta.ai/blog/introducing-muse-glimmer-open-agentic-model}
\bibitem{zuck-letter} M.~Zuckerberg, ``The Future Is For Everyone'' (10 Aug.\ 2026): 6,500-word letter recommitting Meta to open weights; companion X announcement of Muse Glimmer weights and a forthcoming open-weight version of Muse Spark 1.2; contemporaneous coverage in WSJ/NYT/Bloomberg. \url{https://www.meta.com/thefutureisforeveryone/}
\bibitem{muse-aa} Artificial Analysis, ``Muse Glimmer'' (Aug.\ 2026): independent evaluation---Intelligence Index 35 (Gemma 4 31B: 30; Qwen3.6 27B Reasoning: 38; Kimi K2.5: 36); measured hallucination rate 82\% vs 49\% for Qwen3.6 27B; ``Meta's most permissive license yet.'' \url{https://artificialanalysis.ai/articles/muse-glimmer}
\bibitem{qwen-slip} Digital Applied, ``Qwen3.8 open weights: checklist before download'' (7 Aug.\ 2026) and byteiota (9 Aug.\ 2026): the 3 Aug.\ ``next week'' commitment lapsed without a repository; window re-framed to the week of 10 Aug.; fourth consecutive Max-cycle slippage documented. \url{https://www.digitalapplied.com/blog/qwen3-8-open-weights-checklist-before-download}
\bibitem{minimax-geofence} TechTimes (4 Aug.\ 2026): MiniMax H3 tops the open-weight video arena; licence excludes local deployment in the US, EU, UK, and South Korea. \url{https://www.techtimes.com/articles/322904/20260804/minimax-h3-open-weights-exclude-us-eu-uk-korea-local-deployment.htm}
\bibitem{qianwen-platform} Caixin Global (11 Aug.\ 2026): Alibaba opens the Qianwen platform to outside developers for custom AI agents across phones, PCs, and smart glasses; launch partners SF Express, Ziroom, Midea. \url{https://www.caixinglobal.com/2026-08-11/alibaba-opens-qwen-platform-to-outside-developers-in-push-to-build-ai-agent-ecosystem-102473024.html}
\bibitem{mofcom-retaliation} CNBC (6 Aug.\ 2026): Beijing's broadest countermeasures since the October truce---case-by-case export review on drones and components, six US entities sanctioned, national-security probe into imported printers; triggered by FCC Covered List and UFLPA additions in late July. \url{https://www.cnbc.com/2026/08/06/us-china-trump-xi-summit-trade-technology-restrictions.html}
\bibitem{wh-framework} Fortune (4 Aug.\ 2026) / Axios (3 Aug.\ 2026): White House finalizes a voluntary frontier-model evaluation framework in a closed-door session with OpenAI, Anthropic, Microsoft, Meta, NVIDIA and others; document not publicly released. \url{https://fortune.com/2026/08/04/baffling-white-house-wont-publicly-release-ai-model-evaluation-framework-it-reviewed-today-with-openai-anthropic-microsoft-and-others/}
\bibitem{astra2026} OpenAI, ``Ten advances in mathematics'' (Aug.\ 2026): the Astra model's solutions to ten long-open problems in mathematics and quantum complexity, proofs published; $\sim$\$2,000 of inference compute reported. \url{https://openai.com/index/ten-advances-in-mathematics/}
\bibitem{nist-genesis} NIST (4 Aug.\ 2026): NIST joins the Genesis Mission with two AI Economic Security Centers---AI agents for manufacturing (tenfold drone-production scaling goal in two years) and for critical-infrastructure cyber-defense. \url{https://www.nist.gov/news-events/news/2026/08/nist-joins-national-genesis-mission-accelerate-ai-innovation}
\bibitem{coreweave-q2} Benzinga (8 Aug.\ 2026): CoreWeave ahead of Q2---bear market off the July low, 15\% short interest, net loss widened to \$740M, debt above \$25B against \$30--35B capex guidance, Jefferies and Raymond James downgrades; revenue nonetheless doubled y/y to \$2.07B on a \$99B backlog. \url{https://www.benzinga.com/markets/equities/26/08/61062163/coreweave-stock-stuck-in-a-bear-market-as-short-interest-rises-ahead-of-earnings}


% ============ MID-AUGUST 2026: HARNESSES, WEIGHTS, ROADMAP ============
\bibitem{qwen38-weights} Alibaba / Qwen, model repositories \texttt{Qwen/Qwen3.8-2.4T-A95B} (2.4T total / 95B active; 512 experts, 11 activated; 262,144 native context; text-only; 4.89\,TB in 213 BF16 shards) and \texttt{Qwen/Qwen3.8-27B} (Apache 2.0, vision-language). The flagship carries the custom \texttt{qwen3.8-max} licence: separate commercial agreement required where model-as-a-service or AI-work-assistant revenue exceeds US\$50M in any consecutive 12 months; attribution required above 100M MAU or US\$20M monthly revenue. \url{https://huggingface.co/Qwen/Qwen3.8-2.4T-A95B}
\bibitem{qwen38-scmp} South China Morning Post (13 Aug.\ 2026), ``Alibaba adds commercial restrictions to open-weight Qwen3.8-Max AI model.'' Note that several outlets reported the release as uniformly Apache 2.0; the licence file contradicts this for the flagship. \url{https://www.scmp.com/tech/tech-trends/article/3363927/alibaba-adds-commercial-restrictions-open-weight-qwen38-max-ai-model}
\bibitem{dsh-repo} DeepSeek, \emph{DeepSeek Harness} (\texttt{dsh}), MIT licence, developer preview v0.1, released 13 Aug.\ 2026; $\sim$100,200 GitHub stars and 9,400 forks within $\sim$48 hours. \url{https://github.com/deepseek-ai/deepseek-harness}
\bibitem{dsh-official} DeepSeek, DeepSeek Harness documentation: plugin-native architecture (models, tools, skills, sessions, sandboxes, filesystems, loops, orchestration and UI as components; ``no privileged core to patch''), four runtime modes, append-only session logging with resume/fork/replay. \url{https://deepseek.com/harness/en/}
\bibitem{dsh-mcp} DeepSeek Harness MCP client plugin documentation: stdio and streamable-http transports; tools namespaced \texttt{mcp\_\_<server>\_\_<tool>}; MCP \emph{Resources} and \emph{Prompts} not supported; image/audio payloads become placeholders. \url{https://deepseekdocs.com/en/docs/features/mcp}
\bibitem{cordis-paper} Y.~Shi, W.~Zhang and T.~Cui (Peking University; DeepSeek-AI), \emph{A Programming Paradigm for Spatiotemporal Composability}, draft of 13 Aug.\ 2026, 88 pp., self-hosted preprint marked ``under active revision.'' Formalizes revertible effects and reactive coeffects; implements them in Cordis. Cites OpenAI's ``Harness Engineering'' and Anthropic's ``Harness Design for Long-Running Application Development.'' \url{https://github.com/cordiverse/paper}
\bibitem{dsh-scmp} South China Morning Post (14 Aug.\ 2026): the harness launch marks a strategic pivot---``shifting the global competition from model intelligence to one that focuses on how seamlessly an AI agent can connect to real-world software.'' \url{https://www.scmp.com/tech/big-tech/article/3364082/}
\bibitem{dsh-newstack} The New Stack (13 Aug.\ 2026): DeepSeek Harness architecture, model-agnostic provider support (Anthropic, OpenAI, Bedrock, Azure, Gemini, DeepSeek), subagent bridges to Claude Code and OpenAI Codex; external pull requests not accepted. \url{https://thenewstack.io/deepseek-harness-open-source-plugins/}
\bibitem{dsh-venturebeat} VentureBeat (13--14 Aug.\ 2026), ``DeepSeek Harness launches as open-source rival to Claude Code, alongside V4-Pro on API with higher prices.'' \url{https://venturebeat.com/technology/deepseek-harness-launches-as-open-source-rival-to-claude-code-alongside-v4-pro-on-api-with-higher-prices}
\bibitem{dsh-36kr} 36Kr (14 Aug.\ 2026): Cordis originates as the kernel of Koishi, a four-year-old open-source chatbot framework with 4,000+ community plugins; the paper's first author is Koishi's creator. \url{https://eu.36kr.com/en/p/3938795963137411}
\bibitem{harness-effect} M.~Sayed Ali, A.~Novik, A.~Boddupally et al.\ (Writer, Inc.), \emph{The Harness Effect: How Orchestration Design Sets the Token Economics of Enterprise Agentic AI}, arXiv:2607.06906, Jul.\ 2026. 22 locked tasks, six models, five vendors: harness adoption saves 33--61\% holding the model fixed, against 36\% from switching between cheapest and most expensive model. Vendor-authored; n=22. \url{https://arxiv.org/html/2607.06906v1}
\bibitem{claude-code-anatomy} J.~Liu, X.~Zhao, X.~Shang and Z.~Shen (MBZUAI; UCL), \emph{Dive into Claude Code: The Design Space of Today's and Future AI Agent Systems}, arXiv:2604.14228, Apr.\ 2026. Anatomy of a production harness; reports the community 1.6\%/98.4\% code-share estimate as an estimate over a leak-derived bundle, not a measurement. \url{https://arxiv.org/html/2604.14228v1}
\bibitem{harness-register} Comparative harness reporting, 2026: convergent architecture across Claude Code, Codex CLI, Aider and OpenClaw; system-prompt budgets differing by nearly two orders of magnitude; Codex's one-line importer for Claude Code configuration; the duelling positions of B.~Cherny (Anthropic: the harness is ``the thinnest possible wrapper'') and K.~Prasad (Obvix Labs: ``the model is interchangeable---the harness is not''). \url{https://www.techtimes.com/articles/316928/20260521/}
\bibitem{octomind} Octomind, ``Why we no longer use LangChain for building our AI agents'': rigid abstractions, dynamic tool gating unavailable, debugging overhead---``when our team began spending as much time understanding and debugging LangChain as it did building features, it wasn't a good sign.'' No switching cost quantified. \url{https://octomind.dev/blog/why-we-no-longer-use-langchain-for-building-our-ai-agents}
\bibitem{mcp-aaif} Anthropic (9 Dec.\ 2025): donation of the Model Context Protocol to the Linux Foundation's Agentic AI Foundation, co-founded with Block and OpenAI; supporting members Google, Microsoft, AWS, Cloudflare, Bloomberg; 10,000+ active public MCP servers and 97M+ monthly SDK downloads at donation. \url{https://www.anthropic.com/news/donating-the-model-context-protocol-and-establishing-of-the-agentic-ai-foundation}
\bibitem{a2a-year} Linux Foundation (9 Apr.\ 2026): the Agent2Agent protocol passes 150+ organizations, 22,000+ GitHub stars and five production SDKs in its first year. \url{https://www.linuxfoundation.org/press/a2a-protocol-surpasses-150-organizations-lands-in-major-cloud-platforms-and-sees-enterprise-production-use-in-first-year}
\bibitem{agent-plugins} Agent Plugins 1.0 (6 Aug.\ 2026): vendor-neutral package format for distributing Agent Skills and MCP servers across clients; technical steering committee of core maintainers from Amazon, Cursor, Microsoft, OpenAI and Vercel, with Google joining the same day; independently governed. \url{https://developers.googleblog.com/agent-plugins-package-your-skills-tools-and-more/}
\bibitem{agent-plugins-newstack} The New Stack (8 Aug.\ 2026) and contemporaneous analysis of Agent Plugins 1.0: the spec standardizes packaging, not distribution---``if useful behavior ends up inside vendor namespaces, the standard will look open while lock-in just moves elsewhere'' (A.~Dalloul, EmpirioLabs); Anthropic, whose format the standard is based on, did not join. \url{https://thenewstack.io/agent-plugins-open-standard/}
\bibitem{cac-agents} Cyberspace Administration of China (17 Jul.\ 2026, WAIC), \emph{Global Cooperation Initiative on Agent Mutual Trust, Interconnection, and Interoperability}: calls for ``open, non-discriminatory, and transparent international interoperability standards for agents,'' agent identity systems and shared multi-agent infrastructure; names no existing protocol. \url{https://www.geopolitechs.org/p/chinas-cac-releases-global-cooperation}
\bibitem{opensandbox} OpenSandbox: agent execution sandbox released by Alibaba in March 2026, Apache 2.0, subsequently moved to a neutral organization with published governance, an OpenSSF best-practices badge and a CNCF landscape listing; gVisor/Kata/Firecracker isolation; listed integrations include Claude Code, Gemini CLI, LangGraph, Google ADK and Qwen Code. \url{https://github.com/opensandbox-group/OpenSandbox}
\bibitem{langchain-survey} LangChain, \emph{State of Agent Engineering} (n=1,340; fielded Nov.--Dec.\ 2025): 57\% have agents in production, 67\% among 10,000+ employee enterprises; over three-quarters run multiple models; the survey does not measure framework-switching rates. \url{https://www.langchain.com/state-of-agent-engineering}
\bibitem{grok46} SpaceXAI (xAI LLC), Grok 4.6 model card, rev.\ 12 Aug.\ 2026: closed weights, 500K context, \$2/\$6 per Mtok, same 1.5T base as Grok 4.5 with gains from post-training; Artificial Analysis Intelligence Index 61, tying GPT-5.6 Sol and two points below Claude Opus 5. \url{https://media.x.ai/v1/website/card-4p6-4cd2dc57.pdf}
\bibitem{glm53-early} the-decoder (14 Aug.\ 2026): Zhipu releases GLM-5.3 via API claiming the strongest open-weights coding model; weights withheld at announcement pending a security review, and published on 28 Aug.\ 2026; 743B MoE on the GLM-5.2 base; all benchmarks vendor-reported, and it trails GPT-5.6 Sol on Terminal-Bench 3.0 and DeepSWE v1.1. \url{https://the-decoder.com/zhipu-ai-releases-glm-5-3-claims-its-the-strongest-open-weights-coding-model/}
\bibitem{nemotron4} Techzine (12 Aug.\ 2026), relaying \emph{The Information}: NVIDIA reported to be in the final training phase of a trillion-parameter ``Nemotron 4'' intended to compete with the most powerful open-source models; NVIDIA states final training is not complete; licence unannounced; possible availability late autumn 2026. \url{https://www.techzine.eu/news/analytics/143552/nvidia-is-building-nemotron-4-with-at-least-1-trillion-parameters/}
\bibitem{gemini37} SiliconANGLE (13 Aug.\ 2026): Google launches Gemini 3.7 Flash for coding and agents at \$0.75 per Mtok input. \url{https://siliconangle.com/2026/08/13/google-launches-gemini-3-7-flash-coding-ai-agent-projects/}
\bibitem{deepseek-v4pro} Reuters via Yahoo Finance (13 Aug.\ 2026): DeepSeek takes V4-Pro to general availability, raises API prices for V4 Pro and V4 Flash, introduces peak/off-peak pricing; reported plans to at least double headcount and to expand private hiring of chip-design engineers toward an in-house accelerator reducing reliance on both NVIDIA and Huawei. \url{https://finance.yahoo.com/news/deepseek-releases-official-v4-pro-125608280.html}
\bibitem{nvidia-rubin-pod} NVIDIA, ``NVIDIA Vera Rubin POD: seven chips, five rack-scale systems, one AI supercomputer'' (16 Mar.\ 2026, mod.\ 4 Aug.\ 2026): Rubin in full production 2H~2026; 288\,GB HBM4, $\sim$13\,TB/s, NVLink~6. \url{https://developer.nvidia.com/blog/nvidia-vera-rubin-pod-seven-chips-five-rack-scale-systems-one-ai-supercomputer/}
\bibitem{gtc2026-roadmap} Tom's Hardware (17 Mar.\ 2026), reporting the GTC~2026 keynote: Feynman in 2028 with the Rosa CPU, LP40 inference processor, BlueField-5 and CX10, and Kyber racks in both copper and co-packaged-optical scale-up; NVIDIA's headline metrics stated as tokens per watt and cost per token rather than FLOPS. \url{https://www.tomshardware.com/pc-components/gpus/nvidia-updates-data-center-roadmap-with-rosa-cpu-and-stacked-feynman-gpus-optical-nvlink-groq-lpus-with-nvfp4-and-nvlink-also-on-deck}
\bibitem{kyber-slip} CNBC and Tom's Hardware (6 Jul.\ 2026), reporting SemiAnalysis: the Kyber NVL144 rack for Rubin Ultra slips from 2027 to 2028 on a 78-layer orthogonal-midplane signal-integrity problem; the NVL72$\times$2 stopgap cancelled after customer pushback; NVIDIA responds only that ``our roadmap is intact.'' \url{https://www.cnbc.com/2026/07/06/nvidia-kyber-rack-system-delays-manufacturing-taiwan-rubin-chips-.html}
\bibitem{semianalysis-gtc} SemiAnalysis (24 Mar.\ 2026), GTC~2026 analysis: attention--FFN disaggregation, KV-cache-bounded attention utilization, and new storage tiers treating KV cache as a first-order constraint. \url{https://newsletter.semianalysis.com/p/nvidia-the-inference-kingdom-expands}
\bibitem{hbm-gap} Grunewald and Akula, AI Frontiers (2 Feb.\ 2026), on HBM as the binding constraint under export controls, with CXMT at HBM2-class and early HBM3 sampling; read together with SemiAnalysis on the Ascend ramp being memory-bound rather than logic-bound. Note that die counts and stack counts differ by roughly an order of magnitude and must not be compared directly. \url{https://ai-frontiers.org/articles/high-bandwidth-memory-critical-gaps-us-export-controls}


% ============ AUGUST 2026: FINANCING PLATFORMS AND THE COMMODITY TURN ============
\bibitem{nvidia-500b-analysis} Analyst reaction to the financing platforms, 11--12 Aug.\ 2026: F.~Wang (Hedgeye)---``in effect, they made Nvidia's product cheaper without really cutting GPU prices''; J.~Moore (Morgan Stanley) on third-party capital alleviating circularity; UBS and Seaport on circular financing; the telecom comparison (nine suppliers, $\sim$\$25.6B of vendor financing by end-2000; 24 of the 30 largest listed carriers bankrupt) with the disanalogy that Lucent and Nortel held their paper on balance sheet. \url{https://fortune.com/2026/08/11/nvidia-500-billion-ai-financing/}
\bibitem{nvidia-backstop-cut} Wall Street Journal via Reuters (14 Aug.\ 2026): NVIDIA scales its contemplated backstop of OpenAI's Ohio campus back from roughly \$250B to under \$120B, covering only the first phase, after investors raised concerns about risk exposure tied to large financing commitments. \url{https://finance.yahoo.com/technology/ai/articles/nvidia-scales-back-250-billion-234356524.html}
\bibitem{nvidia-backstop-floor} Reporting on NVIDIA's neocloud utilization backstops (from $\sim$1 Jul.\ 2026): a guaranteed utilization floor under deployed GPUs, with NVIDIA renting idle capacity at a predetermined rate in exchange for a share of cloud revenue; disclosed arrangements include Sharon AI (40,000 GB300s) and Firmus Technologies (up to 170,000 GPUs). Rate and revenue share undisclosed; NVIDIA declined to comment. \url{https://capacityglobal.com/news/nvidia-now-agrees-to-rent-back-unused-gpu-capacity-from-neocloud-operators-if-customer-demand-falls-short/}
\bibitem{nvidia-allocation} M.~Kirkwood, Tech Policy Press (18 Dec.\ 2025), on NVIDIA allocation practices, the UK CMA's identification of an interconnected web of over 90 partnerships and strategic investments, and questions over preferential supply and pricing. \url{https://www.techpolicy.press/nvidia-is-building-a-shield-of-concentrated-power/}
\bibitem{nvidia-ncp} NVIDIA Cloud Partner requirements: partners must expose newly delivered capacity through a centralized API and report delivered, healthy, reserved and in-use accelerators per tenant, with strict logical and physical isolation of NVIDIA-reserved capacity. No prohibition on resale or subletting of data-centre capacity is stated. \url{https://docs.nvidia.com/dsx/ncp/nvidia-requirements-for-ai-clouds/home\#capacity-and-fleet-management}
\bibitem{openai-exits} Coverage of OpenAI's August 2026 executive departures: B.~Lightcap (COO 2022--spring 2026, latterly special projects) announced his departure 11 Aug.; D.~Dresser, first chief revenue officer, announced hers 13 Aug.\ after $\sim$8 months, with D.~Rajic named successor; CFO S.~Friar briefed investors 14 Aug.\ reporting an annualized run rate near \$40B and enterprise revenue overtaking consumer. \url{https://techcrunch.com/2026/08/11/brad-lightcap-openais-longtime-coo-is-leaving-to-start-something-new/}
\bibitem{openai-s1} OpenAI publicly disclosed on 8 Jun.\ 2026 that it had confidentially submitted a draft registration statement (no public S-1 filed; no listing date set): ``We recently submitted a confidential S-1. We expect it to leak so we're just announcing it. We have not decided on timing yet.'' Valuation cited at $\sim$\$852B. \url{https://www.cbsnews.com/news/openai-files-confidential-initial-public-offering/}
\bibitem{cme-compute} CME Group and Silicon Data (11 Aug.\ 2026): compute futures to launch 5 Oct.\ 2026, listed on NYMEX, each contract representing one month of GPU rental cost referencing H100 and B200. P.~Keavey (CME): ``Compute has become the currency of the AI age.'' \url{https://www.cmegroup.com/media-room/press-releases/2026/8/11/cme_group_and_silicondatatolaunchcomputefuturesonoctober5tounloc.html}
\bibitem{ice-ornn} Intercontinental Exchange and Ornn (19 May 2026): cash-settled GPU compute futures on the Ornn Compute Price Index, built exclusively from printed transactions across H100, H200, B200 and RTX 5090. \url{https://ir.theice.com/press/news-details/2026/ICE-and-Ornn-to-Launch-GPU-Compute-Futures-Contracts/default.aspx}
\bibitem{ice-nativx} ICE and NATIVX (1 Jul.\ 2026): energy-normalized compute futures on the COIL index, which ``strips out the noise introduced by regional power cost disparities,'' with sub-indices COIL-T (training), COIL-I (inference), COIL-G (graphics) and COIL-CO (connectivity). \url{https://www.businesswire.com/news/home/20260701837320/en/ICE-and-NATIVX-to-Launch-Energy-Normalized-Compute-Futures-Contracts}
\bibitem{kalshi-compute} Kalshi (14 Jul.\ 2026): compute forward curves for B200, H200 and A100 derived from its own weekly and monthly GPU price contracts. T.~Mansour: ``Compute is the new oil. Like every commodity before it, it needs a real derivatives market.'' Volume and open interest undisclosed. \url{https://news.kalshi.com/p/compute-forward-curves}
\bibitem{architect-futures} Architect Financial Technologies and Ornn (21 Jan.\ 2026): perpetual futures on GPU and DRAM prices. B.~Harrison: ``there is an urgent cross-industry need to establish standardized derivatives contracts and centralized order books for compute.'' \url{https://www.prnewswire.com/news-releases/architect-financial-technologies-partners-with-compute-index-provider-ornn-to-launch-exchange-traded-futures-on-gpu-and-ram-prices-302666613.html}
\bibitem{silicon-data} Silicon Data: daily GPU rental indices since 20 May 2025 built on $\sim$3.5 million pricing observations, distributed via Bloomberg (ticker SDH100RT); observed dispersion of $\sim$\$2.74/hr (neocloud) against $\sim$\$7.20/hr (hyperscaler) for H100; an LLM token-expenditure index published alongside; forward curves in backwardation (B200 spot $\sim$\$5.62 vs $\sim$\$5.17 at 36 months). Index-provider data, self-interested. \url{https://www.silicondata.com/products/silicon-index}
\bibitem{compute-commodity} D.~Thomas, Benzinga (13 Aug.\ 2026), reporting analysis by R.~Goswami (Semafor) on the CME compute-futures launch: the steel-benchmark rebuttal to the heterogeneity objection, and the market-structure objection---``there's no other game in town except for Nvidia,'' an oil market ``run entirely by Standard Oil''; NVIDIA ``has little incentive to encourage greater price transparency or falling compute costs.'' \url{https://finance.yahoo.com/technology/ai/articles/why-nvidia-may-stand-way-153004691.html}
\bibitem{variant-preconditions} C.~Shack and A.~Levin, Variant Fund (14 May 2026): five preconditions for a compute futures market, of which supply-side fragmentation and a standardized tradable unit are absent; the top four hyperscalers hold $\sim$78\% of global self-built critical IT power and 69\% of H100 supply; inference is more fungible than training and projected above 65\% of AI compute demand by 2029. \url{https://variant.fund/articles/compute-as-commodity-framework-gpu-futures-markets-emerge/}
\bibitem{friedman-basis} D.~Friedman (17 Jul.\ 2026), ``Compute futures won't make compute fungible'': futures ``make its non-fungibility legible''---``the benchmark doesn't eliminate differences within the physical market. It gives those differences something to be priced against.'' Five basis categories: hardware, configuration/topology, location, contractual term, counterparty. \url{https://davefriedman.substack.com/p/compute-futures-wont-make-compute}
\bibitem{china-compute-exchanges} China's computing-power trading platforms: Shanghai (trial operation 26 Apr.\ 2023, via SHIXP) and Yinchuan, Ningxia (24 Feb.\ 2023, co-built with the Beijing International Big Data Exchange; 27 participating institutions including Huawei, Alibaba Cloud, SenseTime, Baidu); CAICT's national China Computing Power Platform under MIIT with 1,000+ registered enterprise users, 100+ service providers and 10 regional sub-platforms. None publishes standardized contracts, clearing, or transaction volumes. \url{https://www.yicaiglobal.com/news/china-debuts-first-computing-power-trading-hub}
\bibitem{china-compute-network} People's Daily (5 Jun.\ 2026): 1.37 million PFLOPS of intelligent computing under national monitoring, $\sim$72\% of the national total, with 80\%+ of intelligent computing inside the eight national hubs; the 15th Five-Year Plan extends East Data West Computing; standardized interconnection targeted for 2028. \url{https://en.people.cn/n3/2026/0605/c90000-20464201.html}
\bibitem{china-tokens} TechNode (17 Apr.\ 2026), reporting Mao Shengyong, deputy head of the National Bureau of Statistics, at a State Council Information Office briefing: China's daily token usage exceeded 140 trillion in March 2026, up over 40\% from end-2025. \url{https://technode.com/2026/04/17/china-authority-says-daily-ai-token-usage-exceeds-140-trillion-in-march-up-over-40-vs-end-2025/}
\bibitem{token-prices} Comparative token pricing, mid-2026: DeepSeek V4-Pro at \$0.435/\$0.87 per Mtok against GPT-5.6 Sol at \$5/\$30 and Claude Opus 4.8 at \$5/\$25 ($\sim$34$\times$ on output); DeepSeek V4-Flash at \$0.14/\$0.28; Qwen3-8B at \$0.01/\$0.01. \url{https://www.cgphbanquedaffaires.com/post/china-has-found-its-next-commodity-to-export-the-token}
\bibitem{inference-price-floor} South China Morning Post (10 Aug.\ 2026), citing Jefferies and Silicon Data: blended enterprise inference cost fell to \$1.16--1.18 per million tokens on 6--8 Aug.\ 2026, a 2026 low, from \$1.45 in late July and \$2.04 on 31 May, driven by price wars and Chinese open-source models. \url{https://www.scmp.com/tech/tech-trends/article/3363549/enterprise-ai-costs-hit-2026-low-driven-price-wars-chinese-open-source-models-research}
\bibitem{china-token-export} China Daily (1 Jun.\ 2026), ``Tokens are China's next goldmine'': the export model sells inference by API while electricity, servers and models remain domestic. Liu Xiaoge, State Council Development Research Center: token exports ``provide a new way to offset risks of traditional trade frictions and drive growth in services trade.'' \url{https://global.chinadaily.com.cn/a/202606/01/WS6a1cc2c4a310d6866eb4baae.html}
\bibitem{china-idle} Reuters (Jul.\ 2025) and ASPI Strategist (6 May 2026): utilization at Chinese AI data centres of 20--30\% against a 60\% target; 633 hyperscale facilities built under East Data West Computing by 2024; 100+ projects scrapped in 18 months; western-region capacity failing latency, fibre and talent tests for inference; NVIDIA/CUDA and Ascend/CANN estates that cannot be integrated into a single seamless cloud environment. \url{https://www.aspistrategist.org.au/abundant-electricity-isnt-enough-chinas-overbuilt-ai-computing-power-is-underused/}
\bibitem{volume-revenue} Baiguan (9 Jul.\ 2026), on neutral-router economics: Chinese models at roughly 61\% of token consumption while Anthropic holds $\sim$12\% of volume and captures $\sim$46\% of revenue; large enterprises still overwhelmingly buy from Anthropic, OpenAI, Azure and Google Cloud; China's 2017 National Intelligence Law as a persistent procurement obstacle. \url{https://www.baiguan.news/p/china-ai-token-export-openrouter-market-share-chinese-models-performance-gap-cost-efficiency-semiconductor-energy-self-sufficiency-investment-opportunities}
\bibitem{value-capture} SemiAnalysis, ``AI Value Capture'': inference gross margins rising from under 40\% to over 70\% across 2025--26 while the compute layer commoditized, with the surplus concentrating at the model layer; neoclouds earning close to their minimum IRR hurdle at prevailing rental rates. \url{https://newsletter.semianalysis.com/p/ai-value-capture-the-shift-to-model}


% ============ LOCALITY, LATENCY, AND THE MEMORY-SCARCITY MODEL ============
\bibitem{matsuoka-memory} S.~Matsuoka, \emph{Memory Scarcity, Open Models, and the Restructuring of the AI Industry, 2026--2030: A quantitative scenario analysis of inference economics, training-cost divergence, and infrastructure solvency}, arXiv:2607.07207, 8 Jul.\ 2026. Entrant--incumbent cost gap 3.2$\times$ (2026), 1.9$\times$ (2027), re-widening to 3--4$\times$ by 2029--30 via a depreciation conveyor; training bifurcating into an \$18--38B frontier tier and a $\sim$\$5M distillation tier by 2030; solvency corridor requiring $\sim$2$\times$ annual token-demand growth for four years; scenarios Rotating Landlord Oligopoly 25\%, Commoditization Crash 25\%, Jevons Absorption 20\%, System-Layer Re-differentiation 18\%, Geopolitical Bifurcation 12\%. Author's own companion analysis. \url{https://arxiv.org/abs/2607.07207}
\bibitem{toms-latency} Tom's Hardware (24 Jul.\ 2025), the proximate English source for the ``20\,ms'' claim: ``Beijing established a target of a network latency of 20 milliseconds, which is necessary for real-time applications like financial services or high-frequency trading. However, many western-region data centers fail to meet this benchmark.'' Note the internal inconsistency: the primary Chinese document assigns high-frequency trading to a separate 10\,ms tier. \url{https://www.tomshardware.com/desktops/servers/china-is-developing-nation-spanning-network-to-sell-surplus-data-center-compute-power-latency-disparate-hardware-are-key-hurdles}
\bibitem{ndrc-2021} National Development and Reform Commission et al.\ (May 2021), \emph{Implementation Plan for Computing Hubs in the National Integrated Big Data Centre Collaborative Innovation System}: end-to-end network latency ``in principle within 20 milliseconds'' \emph{one-way} (danxiang) for the eastern hub clusters, serving industrial internet, finance and securities, disaster warning, telemedicine, video calling and AI inference; a separate 10\,ms tier for city-edge facilities serving high-frequency trading, VR/AR, ultra-high-definition video and connected vehicles; western nodes assigned to ``back-end processing, offline analysis, storage backup''---non-real-time demand. Predates ChatGPT by eighteen months.
\bibitem{nda-latency} State Council of the People's Republic of China (29 Aug.\ 2024), quoting Liu Liehong, head of the National Data Administration: ``Network latency between the eastern and western hubs has generally met the 20-millisecond (ms) requirement.'' \url{https://english.www.gov.cn/news/202408/29/content_WS66d03a1ac6d0868f4e8ea539.html}
\bibitem{china-latency-measured} Published intra-China hub latencies: Gui'an to Chengdu/Chongqing 6\,ms, to Guangzhou/Shenzhen 10\,ms, to the Yangtze Delta and Beijing-Tianjin-Hebei within 18\,ms; Zhongwei (Ningxia) to Beijing 8--10\,ms one-way and to Shanghai within 15\,ms one-way, on 18\,Tbps of export bandwidth with direct fibre to 26 cities; Inner Mongolia to Beijing-Tianjin-Hebei within 5\,ms. Chinese official and state-adjacent sources; no independent Western measurement of intra-China backbone latency was located. \url{https://news.qq.com/rain/a/20250217A07J1N00}
\bibitem{m890-ulanqab} Global Times and TechNode (12 Aug.\ 2026): Alibaba Cloud's Lingjun M890 supernode enters commercial service with first deployment in Ulanqab, Inner Mongolia, designed for \emph{inference} on mixture-of-experts models up to ten trillion parameters, already serving live commercial traffic for Kimi K3 and Qwen3.8-Max. \url{https://www.globaltimes.cn/page/202608/1368093.shtml}
\bibitem{alibaba-endpoints} Alibaba Cloud Model Studio: published public endpoints in Beijing, Hong Kong, Singapore, Frankfurt, Tokyo and US (Virginia). \url{https://www.alibabacloud.com/help/en/model-studio/regions/}
\bibitem{ttft-budget} Spheron, \emph{LLM Inference SLO Engineering: TTFT, ITL, and P99 Latency Budgets} (15 May 2026): a 400\,ms P99 TTFT budget allocating 20\,ms to network ingress, 50\,ms to queueing, 280\,ms to prefill, 10\,ms to tokenization, 15\,ms to gateway and 25\,ms to safety filtering; inline code completion at 100\,ms P99 TTFT (``must feel local to IDE''); tightening a target from 500\,ms to 200\,ms raising infrastructure spend $\sim$35\%. \url{https://www.spheron.network/blog/llm-inference-slo-ttft-itl-latency-budget-guide-2026/}
\bibitem{ttft-measured} Digital Applied, \emph{AI Model Latency Benchmarks 2026} (23 Apr.\ 2026), 10,000 samples per model and provider at 1,024-token input: GPT-5.5 TTFT P50 1.12\,s at 92 output tok/s; Claude Opus 4.7 0.85\,s; a reasoning model at medium effort 8.4\,s; extended thinking 28\,s; specialist inference hardware 0.16--0.18\,s. \url{https://www.digitalapplied.com/blog/ai-model-latency-benchmarks-2026-ttft-throughput}
\bibitem{rtt-matrix} Measured AWS inter-region round-trip latency: Tokyo to Oregon 99\,ms; Tokyo to Northern Virginia 149\,ms; Seoul to Oregon 120\,ms; Singapore to Oregon 173\,ms; Frankfurt to Northern Virginia 92\,ms. \url{https://latency.bluegoat.net/}
\bibitem{openai-residency} OpenAI Help Center, \emph{Data residency and inference residency} (updated Aug.\ 2026): inference residency---``model inference (GPU execution) on in-scope customer content happens exclusively in-region''---available for Europe (EEA and Switzerland), the United States and the United Arab Emirates only. Australia, Canada, India, Japan, Singapore, South Korea and the United Kingdom receive data residency only. \url{https://help.openai.com/en/articles/9903489-data-residency-and-inference-residency-for-chatgpt}
\bibitem{anthropic-residency} Anthropic platform documentation, data residency: inference geography defaults to \texttt{global}---``inference may run in any available geography for optimal performance and availability''---with only \texttt{us} and \texttt{global} available. \url{https://platform.claude.com/docs/en/manage-claude/data-residency}
\bibitem{azure-global} Microsoft Learn, deployment types in Foundry Models: Global Standard ``use[s] Azure's global infrastructure to dynamically route traffic to available datacenters''; ``For most workloads, start with Global Standard\dots\ has the lowest price.'' \url{https://learn.microsoft.com/en-us/azure/foundry/foundry-models/concepts/deployment-types}
\bibitem{vertex-global} Google Cloud (28 Jul.\ 2025), global endpoint for Claude models on Vertex AI: ``dynamically route your requests to any region with available capacity''; latency ``may be higher in some cases''; customers with data-residency requirements should use regional endpoints because the global endpoint ``does not guarantee that requests will be processed in any specific location.'' \url{https://cloud.google.com/blog/products/ai-machine-learning/global-endpoint-for-claude-models-generally-available-on-vertex-ai}
\bibitem{bedrock-cris} AWS (24 Feb.\ 2026): global cross-Region inference for Claude models with source regions Singapore, Jakarta, Taipei, Thailand and Malaysia, routing to more than twenty commercial regions worldwide while data at rest remains in the source region. \url{https://aws.amazon.com/blogs/machine-learning/global-cross-region-inference-for-latest-anthropic-claude-opus-sonnet-and-haiku-models-on-amazon-bedrock-in-thailand-malaysia-singapore-indonesia-and-taiwan}
\bibitem{gorgo} Toniolo, Dinesh and Thorstenson, \emph{GORGO: Maximizing KV-Cache Reuse While Minimizing Network Latency in Cross-Region LLM Load Balancing}, arXiv:2602.11688, Feb.\ 2026: measured inter-region latencies of 281\,ms (US West to Germany) and 3\,ms (intra-region), with a worked case in which routing to Germany produced a lower total TTFT (296\,ms) than serving locally under load (305\,ms); 2.5$\times$ faster median TTFT than least-load routing. \url{https://arxiv.org/html/2602.11688v1}
\bibitem{prefill-service} R.~Qin, W.~He, Y.~Wang et al.\ (Moonshot AI; Tsinghua University), \emph{Prefill-as-a-Service: KVCache of Next-Generation Models Could Go Cross-Datacenter}, arXiv:2604.15039, 16 Apr.\ 2026: long-context prefill offloaded to standalone clusters with KV cache transferred over commodity Ethernet; $\sim$170\,Gbps for a 1T-parameter model at 32K context, below 100\,Gbps at 128K; 54\% higher serving throughput than homogeneous prefill--decode. \url{https://arxiv.org/html/2604.15039v1}
\bibitem{kv-internet} S.~Ray, N.~Feamster and J.~Jiang (University of Chicago), \emph{An Internet for the KV Cache}, arXiv:2608.01526, 2 Aug.\ 2026: transferring a 1\,GB KV cache over a 10\,Gbps link costs $\sim$0.8\,s and $\sim$\$0.09 against $\sim$15\,s and \$0.47 to recompute on 8$\times$B200---19$\times$ faster and 5.2$\times$ cheaper to ship than to recompute. \url{https://arxiv.org/html/2608.01526v1}
\bibitem{distserve} Zhong, Liu, Chen et al., \emph{DistServe: Disaggregating Prefill and Decoding for Goodput-Optimized Large Language Model Serving}, arXiv:2401.09670 (OSDI '24). Table 1 service-level objectives: code completion TTFT 125\,ms; chatbot 250\,ms--2.5\,s; summarization 15\,s. KV-cache transmission under 0.1\% of total latency. \url{https://arxiv.org/pdf/2401.09670}
\bibitem{voice-latency} Published real-time voice latency budgets (2026): $\sim$500--600\,ms total for conversational speech, of which 30--80\,ms is network; measured production voice-to-voice medians of 680--850\,ms; a co-located stack reducing network and SIP from 150\,ms to 45\,ms and total from 1,210\,ms to 450\,ms. \url{https://thepromptbench.com/voice-and-realtime/latency-budgets-for-realtime-voice/}
\bibitem{prompt-caching} Anthropic (13 Mar.\ 2025): prompt caching reduces ``costs by up to 90\% and latency by up to 85\% for long prompts.'' OpenAI's documentation notes that caching works by routing to servers that recently processed the same prompt---creating region stickiness independent of propagation delay. \url{https://claude.com/blog/token-saving-updates}
\bibitem{equinix-rag} Equinix (30 Apr.\ 2025), on retrieval-augmented generation: the inference platform must be ``tightly coupled to the data platform'' so that retrieval reaches enterprise data over low-latency connectivity---an argument about proximity to the corpus, not to the user. \url{https://blog.equinix.com/blog/2025/04/30/what-is-retrieval-augmented-generation-rag-and-where-should-you-do-it/}
\bibitem{sina-inference} Sina Finance (14 Aug.\ 2026), ``China's new computing map: training migrates west, inference sinks east'': inference-centric compute moves east because the eastern regions are the main point of explosion of AI inference \emph{demand}. The article states no latency figure. \url{https://finance.sina.com.cn/roll/2026-08-14/doc-ininhhrn4148288.shtml}


% ============ MEMORY: PRICES, HBF, ZAM, AND THE LOCAL-INFERENCE WALL ============
\bibitem{trendforce-dram} TrendForce contract-price series (2026): conventional DRAM +93--98\% QoQ in 1Q26, +58--63\% in 2Q26, +13--18\% in 3Q26, following $\sim$172\% across 2025; DRAM industry revenue \$97B in 1Q26, +81\% QoQ. \url{https://www.trendforce.com/presscenter/news/20260601-13070.html}
\bibitem{dram-street} Street and spot memory pricing, 2026: 32\,GB DDR5-6000 kits from $\sim$\$85 (mid-2025) to \$369--432; 64\,GB DDR4-2933 ECC RDIMMs from $\sim$EUR~98 to EUR~500 used and EUR~835 new; DDR5 spot +307\% Sep--Nov 2025. Trackers diverge by up to 50\%; treat as a range. \url{https://capitalandcompute.net/memory-prices/ddr5/}
\bibitem{hbm-pricing} Seoul Economic Daily (12 Jul.\ 2026): HBM3E at $\sim$\$1.5--1.6 per Gb and HBM4 at $\sim$\$2 per Gb in 2H26, forecast \$4--5+ per Gb in 2027; HBM4 taking 4--6 months to produce against 3--3.5 for DDR5, at low initial yields. \url{https://en.sedaily.com/finance/2026/07/12/hbm4-prices-to-double-next-year-as-samsung-sk-hynix-keep}
\bibitem{mac-studio-512} Tom's Hardware and MacRumors (5--6 Mar.\ 2026): Apple withdraws the 512\,GB Mac Studio configuration and raises the 96\,GB-to-256\,GB upgrade from \$1,600 to \$2,000; 256\,GB deliveries slip to May. \url{https://www.tomshardware.com/tech-industry/apple-pulls-512-mac-studio-upgrade-option}
\bibitem{kimi-k3-card} Moonshot AI, Kimi K3 model card: 2.8T total / 104B active, 16 of 896 experts routed per token, native MXFP4 from quantization-aware fine-tuning; published download $\sim$1.4\,TB against $\sim$5.6\,TB at FP16. Weights released 27 Jul.\ 2026 under the custom Kimi K3 licence. \url{https://huggingface.co/moonshotai/Kimi-K3}
\bibitem{deepseek-hw} DeepSeek-AI, \emph{Insights into DeepSeek-V3: Scaling Challenges and Reflections on Hardware for AI Architectures}, arXiv:2505.09343, May 2025: KV cache per token of 70.272\,KB for DeepSeek-V3 multi-head latent attention against 327.680\,KB for Qwen-2.5 72B and 516.096\,KB for Llama-3.1 405B (both grouped-query); memory demand growing $>$1000\% per year against HBM capacity under 50\%. \url{https://arxiv.org/html/2505.09343v1}
\bibitem{hbf-ocp} SanDisk and SK hynix (3 Aug.\ 2026): first Open Compute Project technical specification for High Bandwidth Flash, following the August 2025 standardization memorandum and a February 2026 consortium launch. Consortium members SanDisk, SK hynix, Google and Tenstorrent; NVIDIA, AMD, Samsung, Micron and Kioxia absent. \url{https://www.sandisk.com/company/newsroom/press-releases/2026/2026-08-03-Sandisk-and-sk-hynix-advance-global-standardization-of-hbf}
\bibitem{hbf-spec} HBF specification as reported at FMS 2026: up to 512\,GB per stack in 8-Hi and 16-Hi NAND; three bandwidth grades from $\sim$0.4 to 3.0\,TB/s; UCIe host interface at up to 64\,GT/s over 64 lanes; NAND dies stacked with TSVs in an HBM4-compatible footprint with a CBA logic die scheduling parallel sub-array reads. Evaluation samples to selected customers promised by end-2026. \url{https://www.tomshardware.com/pc-components/ssds/sandisk-and-sk-hynix-unveil-hbf-spec-up-to-16-hi-nand-stacks-3-tb-s-bandwidth-ucie}
\bibitem{hbf-critique} Contemporaneous criticism of the HBF specification (Aug.\ 2026): the OCP document carries no version number, no conformance procedure and no certification scheme, and omits latency, power consumption, write speed, error correction, endurance and cost; the stated bandwidth and capacity figures are the two vendors' development targets rather than conformance requirements. \url{https://www.igorslab.de/en/512-gb-up-to-3-tb-s-hbf-closes-gap-between-hbm-and-ssd/}
\bibitem{hbf-analysis} Chipstrat, \emph{High Bandwidth Flash: the full report} (7 Jul.\ 2026): HBF latency in microseconds against HBM4's $\sim$100\,ns; minimum read granularity $\sim$128$\times$ larger; power $\sim$2$\times$ HBM4; cost per gigabyte $\sim$10$\times$ lower but cost per \emph{stack} comparable because advanced packaging dominates; suitable for weights, unsuitable for KV cache, activations or training. \url{https://www.chipstrat.com/p/high-bandwidth-flash-the-full-report}
\bibitem{saimemory} SoftBank Corp.\ (3 Feb.\ 2026) and SAIMEMORY Corporation (22 Apr.\ 2026): SAIMEMORY incorporated December 2024 as a SoftBank subsidiary, joint venture with Intel; Series A from Fujitsu, the Development Bank of Japan, RIKEN and SoftBank; NEDO grant under the Post-5G / advanced semiconductor programme; prototype targeted for the fiscal year ending March 2028, commercialization FY2029. \url{https://www.softbank.jp/en/corp/news/press/sbkk/2026/20260203_01/}
\bibitem{zam-tech} ZAM (Z-Angle Memory): stacked DRAM built on Intel's Next Generation DRAM Bonding, validated under a US Department of Energy Advanced Memory Technology programme; a ``via-in-one'' TSV replacing HBM's many TSVs to free die area; claimed $\sim$40--50\% power reduction against comparable HBM and a target of up to 512\,GB per chip. \url{https://www.blocksandfiles.com/ai-ml/2026/02/04/intel-and-softbank-aim-zam-at-hbm/4090386}
\bibitem{zam-vlsi} Powerchip, Intel and SAIMEMORY at the VLSI Symposium (Jun.\ 2026): functional validation of a nine-layer fusion-bonded DRAM stack; $\sim$0.25\,Tb/s/mm$^2$ bandwidth density, $<$0.7\,pJ/bit data-movement energy, $\sim$3\,\textmu m substrates, $\sim$13.7k TSVs per layer, 171\,mm$^2$ footprint at $\sim$1.125\,GB per layer---i.e.\ a demonstrated module of roughly 10\,GB against the 512\,GB target. \url{https://www.trendforce.com/news/2026/05/22/news-psmc-joins-intel-saimemory-to-demo-9-layer-fusion-bonded-via-in-one-architecture-for-high-bandwidth-3d-memory/}
\bibitem{lpddr6} JEDEC JESD209-6 (LPDDR6), published 9 Jul.\ 2025; Samsung has sampled LPDDR6X to Qualcomm for the AI250, with commercial production indicated for 2H27. The extension roadmap is under JEDEC standardization (see \cite{jedec-lpddr6-roadmap}).
\bibitem{qualcomm-ai200} Qualcomm (27 Oct.\ 2025): AI200 with 768\,GB of LPDDR per card, direct liquid cooling, 160\,kW racks, for 2026; AI250 adding near-memory compute with a claimed $>$10$\times$ effective bandwidth, for 2027. \url{https://www.qualcomm.com/news/releases/2026/qualcomm-ai200-ai250}
\bibitem{ssd-energy} Kyung, Yun and Ahn (Seoul National University), \emph{SSD Offloading for LLM Mixture-of-Experts Weights Considered Harmful in Energy Efficiency}, arXiv:2508.06978, Aug.\ 2025: flash reads at $\sim$102.4\,pJ/bit against DRAM's $\sim$4.2; SSD offloading raising per-token energy 3.8--12.5$\times$ for Mixtral and 4.7--9.8$\times$ for DeepSeek-R1, with SSD access exceeding 73--80\% of total energy; prefetching hides latency (1.25--1.32$\times$) but cannot hide energy; flash requires $\sim$10$\times$ improvement to $\sim$10\,pJ/bit to compete. \url{https://arxiv.org/html/2508.06978v1}
\bibitem{skhynix-shortage} SK hynix chief executive, 11 Jul.\ 2026, on the day of the company's Nasdaq listing: ``We forecast that next year will be the worst year in the industry's history from the supply perspective,'' and ``we still forecast that customer demand will remain higher than our supply capacity even beyond 2030.'' \url{https://www.tomshardware.com/pc-components/dram/sk-hynix-says-2027-will-be-the-worst-year-for-memory-shortage-forecasts-crunch-to-last-until-2030-ceo-shares-grim-outlook-on-the-day-sk-hynix-gets-listed-on-nasdaq}
\bibitem{nand-surplus} TrendForce (21 Jul.\ 2026): NAND flash in a 4--5\% supply deficit for 2026 with the balance turning positive in 2H27, against server DRAM contract prices forecast to rise every quarter from 2H26 through 2H27 on RDIMM bit-supply growth of only 15--20\% year over year. \url{https://www.trendforce.com/presscenter/news/20260721-13148.html}
\bibitem{cpx-cancel} NVIDIA Rubin CPX---announced 9 Sep.\ 2025 as a capacity-first inference GPU with 128\,GB of GDDR7 rather than HBM---removed from the roadmap at GTC 2026, with supply-chain reporting indicating no memory or substrate orders placed and substrate suppliers treating it as effectively cancelled; NVIDIA indicating the idea may return with Feynman. \url{https://wccftech.com/nvidia-rubin-cpx-is-off-the-roadmap-but-may-return-with-feynman-in-2028/}
\bibitem{groq-deal} NVIDIA acquires Groq's assets and acquihires its leadership, announced 25 Dec.\ 2025 (initial reporting described the transaction as a technology licence plus hiring; later trade coverage and NVIDIA's own product integration describe an asset purchase with an acquihire, structured to avoid extended regulatory review); reported at $\sim$\$20B, the company's largest transaction on record. Groq's architecture is SRAM-based---roughly 500\,MB per chip at $\sim$150\,TB/s---i.e.\ the opposite of capacity-first. \url{https://www.datacenterdynamics.com/en/news/nvidia-to-license-tech-from-ai-inference-chip-company-groq-hire-its-leadership/}
\bibitem{menlo-openweight} Menlo Ventures enterprise LLM market surveys: open-weight share of enterprise workloads falling from 19\% (late 2024) to 13\% (mid-2025) to $\sim$11\% (2026), while enterprise LLM API spend rose from \$3.5B to \$8.4B; constraints reported as licensing terms, governance and provenance, support and total cost of ownership rather than hardware. \url{https://menlovc.com/perspective/2025-mid-year-llm-market-update/}
\bibitem{ymtc-nand} Counterpoint Memory \& Storage Tracker (12 Aug.\ 2026): YMTC enters the top three NAND makers by shipment share for the first time at 14\%, displacing Kioxia, behind Samsung (25\%) and SK hynix (22\%); enterprise SSDs absorbing 48\% of NAND bits in Q2 2026, up from 26\% a year earlier. YMTC's Xtacking is wafer-to-wafer hybrid bonding, in mass production since 2018---the process HBF requires. \url{https://www.tomshardware.com/tech-industry/ymtc-breaks-into-the-top-three-nand-makers-for-the-first-time}


\bibitem{jedec-lpddr6-roadmap} JEDEC (22 Apr.\ 2026), ``JEDEC Previews LPDDR6 Roadmap Expanding LPDDR into Data Centers and Processing-in-Memory'': narrower per-die interface enabling higher capacities, ``512\,GB density on the horizon,'' an LPDDR6 PIM standard in development, and an ``LPDDR6-based SOCAMM2 module standard in development'' with ``a clear upgrade path from today's LPDDR5X SOCAMM2 modules.'' \url{https://www.jedec.org/news/pressreleases/jedec%C2%AE-previews-lpddr6-roadmap-expanding-lpddr-data-centers-and-processing-memory}
\bibitem{sandisk-hbf-facts} SanDisk, \emph{High Bandwidth Flash Fact Sheet} (Jul.\ 2025): 512\,GB per 16-die stack at 256\,Gb per die; 1.6\,TB/s read bandwidth in generation 1; generation 2 ``set to reach up to 1\,TB'' at over 2\,TB/s and 0.8$\times$ the power; generation 3 at 1.5\,TB and 3.2\,TB/s at 0.64$\times$; ``non-volatile, no refresh power needed''; ``one of the lowest cost per bit of any memory technology''; ``similar cost'' to HBM per stack; higher latency than HBM acknowledged and unquantified. \url{https://documents.sandisk.com/content/dam/asset-library/en_us/assets/public/sandisk/collateral/company/Sandisk-HBF-Fact-Sheet.pdf}


% ============ APPLIANCE-VS-RUBIN COST MODEL ANCHORS ============
\bibitem{kimi-k3-arch} Moonshot AI, \emph{Kimi K3} technical report, arXiv:2607.24653 (Jul.\ 2026), and the published model configuration: 93 layers (69 Kimi Delta Attention + 24 MLA), 896 experts with 16 routed and 2 shared, MXFP4 experts; MLA KV cache $\sim$27\,KB per token at BF16 (kv\_lora\_rank 512 + rope 64), KDA state $\sim$207\,MB per sequence. \url{https://arxiv.org/pdf/2607.24653}
\bibitem{vllm-k3} vLLM project blog (27 Jul.\ 2026), Kimi K3 deployment guide: ``2K+ tokens/GPU-s'' at the high-throughput operating point on GB300 NVL72; single-user 111--118 tok/s at TP8--16, 331--370 with speculative decoding; minimum deployment 8$\times$B300 or 16$\times$B200. \url{https://vllm.ai/blog/2026-07-27-k3}
\bibitem{inferencex-kimi} SemiAnalysis InferenceX (19 May 2026): Kimi K2.5 (1T/32B) on B200 NVFP4 at TP4---3,203 tok/s per GPU at 32 tok/s per user, 3,879 at 90; H200 INT4 at 991--1,175. And GB200 NVL72 on DeepSeek-R1 FP4: 4,130 tok/s per GPU at 125 tok/s per user, peak 14,659 at $\sim$18. \url{https://inferencex.semianalysis.com/blog/b200-nvfp4-vs-h200-int4-kimi-k2-vllm-perf-per-dollar}
\bibitem{vr200-price} Morgan Stanley estimate for the Vera Rubin VR200 NVL72, reported May 2026: $\sim$\$7.8M per rack, of which GPUs $\sim$\$4.0M ($\sim$\$55k each), memory $\sim$\$2.0M (20.7\,TB HBM4 plus 54\,TB LPDDR5X), Vera CPUs $\sim$\$180k---labelled a bill of materials but stated to be ``representative of what cloud service providers pay, not what it costs Nvidia itself to build one'' (PC Gamer). Tom's Hardware (24 Mar.\ 2026): ODM quotes for VR200 NVL72 \$5--7M without warranty, ``as much as \$8.8M''; NVIDIA-supplied trays $\sim$90\% of a system's cost, integrator margins thin; GB200 NVL72 \$2.8--3.4M and GB300 NVL72 \$6--6.5M for comparison. Loop Capital (A.~Baruah) put GB300 NVL72 at \$3.7--4.0M and SemiAnalysis GB200 NVL72 at $\sim$\$3.1M bare / \$3.9M all-in---the dispersion is the point. \url{https://www.pcgamer.com/hardware/a-single-nvidia-vera-rubin-rack-is-estimated-to-cost-usd7-803-148-with-over-usd2-million-of-that-figure-spent-on-memory-alone/}; \url{https://www.tomshardware.com/tech-industry/artificial-intelligence/price-of-nvidias-vera-rubin-nvl72-racks-skyrockets-to-as-much-as-usd8-8-million-apiece-but-server-makers-margins-will-be-tight-nvidia-is-moving-closer-to-shipping-entire-full-scale-systems}
\bibitem{vr200-power} VR200 NVL72 rack power estimates: $\sim$190\,kW at 1.8\,kW per GPU (Max-Q) to $\sim$230\,kW at 2.3\,kW (Max-P); 50\,PFLOPS FP4 inference per GPU. \url{https://hashrateindex.com/blog/nvidia-vera-rubin-nvl72-specs-breakdown/}
\bibitem{dgx-spark-price} NVIDIA DGX Spark: GB10 Grace Blackwell, 128\,GB LPDDR5X at 273\,GB/s, ``up to 1 PFLOP FP4'' (sparse); Founders Edition \$3,999 at launch (Oct.\ 2025), raised to \$4,699 on 27 Feb.\ 2026 citing memory shortages; ASUS Ascent GX10 partner unit \$2,999 at launch. \url{https://www.tomshardware.com/desktops/mini-pcs/nvidia-dgx-spark-gets-18-percent-price-increase-as-memory-shortages-bite-founders-edition-now-usd4-699-up-from-usd3-999}
\bibitem{framework-desktop} Framework Desktop with AMD Ryzen AI Max+ 395 and 128\,GB LPDDR5X, \$1,999 at announcement (Feb.\ 2025); GMKtec EVO-X2 with the same part at \$1,799.99 early-bird (May 2025). Post-surge: AMD's own Ryzen AI Halo 128\,GB desktop at \$3,999 (May 2026). \url{https://www.techspot.com/news/106934-framework-first-ever-desktop-pc-packs-amd-ryzen.html}
\bibitem{nand-spot} TrendForce memory spot update (16 Jul.\ 2025): 512Gb TLC NAND wafer spot \$2.673, i.e.\ $\sim$\$0.042 per GB; by 9 Aug.\ 2026 \$21.125, i.e.\ $\sim$\$0.33 per GB. \url{https://www.trendforce.com/news/2025/07/16/insights-memory-spot-price-update-prices-slightly-cool-off-after-recent-surge-with-ddr4-leading-the-drop}
\bibitem{lpddr-presurge} LPDDR5X pricing: $\sim$\$3.3 per GB in Q1 2025 (derived from SemiAnalysis's ``3$\times$ since Q1 2025 to $\sim$\$10/GB'' via Wccftech); TrendForce 2Q26 at \$19.3--19.8 per GB with long-term agreements signed at \$7.8--21; +58--63\% QoQ in 1Q26 and +78--83\% in 2Q26. \url{https://www.trendforce.com/presscenter/news/20260514-13044.html}
\bibitem{cowos-cost} Advanced-packaging cost estimates (2025--26, citing JPMorgan, Morgan Stanley and TrendForce): CoWoS-S $\sim$\$750 per H100-class chip ($\sim$2,500\,mm$^2$); CoWoS-L \$1,000--1,100 per B200-class chip; processed CoWoS wafer \$10--12k; interposer 50--70\% of packaging cost. \url{https://siliconanalysts.com/analysis/cowos-packaging-cost-chiplet-vs-monolithic-2026}
\bibitem{matsuoka-dc-survey} S.~Matsuoka, \emph{Datacenter Construction for State-of-the-Art GPU Machines: A survey with three reference architectures, explicit cost scopes, evidence-graded sources, P10/P50/P90 ranges, escalation to 2030, an independent Japanese reconstruction, and a costed lean reference design}, Rev.\ 5, RIKEN Center for Computational Science technical report, Aug.\ 2026 (all figures from public sources or the author's modelling; distributed without restriction). Facility-only cost per watt of IT load, excluding land and IT, at 50\,MW: architecture A (resilient inference/cloud, Tier-III-capable) \$11.9/W P50 with a \$8.5--17.7 P10--P90 band; B (AI-training-optimised, warm-water, chillerless) \$9.0/W; C (interruptible/phased) \$8.4/W; utility interconnection the widest single driver; large-load interconnection 2--7 years by market; discrete-block penalty for capacity built ahead of load. Author's own companion survey.
\bibitem{epoch-dc-tco} Epoch AI, ``Total cost of ownership of a one-gigawatt AI data center'' (14 May 2026): \$38B capex, of which \$26B compute (servers and network) and \$12B construction (facility, land, utility works); annualized \$8.5B/yr, servers $\sim$60\%; opex \$0.9B/yr of which energy \$0.6B; IT life 5 years, facility 14. \url{https://epoch.ai/data-insights/ai-datacenter-cost-breakdown}
\bibitem{jll-dc-2026} JLL, \emph{2026 Global Data Center Outlook} (Jan.\ 2026): shell-and-core \$11.3M/MW global average (\$10.7M in 2025), AI technology fit-out up to \$25M/MW on top; after ten years $\sim$65\% of facilities have received significant renovation, after fifteen nearly all; tenant IT refreshed roughly every five years. \url{https://www.jll.com/en-us/insights/market-outlook/data-center-outlook}
\bibitem{dc-facility-costs} Turner \& Townsend, \emph{Data Centre Construction Cost Index 2025}: Silicon Valley \$13.3/W, Tokyo \$15.2, New Jersey \$12.9, Phoenix \$9.8 (shell and MEP; excludes utility works, site works and IT); liquid-cooled premium 7--10\% over air. \url{https://reports.turnerandtownsend.com/data-centre-construction-cost-index-2025/data-centre-cost-trends}
\bibitem{dc-reit-lives} Equinix Form 10-K FY2025 (Feb.\ 2026): buildings 12--60 years, core systems 3--40, leasehold improvements 12--40; Digital Realty Form 10-K FY2025 (Feb.\ 2026): buildings and improvements 5--39 years, machinery and equipment 7--15. \url{https://www.sec.gov/Archives/edgar/data/1101239/000110123926000032/R52.htm}
\bibitem{kpmg-dc-opex} KPMG, \emph{Benchmarking CapEx and OpEx in the Global Data Centre Market 2026} (May 2026): 30\,MW greenfield cloud facility, non-electric opex \$7.1--10.6M/yr across four European markets ($\sim$\$237--353 per kW-IT per year: engineering labour, facilities O\&M, security, insurance, property tax, IT and network, cleaning), labour $\sim$40\% of opex; construction capex \$6.7--8.5M/MW. \url{https://assets.kpmg.com/content/dam/kpmgsites/ie/pdf/insights/strategy/ie-global-data-centre-cost-benchmarks.pdf.coredownload.inline.pdf}
\bibitem{dc-equipment-lives} Schneider Electric White Paper 231, ``Battery technology for data centers'': VRLA calendar life 4--6 years vs.\ Li-ion 10--15, VRLA replaced 2--3 times over a fifteen-year facility life; ASHRAE equipment service-life table (Applications Handbook, as reproduced): centrifugal chillers $\sim$23 years, cooling towers 20, transformers 30, pumps 20; three-phase UPS modules commonly 10--12 years. \url{https://download.se.com/files?p_Doc_Ref=SPD_VAVR-A5VTU5_EN}
\bibitem{meta-hyperion-bond} IFR: Beignet Investor LLC (Meta/Blue Owl Hyperion joint venture) \$27.3B senior secured amortising 144A bond, 6.581\% coupon, T+225\,bp, priced 16 Oct.\ 2025, maturity May 2049, rated A+, PIMCO anchor. \url{https://www.ifre.com/ifr-awards/2327933/financing-package-blue-owl-capitalbeignet-investors-us27.3bn-23.6-year-bond}
\bibitem{coreweave-notes} CoreWeave investor releases: \$2.0B 9.25\% senior notes due 2030 (closed 27 May 2025); \$1.75B 9.00\% senior notes due Feb.\ 2031 at par (22 Jul.\ 2025); \$1.0B 9.75\% senior notes due Oct.\ 2031 (16 Apr.\ 2026); \$8.5B delayed-draw term loan at SOFR+225\,bp (31 Mar.\ 2026). \url{https://investors.coreweave.com/news/news-details/2026/CoreWeave-Announces-Pricing-of-1000-million-of-9-750-Senior-Notes-due-2031/default.aspx}
\bibitem{cbre-h2-2025} CBRE, \emph{North America Data Center Trends H2 2025} (26 Feb.\ 2026) and Northern Virginia market release (18 Mar.\ 2026): NoVA 10\,MW+ requirements asking \$155--185 per kW per month, vacancy 0.5\%, 1,102\,MW absorbed in 2025; primary-market 250--500\,kW average \$195.94/kW-mo ($+$6.5\% y/y), wholesale \$196.25 (record); power passed through at cost by market convention. \url{https://www.cbre.com/insights/books/north-america-data-center-trends-h2-2025}
\bibitem{rail-us} Association of American Railroads, \emph{A Short History of American Freight Railroads}: mileage peaked in 1916; $\sim$254,000 miles operated by 1917; more than 70,000 miles ($\sim$30\% of the network) in receivership by 1937; Federal Railroad Administration, Freight Rail Overview (updated Feb.\ 2025): ``almost 140,000 route miles.'' \url{https://www.aar.org/wp-content/uploads/2018/05/AAR-Short-History-American-Freight-Railroads.pdf}
\bibitem{rail-uk} Office of Rail and Road, \emph{Rail infrastructure and assets 2024--25}: route length open for traffic in Great Britain 15,747\,km (9,785 miles) at 31 Mar.\ 2025; network $\sim$18,000 route miles in 1963; the Beeching report proposed 5,000 route miles and 2,363 stations for closure (National Archives). \url{https://dataportal.orr.gov.uk/media/3a1mbkaz/rail-infrastructure-and-assets-2024-25.pdf}
\bibitem{jr-hokkaido} Hokkaido Railway Company (JR Hokkaido), company profile and operating-area pages: 2,240.5\,km of route at 1 Apr.\ 2026 (trunk 1,378.9, regional 861.6); JR Hokkaido, ``Sections that are difficult for the company to maintain on its own'' (18 Nov.\ 2016): 13 sections, 1,237.2\,km, of a $\sim$2,569\,km network; JR Hokkaido press releases of 18 Mar.\ and 15 Apr.\ 2026 listing bus conversions 2019--2026 (Sekisho Shin-Yubari--Yubari, Sassho, Hidaka Mukawa--Samani, Rumoi in two stages, Nemuro Furano--Shintoku; $\sim$312\,km) and the eight remaining ``yellow'' sections (925.7\,km) under discussion; the $\sim$4,000\,km JNR-era peak is a widely repeated figure (Traffic News, Oct.\ 2016) without an official series. \url{https://www.jrhokkaido.co.jp/corporate/company/com_02.html}
\bibitem{rail-1893} S.~Daggett, \emph{Railroad Reorganization} (Harvard, 1908): in 1893 alone more than 27,000 miles with almost \$2B of capitalization were taken over by the courts, and 40,503 miles were sold at foreclosure in 1894--98; W.~N.~Leonard, ``The Decline of Railroad Consolidation,'' \emph{J.\ Econ.\ Hist.} 9 (1949): about one-fourth of all rail mileage in receivership after the Panic of 1893. \url{https://www.gutenberg.org/cache/epub/55397/pg55397-images.html}
\bibitem{dc-cancellations} DCD (27 Mar.\ 2025): TD Cowen reports Microsoft cancelled or deferred up to 2\,GW of data-centre projects in the US and Europe over six months, ``point[ing] to data center oversupply relative to its current demand forecast''; DCD (6 Mar.\ 2026): Oracle and OpenAI drop plans to expand the flagship Abilene Stargate site (600--800\,MW), citing financing and OpenAI's frequently changing demand forecasting, the 4.5\,GW agreement remaining intact. \url{https://www.datacenterdynamics.com/en/news/oracleopenai-drop-plans-to-expand-flagship-abilene-stargate-site-meta-in-talks-to-pick-up-crusoe-capacity-with-nvidias-help/}
% ---- Appendix: High Bandwidth Flash dossier ----
\bibitem{hbf-xlflash} Toshiba Memory (now Kioxia), XL-FLASH storage-class memory introduced at Flash Memory Summit, Aug.\ 2019: SLC, $\sim$5\,\textmu s read latency, 128\,Gb die; EE Times, ``Flash stacks slowly rise'' (7 Aug.\ 2019); Kioxia XL-FLASH CXL module samples Aug.\ 2026. \url{https://www.eetimes.com/flash-stacks-sloooowly-rise/}
\bibitem{hbf-optane} ServeTheHome (Jul.\ 2022): Intel winds down the Optane business with a \$559M inventory impairment in Q2 2022 results. \url{https://www.servethehome.com/intel-optane-559m-impairment-with-q2-2022-wind-down/}
\bibitem{hbf-znand} Tom's Hardware (8 Aug.\ 2025): Samsung revives Z-NAND at FMS 2025 as a GPU-direct storage tier, claiming 15$\times$ conventional NAND performance and 80\% lower power. \url{https://www.tomshardware.com/pc-components/storage/samsungs-revived-z-nand-targets-15x-performance-increase-over-traditional-nand-once-a-competitor-to-intels-optane-z-nand-makes-a-play-for-ai-datacenters}
\bibitem{hbf-optane-analogy} XenoSpectrum (31 Jul.\ 2026), ``Sandisk HBF and the Optane lesson'': two-vendor standard risk; ``excellent technical specifications alone don't guarantee adoption.'' \url{https://xenospectrum.com/en/sandisk-hbf-optane-standardization/}
\bibitem{sandisk-hbf-blog} SanDisk, ``Scaling beyond the wall: inside Sandisk's High Bandwidth Flash for AI'' (6 Aug.\ 2025): CBA architecture, massively parallel read paths; Llama 3.1 405B weight-read simulation ``within 2.2\% of unlimited-capacity HBM''; samples 2H CY2026, first AI-inference devices early 2027; ``it is bandwidth that matters, not latency.'' \url{https://www.sandisk.com/company/newsroom/blogs/2025/scaling-beyond-the-wall-inside-sandisks-high-bandwidth-flash-for-ai}
\bibitem{hbf-semieng} Semiconductor Engineering, ``Flash getting stacked: high-bandwidth version'' (14 May 2026): SanDisk's Cynthia Hsu on choosing OCP over JEDEC (``workstreams are highly goal-oriented \ldots iterate on the specification in real-time''); roadmap Gen.\ 1/2/3 = 1$\times$/1.5$\times$/2$\times$ capacity, 1$\times$/1.45$\times$/2$\times$ bandwidth, 1$\times$/0.8$\times$/0.64$\times$ energy; 17-die stack (16 core + base); Jim Handy: ``Most of the speed you lose is in the write cycle.'' \url{https://semiengineering.com/flash-getting-stacked-high-bandwidth-version/}
\bibitem{hbf-toms-2025} Tom's Hardware, A.~Shilov (13 Feb.\ 2025): SanDisk's HBF at its 2025 Investor Day---512\,GB per 16-high stack of 256\,Gb dies, up to 4\,TB ``VRAM'' per GPU, ``match the bandwidth of HBM \ldots 8 to 16 times capacity at a similar cost point''; ``minor protocol changes'' from HBM; caveats on latency, endurance, block granularity and thermal stress near a GPU. \url{https://www.tomshardware.com/pc-components/dram/sandisks-new-hbf-memory-enables-up-to-4tb-of-vram-on-gpus-matches-hbm-bandwidth-at-higher-capacity}
\bibitem{hbf-bf-2025} Blocks \& Files (12 and 25 Feb.\ 2025): SanDisk investor-day roadmap; Ilkbahar's edge framing of a 512\,GB HBF die subdivided into eight 64\,GB mini-arrays for a mixture-of-experts model on a smartphone. \url{https://blocksandfiles.com/2025/02/12/sandisk-spills-its-technolgy-futures-beans/}
\bibitem{hbf-tab} SanDisk (BusinessWire, 24 Jul.\ 2025): HBF Technical Advisory Board formed---David Patterson (chair; Pardee Professor Emeritus, UC Berkeley; Google Distinguished Engineer) and Raja Koduri (Oxmiq Labs; formerly AMD and Intel); Patterson: HBF ``shows promise \ldots delivering unprecedented memory capacity at high bandwidth.'' \url{https://www.businesswire.com/news/home/20250724968078/en/Sandisk-Forms-HBF-Technical-Advisory-Board-to-Guide-Development-and-Strategy-for-High-Bandwidth-Flash-Memory-Technology}
\bibitem{hbf-mou} SanDisk investor relations (6 Aug.\ 2025): SanDisk and SK hynix sign memorandum of understanding at FMS 2025 ``to establish specifications for High Bandwidth Flash''; HBF named FMS 2025 Best of Show, Most Innovative Technology. \url{https://investor.sandisk.com/news-releases/news-release-details/sandisk-collaborate-sk-hynix-drive-standardization-high}
\bibitem{hbf-kioxia-module} Kioxia (20 Aug.\ 2025): prototype 5\,TB, 64\,GB/s high-bandwidth flash memory module---PCIe 6.0 x8, daisy-chained controllers over 128\,Gbps PAM4, under 40\,W, developed under NEDO's Post-5G programme. \url{https://apac.kioxia.com/en-apac/about/news/2025/20250820-1.html}
\bibitem{hbf-skh-ocp25} SK hynix Newsroom (27 Oct.\ 2025): next-generation NAND strategy at OCP Global Summit 2025---the ``AIN'' family, of which AIN B ``leverages HBF technology by vertically stacking multiple NAND dies''; joint ``HBF Night'' with SanDisk. \url{https://news.skhynix.com/sk-hynix-presents-next-generation-nand-storage-product-strategy-at-ocp-2025/}
\bibitem{hbf-bf-years} Blocks \& Files (27 Nov.\ 2025), ``High bandwidth flash is years away despite its promise'': two or more years from commercial availability; NVIDIA involvement essential; a hypothetical 16-Hi stack of 321-layer dies is 5,136 layers. \url{https://www.blocksandfiles.com/flash/2025/11/27/high-bandwidth-flash-is-years-away-despite-its-promise/1713383}
\bibitem{hbf-ymtc} DigiTimes (31 Dec.\ 2025), ``YMTC goes full speed toward HBF next''; YMTC's Fan Zengxu at the Xinchen developer conference (26 Dec.\ 2025) on HBF fitting Xtacking and edge-AI inference. \url{https://www.digitimes.com/news/a20251231PD217/ymtc-3d-hbm-flash-technology.html}
\bibitem{hbf-ma-patterson} X.~Ma and D.~Patterson (Google DeepMind), ``Challenges and Research Directions for LLM Inference Hardware,'' arXiv:2601.05047 (8 Jan.\ 2026; IEEE Computer, May 2026): proposes High Bandwidth Flash for ``10$\times$ memory capacity with HBM-like bandwidth''; ballpark stack 512\,GB, 1,638\,GB/s read, under 80\,W, microsecond latency, 4\,KB read granule, low write endurance. \url{https://arxiv.org/abs/2601.05047}
\bibitem{hbf-h3} Ha, Kim and Kim (SK hynix), ``H$^3$: Hybrid Architecture Using HBM and HBF for Cost-Efficient LLM Inference,'' IEEE Computer Architecture Letters 25(1), Jan.\ 2026 (DOI 10.1109/LCA.2026.3660969): B200-class package with 8 HBM3E and 8 HBF stacks; HBF $\sim$16$\times$ HBM capacity, similar bandwidth, microsecond latency, $\sim$100k write cycles, up to 4$\times$ HBM power; 40\,MB SRAM latency-hiding buffer needing $>$80\% hit rate; 1.25$\times$ throughput at 1M context, 6.14$\times$ at 10M, 18.8$\times$ batch, 2.69$\times$ performance per watt; ``suitable for gigantic read-only use cases.'' \url{https://ieeexplore.ieee.org/document/11371745}
\bibitem{hbf-consortium} SanDisk (25 Feb.\ 2026): SanDisk and SK hynix begin global standardization of HBF through a dedicated workstream under the Open Compute Project. \url{https://www.sandisk.com/company/newsroom/press-releases/2026/2026-02-25-sandisk-and-sk-hynix-begin-global-standardization-of-next-generation-memory-solution-high-bandwidth-flash-hbf}
\bibitem{hbf-kioxia-tsv} Blocks \& Files (20 Mar.\ 2026): Kioxia Europe CTO Axel Stoermann on Kioxia's own high-bandwidth flash---up to 32 TSV-stacked dies on a shared interposer beside HBM, ``with Nvidia''; XL-Flash the primary candidate cell; ``The GPU is next to the HBM and then next to HBF, all connected to one and the same interposer.'' \url{https://www.blocksandfiles.com/flash/2026/03/20/kioxia-eyes-gpu-memory-stack-with-high-bandwidth-flash-push/5210520}
\bibitem{hbf-q3call} SanDisk Q3 FY2026 earnings call (30 Apr.\ 2026), transcript: ``still on the timeline \ldots of having the NAND late this year, and looking for more of a system with the controller early to mid next year''; HBF ``not a substitute for an enterprise SSD.'' \url{https://www.fool.com/earnings/call-transcripts/2026/04/30/sandisk-sndk-q3-2026-earnings-transcript/}
\bibitem{hbf-equipment} TrendForce (5 and 22 Jun.\ 2026): Hanmi Semiconductor to deliver first HBF thermo-compression bonders in 2H26, with Hanwha Semitech, ASMPT and K\&S developing HBF tools; SanDisk US patent 12,430,274 B2 bonding a processor onto a CBA NAND tile with HBM stacks on a shared interposer. \url{https://www.trendforce.com/news/2026/06/05/news-hbf-opens-new-equipment-race-hanmi-semiconductor-reportedly-eyes-first-tc-bonder-deliveries-in-2h26/}
\bibitem{hbf-micron-hotinfra} Petrucci, Zacarias and Tanna (Micron), ``Is High-Bandwidth Flash All You Need?'' HotInfra Workshop at ISCA 2026 (28 Jun.\ 2026): roofline of tier-2 options---HBM 3.0\,pJ/bit; on-package HBF 4\,TiB, 8\,TB/s nominal, 20\,\textmu s, 8.5\,pJ/bit, 36\% scattered-read efficiency; on-package LPDDR 2\,TiB, 2\,TB/s, 2.4\,pJ/bit; on DeepSeek-V4-Pro, Kimi K2.6 and Nemotron-3-Ultra at 1M context, on-package LPDDR beats HBF at every operating point (up to 2.1$\times$ on KV-heavy workloads). \url{https://hotinfra.org/2026/papers/hotinfra26-final83.pdf}
\bibitem{hbf-flashaccel} Wang, Xue, Sun, Zhang, Dou, Li and Chen (ICT-CAS / UCAS / IME-CAS), ``FlashAccel: Leveraging High-Bandwidth Flash for High-Throughput LLM Inference,'' arXiv:2607.10186 (11 Jul.\ 2026): HBF stack modelled at 768\,GB/s, 192\,GB, SLC, 4\,\textmu s tR, 4\,KB page, 8\,pJ/bit read, 1M P/E; six stacks plus HBM on an H200-class node give 2.54$\times$ throughput per GPU and 1.93$\times$ energy efficiency at a 100\,ms SLO on Qwen3-235B/480B, DeepSeek-V3 and Llama-3.1-405B, using $\sim$19\,MB hyper-page reads and SRAM prefetch. \url{https://arxiv.org/abs/2607.10186}
\bibitem{hbf-bf-spec} Blocks \& Files (5 Aug.\ 2026), ``Getting flash up close and personal to GPUs'': the OCP HBF Architecture Specification v0.7.0 (High-Level Base Die Specification)---grades 0.384 / 1.536 / 3.072\,TB/s; up to 16 independent full-duplex host channels per stack, 1--4 AXI channels each; each host channel connected to an independent set of NAND dies; developed in about six months. \url{https://www.blocksandfiles.com/flash/2026/08/05/getting-flash-up-close-and-personal-to-gpus/5283327}
\bibitem{hbf-toms-spec} Tom's Hardware (4 Aug.\ 2026): SK hynix and SanDisk unveil the HBF specification at FMS 2026---8-Hi and 16-Hi to 512\,GB, $\sim$0.4--3.0\,TB/s, UCIe up to 64\,GT/s over 64 lanes; the 3\,TB/s grade exceeds a single HBM4 stack ($\sim$2\,TB/s); AMD, Broadcom, Intel, NVIDIA, Marvell, Micron, Qualcomm and Samsung ``show no stated interest.'' \url{https://www.tomshardware.com/pc-components/ssds/sandisk-and-sk-hynix-unveil-hbf-spec-up-to-16-hi-nand-stacks-3-tb-s-bandwidth-ucie}
\bibitem{hbf-samsung-fms26} Samsung Newsroom (5 Aug.\ 2026): next-generation 3D memory vision at FMS 2026---zHBM (DRAM stacked atop the accelerator) and zNAND-O (four to eight V-NAND dies on a TSV base die with UCIe to an on-device NPU, targeted 2028); Samsung distinguishes HBF as ``memory semantic tiering next to GPUs in servers.'' Kioxia stays outside the OCP HBF effort (Technologies.org, Aug.\ 2026). \url{https://news.samsung.com/global/samsung-unveils-next-gen-3d-memory-vision-at-fms-2026-charting-the-future-of-ai-infrastructure}
\bibitem{hbf-investorday} SanDisk 2026 Investor Day deck and press release (13 Aug.\ 2026): ``First HBF memory die taped out''; ``first HBF inference product samples \ldots 2027''; Technical Advisory Board slide adds Jim Keller (Tenstorrent); simulation of 4\,TB HBF at 12.8\,TB/s per GPU against 192\,GB HBM---``4 HBF GPUs = 8 HBM GPUs''; ``100B+ models on your edge device.'' Reuters: Ilkbahar on early samples ``next year.'' \url{https://investor.sandisk.com/news-releases/news-release-details/sandisk-details-growth-strategy-and-long-term-financial-model}
\bibitem{hbf-tapeout} TrendForce (14 Aug.\ 2026), ``Sandisk reportedly tapes out first HBF product, targets 2027 samples and 2028 production'': Mizuho and Goldman commentary; Goldman---HBF-only configuration could halve GPU count for the same tokens; Cantor \$2,900 and Goldman \$2,200 price targets; Mizuho---``unlikely to fully replace HBM.'' \url{https://www.trendforce.com/news/2026/08/14/news-sandisk-reportedly-tapes-out-first-hbf-product-targets-2027-samples-and-2028-production/}
\bibitem{hbf-counterpoint} Counterpoint Research (Aug.\ 2026), ``Sandisk Investor Day: caching out the NAND cycle with contracts'': HBF ``a longer-term play,'' ``unproven for AI workloads,'' absent from SanDisk's FY2028--30 financial model---``if HBF lands, it is upside''; Wedbush (via TheStreet) not building HBF into estimates ``until adoption becomes clearer''; YMTC flagged as a competitive threat. \url{https://counterpointresearch.com/en/insights/sandisk-investor-day-caching-out-the-nand-cycle-with-contracts}
\bibitem{hbf-xeno} XenoSpectrum (15 Aug.\ 2026): critique of the investor-day slides---12.8\,TB/s HBF baseline set against HBM3E-era parts while HBM4E is $\sim$4.1\,TB/s per stack; Qwen3-480B sized at 960\,GB in BF16 (480\,GB FP8, 240\,GB FP4); no endurance specification; cites arXiv ``A Full-Stack Characterization of High-Bandwidth Flash for KV-Centric LLM Serving'' (Peking Univ., Aug.\ 2026). Wccftech carried the same critique. \url{https://xenospectrum.com/en/sandisk-investor-day-hbf-criticism-hbm-bandwidth-quantization-endurance/}
\bibitem{hbf-pku} Z.~Li et al.\ (Peking University), ``A Full-Stack Characterization of High-Bandwidth Flash for KV-Centric LLM Serving,'' arXiv, Aug.\ 2026, as reported by XenoSpectrum and NewMaxx (16 Aug.\ 2026): HBF used for KV cache raises latency 2--5.5$\times$ and lowers throughput 1.1--2.7$\times$; the paper itself was not retrievable at this edition. \url{https://borecraft.com/2026/08/16/hbf-sucks-a-full-stack-characterization-of-high-bandwidth-flash-for-kv-centric-llm-serving/}
\bibitem{hbf-heatwatch} Luo, Ghose, Cai, Haratsch and Mutlu, ``HeatWatch: Improving 3D NAND Flash Memory Device Reliability by Exploiting Self-Recovery and Temperature Awareness,'' HPCA 2018: retention loss in 3D charge-trap NAND accelerates superlinearly with temperature; 70$\to$0\,$^\circ$C slows retention loss 58\%; lifetime improvement 3.85$\times$ with temperature-aware policies. \url{https://www.cs.cmu.edu/~yixinluo/index_files/heatwatch_hpca18.pdf}
\bibitem{hbf-straw} Chun et al.\ (Seoul National University), ``STRAW: A Stress-Aware WL-Based Read Reclaim Technique for High-Density NAND Flash-Based SSDs,'' IEEE Computer Architecture Letters 2024 / arXiv:2501.02517: read disturb is strongly word-line-heterogeneous; stress-aware reclaim reduces reclaim writes by 83.6\%. \url{https://arxiv.org/abs/2501.02517}
\bibitem{hbf-businesskorea} BusinessKorea (7 Aug.\ 2026): FMS 2026 panel---Google DeepMind's Xiaoyu Ma on HBF's ``high power efficiency'' with capacity and bandwidth; SK hynix VP Eui-cheol Lim on HBF pairing with HBM in ``reading-centric inference''; SanDisk VP Rajeev Nagabhirava acknowledging NAND degradation under frequent writes, a risk that ``diminishes if HBF operates primarily in read-heavy roles.'' \url{https://www.businesskorea.co.kr/news/articleView.html?idxno=274396}
\bibitem{hbf-hbm4-price} TweakTown (6 Aug.\ 2025), citing Chosun: SK hynix HBM4 12-Hi supply price near US\$500 per stack, 60--70\% above HBM3E's $\sim$\$300; later reports near \$560. \url{https://www.tweaktown.com/news/106826/sk-hynix-drastically-raises-next-gen-hbm4-supply-price-in-war-of-nerves-with-nvidia/index.html}
\bibitem{hbf-koreaherald} Korea Herald (4 Feb.\ 2026): SK hynix framing HBF as an HBM complement, first generation 16 $\times$ 32\,GB = 512\,GB, mass-production target 2027; Samsung positioning HBF inside a unified memory--storage tier; unattributed market projection of \$1B (2027) to \$12B (2030). \url{https://www.koreaherald.com/article/10669660}
\bibitem{hbf-kaist} Blocks \& Files (19 Jan.\ 2026), ``A window into HBF progress'': KAIST Prof.\ Kim Joung-ho---512\,GB / 1.638\,TB/s per HBF; three-tier HBM--HBF--networked SSD; claims Samsung and SanDisk aim at NVIDIA/AMD/Google integration late 2027 or early 2028; ``HBF market could be bigger than HBM'' by 2038. \url{https://www.blocksandfiles.com/ai-ml/2026/01/19/a-window-into-hbf-progress-engineering-prof-talks-of-decade-ahead/4090360}
\bibitem{hbf-bf-window} Blocks \& Files (8 Jan.\ 2026), ``Sandisk and high bandwidth flash efforts'': NVIDIA ``has not yet expressed any public interest''; ``Get Nvidia on board \ldots and it's a done deal.'' \url{https://www.blocksandfiles.com/ai-ml/2026/01/08/sandisk-and-high-bandwidth-flash-efforts/4090532}
\bibitem{hbf-register} The Register, T.~Mann (30 Jul.\ 2026): HBF trades write endurance and microsecond latency for capacity; 256\,Gb dies, 512\,GB per 16-Hi, ``14$\times$ more than HBM4''; a proposed 3.12\,TB package from two HBM and six HBF stacks; wear could make ``that \$50,000-plus GPU of yours a paperweight''; the Optane precedent for proprietary memory. \url{https://www.theregister.com/storage/2026/07/30/gpus-could-explode-to-multiple-tb-with-new-storage-inspired-memory-tech/5281363}
\bibitem{hbf-spectrum} IEEE Spectrum (14 Jul.\ 2026), ``High Bandwidth Flash'': NAND die 4.8\,GB/s, DDR5 DIMM 70.4\,GB/s, HBM4E 3.6\,TB/s per stack; SK hynix SVP Hoshik Kim on HBF reducing the number of accelerators required; Jim Handy: ``Flash is atrociously slow for writes, but for reads, it can be coaxed to go pretty fast.'' \url{https://spectrum.ieee.org/high-bandwidth-flash}
\bibitem{hbf-techradar} TechRadar (6 Aug.\ 2026): HBF ``exists only on paper''; NVIDIA absent; commercialization late 2027--2028, meaningful datacentre demand around 2030; SK hynix gave no HBF production date. \url{https://www.techradar.com/pro/sk-hynix-and-sandisk-reveal-first-specs-for-high-bandwidth-flash-but-dont-get-too-excited-you-probably-wont-be-able-to-buy-it-anytime-soon}
\bibitem{hbf-tradethepool} TradeThePool (5 Aug.\ 2026), SanDisk Q4 earnings preview: HBF performance evidence ``is simulation rather than silicon''; long-term agreements ``untested through a real downturn.'' \url{https://tradethepool.com/pre-market-prep/sandisk-stock-sndk-q4-earnings-preview-the-hbf-bet/}
\bibitem{hbf-hbm-roadmap} Tom's Hardware, HBM roadmaps for Micron, Samsung and SK hynix to HBM4 and beyond (2026): HBM4 36--48\,GB per stack, JEDEC ceiling 64\,GB; HBM4E samples May 2026, volume late 2027; per-stack bandwidth 3.6--4.1\,TB/s. \url{https://www.tomshardware.com/tech-industry/semiconductors/hbm-roadmaps-for-micron-samsung-and-sk-hynix-to-hbm4-and-beyond}
\bibitem{hbf-inflash} In-flash and near-flash compute for LLM inference: Lee et al.\ (SNU), ``AiF: Accelerating On-Device LLM Inference Using In-Flash Processing,'' ISCA 2025 (14.6$\times$ over SSD offload); Deng et al., ``KVNAND,'' arXiv:2512.03608 (Dec.\ 2025); Hao et al., ``NVLLM: 3D NAND-centric edge LLM,'' arXiv:2604.25699 (2026, $\sim$65 tok/s on a 30B model, 5.6$\times$ energy reduction); Jang et al.\ (Korea Univ.), ``Dissecting and Re-architecting 3D NAND Flash PIM Arrays,'' arXiv:2511.12860 (Nov.\ 2025). \url{https://arxiv.org/abs/2604.25699}
\bibitem{hbf-huxiu} Huxiu (4 Aug.\ 2026), Beijing consultancy analysis of HBF as a ``lane-change'' opportunity: China lacks high-end HBM production capability but YMTC's Xtacking and hybrid-bonding strengths suit HBF; 294-layer part in production above 90\% yield, $\sim$13\% unit share in Q1 2026, 200k to 300k wafers per month by 2027; hybrid-bonding licence to Samsung; constraint is tooling above $\sim$300 layers under the entity list; NVIDIA the decider. \url{https://www.huxiu.com/article/4880566.html}
\bibitem{hbf-crazyrich} CrazyRich research memo (19 Feb.\ 2026), ``Beyond HBM: the rise of HBF'': secondary cost estimates (HBF \$3--5/GB against HBM \$20--30/GB); reports YMTC ``not pursuing HBF directly'' and partnering CXMT on HBM. \url{https://crazyrich.ai/reports/memos/beyond-hbm-rise-of-hbf-ai-memory}
\bibitem{hbf-postech} Son, Park, Cho, Ham, Mutlu, Lee, Kim and Park (POSTECH and ETH Zurich), ``Exploring High-Bandwidth Flash for Modern LLM Inference: Opportunities and Challenges,'' IEEE Computer Architecture Letters 25(2), Jan.\ 2026 (DOI 10.1109/LCA.2026.3705817): HBF as main GPU memory ``can significantly improve batch size, throughput and flexibility'' and cut GPU count, ``but only if HBM-equivalent read performance and substantial endurance improvement are achieved.'' \url{https://doi.org/10.1109/LCA.2026.3705817}
\bibitem{hbf-haven} Hsu, Xu, Liu, Rosing and Yu (UCSD and Georgia Tech), ``HAVEN: HBF-Augmented Vector Engine for ANN search,'' arXiv:2603.01175 (1 Mar.\ 2026): HBF holding a full-precision vector database on the GPU for RAG; up to 20$\times$ throughput and 40$\times$ latency improvement over GPU-DRAM and GPU-SSD baselines; 4\,KB page baseline. \url{https://arxiv.org/abs/2603.01175}
\bibitem{hbf-isscc} NAND die-level state of the art: Kioxia and SanDisk, ISSCC 2025 (1\,Tb TLC CBA at 4.8\,Gb/s I/O, 29\% read-energy improvement) and ISSCC 2026 (BiCS10 2\,Tb QLC, 332 active layers, $>$85\,MB/s program); Samsung V10 at 5.6\,Gb/s per pin (ISSCC 2025); SK hynix 321-layer 2\,Tb QLC at 75\,MB/s program (ISSCC 2025)---i.e.\ per-die programme rates roughly fifty times below read rates. \url{https://www.sandisk.com/company/newsroom/press-releases/2025/kioxia-and-sandisk-unveil-next-generation-3d-flash-memory-technology}
\bibitem{hbf-offload} The flash/SSD-offloading lineage: Sheng et al., FlexGen, ICML 2023 (arXiv:2303.06865); Alizadeh et al.\ (Apple), ``LLM in a flash,'' ACL 2024 (arXiv:2312.11514); Eliseev and Mazur, arXiv:2312.17238; Xue et al., MoE-Infinity, arXiv:2401.14361; Cao et al., MoE-Lightning, ASPLOS 2025 (arXiv:2411.11217); He et al., ExpertFlow, arXiv:2410.17954; Tang et al., HOBBIT, arXiv:2411.01433; Yu et al., PreScope, arXiv:2509.23638; Kim et al., FlashMoE, arXiv:2601.17063; Xue et al.\ (SJTU), PowerInfer-2, arXiv:2406.06282 (47B MoE from smartphone flash at 11.68 tok/s); Microsoft, ZeRO-Inference (Sep.\ 2022, OPT-30B at 30 tok/s with NVMe offload). \url{https://arxiv.org/abs/2312.11514}
\bibitem{hbf-routing} Mixture-of-experts routing locality: Liang et al.\ (Fudan, Huawei), ``Not All Models Suit Expert Offloading: On Local Routing Consistency of Mixture-of-Expert Models,'' arXiv:2505.16056 (2025--26; 20 MoE models; cache of $\approx$2$\times$ active experts near-optimal; shared-expert designs lower locality); Oncescu, Dao et al., ``Opportunistic Expert Activation,'' arXiv:2511.02237 (Nov.\ 2025; batch-aware rerouting cuts MoE-layer decode latency 15--39\%); Huang et al.\ (Meta), ``Towards MoE Deployment,'' arXiv:2303.06182; Yuan and Long, ``Economics of AI Decoding Chips,'' arXiv:2607.13068 (Jul.\ 2026; the ``MoE tax''). \url{https://arxiv.org/abs/2505.16056}
\bibitem{dots3-note} Xiaohongshu dots studio, \emph{dots3-note Preview} (released 14 Aug.\ 2026): 280B total / 16B active parameters (45 MoE layers, 256 routed + 1 shared expert, top-8; 7B/1.2B-active MoE vision encoder; 800M audio encoder); text, image, video and audio input; 512K context; Apache~2.0; weights on Hugging Face (\texttt{dots-studio/dots3-note-prev}, FP8 variant) and ModelScope; self-reported SWE-bench Verified 78.4, SWE-bench Pro 61, SWE-bench Multilingual 75.7, MMMU-Pro 79.1; TEMPO (``Test-time-scaled Value Estimation with Macro-step Policy Optimization'') described in the release material; technical report ``coming soon.'' \url{https://huggingface.co/dots-studio/dots3-note-prev}
\bibitem{dots3-note-press} 36Kr (14 Aug.\ 2026) and NetEase/ifeng coverage of the dots3-note release: TEMPO divides trajectories of ten to twenty hours into macro-steps; at each boundary the same model switches from actor to critic and estimates the expected return of the current state through reasoning, tool use and test-time scaling, producing dense intermediate training signal; ARC-AGI-3 $\approx$0.35 under a \$500 budget; VibeSearchBench and VibeLifeBench introduced. \url{https://eu.36kr.com/en/p/3938759517896072}
\bibitem{dots-imo} South China Morning Post (22 Jul.\ 2026): Xiaohongshu's dots-note-3.0 becomes the first AI model to earn a perfect 42/42 at the International Mathematical Olympiad. \url{https://www.scmp.com/tech/article/3361482/worlds-first-ai-model-earn-perfect-score-maths-olympiad-comes-chinas-rednote}
\bibitem{dots3-ascend} IT Home (16 Aug.\ 2026): Huawei Ascend completes ``0-day'' inference deployment and optimization of dots3-note on Atlas 800 A3 and Atlas 900 A3 SuperPoD via vLLM-Ascend; vLLM-Ascend project documentation carries a Dots3-Note deployment tutorial validated on Atlas 800I A3. \url{https://docs.vllm.ai/projects/ascend/en/latest/tutorials/models/Dots3-Note.html}
\bibitem{rednote-ipo} Wall Street Journal via Investing.com (17 Jun.\ 2026) and Bloomberg (15 Jun.\ 2026): Xiaohongshu preparing a Hong Kong IPO at a valuation above \$70B, recent private transactions implying more than \$50B; more than 400 million monthly active users; 2025 net income above \$2B. \url{https://www.investing.com/news/stock-market-news/chinas-xiaohongshu-prepares-hong-kong-ipo-at-over-70-bln-valuation--wsj-4746470}
\bibitem{mimo-v25} Xiaomi MiMo, ``MiMo-V2.5 open-sourced'' (27 Apr.\ 2026): MiMo-V2.5 (310B/15B active MoE, native omni-modal text/image/video/audio input, 1M context) and MiMo-V2.5-Pro (1.02T/42B active, text, 1M context, 27T training tokens), both MIT; MiMo-V2-Flash (309B/15B, MIT) Dec.\ 2025. \url{https://huggingface.co/XiaomiMiMo/MiMo-V2.5}
\bibitem{vercel-gateway} Vercel, ``Open-weight models surge to 29\% of volume, price per token flattens,'' AI Gateway production index (13 Jul.\ 2026, June 2026 data): ``Open-weight models ran 29\% of gateway tokens in June on just under 4\% of spend''; $\sim$11\% of tokens in April; DeepSeek 22.6\% of June tokens, within two points of Google; GLM 5.2 took 76\% of its family's tokens within two weeks of launch; Anthropic 61\% of spend on 32\% of tokens; roughly one in eight enterprise customers runs an open-weight model. \url{https://vercel.com/blog/ai-gateway-production-index-july-2026}
\bibitem{nvme-presurge} 1\,TB NVMe SSD retail, pre-surge: $\sim$\$65--85 in Q1 2025 rising to \$75--100 by Q1 2026; a 1\,TB drive $\sim$\$45 in late 2025 had reached $\sim$\$90 by mid-2026; a 2\,TB drive $\sim$\$120--150 a year earlier (mid-2025) reached \$300--480 by Jul.\ 2026. Storage Disk Prices, SSD Price History; Tech Insider, ``SSD Prices Double as NAND Shortage Hits Gaming PCs.'' \url{https://storagediskprices.com/ssd-price-history/}, \url{https://tech-insider.org/ssd-prices-nand-shortage-2026/}
\bibitem{eia-tariff} US Energy Information Administration, Electric Power Monthly (May 2026): average industrial price 8.71\,\textcent/kWh, commercial 13.54\,\textcent/kWh. \url{https://www.eia.gov/electricity/monthly/update/end-use.php}
\bibitem{api-prices-aug26} API list prices, August 2026, per million tokens input/output: Claude Opus 5 \$5/\$25 (Sonnet 5 \$2/\$10); GPT-5.6 Sol \$5/\$30, Luna \$0.20/\$1.20 after the 30 Jul.\ cut; Gemini 3.1 Pro \$2/\$12 to 200K context; Grok 4.6 \$2/\$6; Kimi K3 \$3/\$15 (\$0.30 cached); Qwen3.8-Max \$2/\$6; DeepSeek V4-Pro from 16 Aug.\ \$1.32/\$3.96 peak and \$0.66/\$1.98 off-peak, V4-Flash \$0.44/\$1.32 peak. Vendor pricing pages. \url{https://platform.kimi.ai/docs/pricing/chat-k3}

% ---------------- Hot Chips 2026, the Western mirror, GLM-5.3 ----------------
\bibitem{hc2026} Hot Chips 2026, Stanford University, 23--25 August 2026; 31 technical presentations across memory, server CPU, GPU/accelerator, network and tutorial tracks. Deck numbers cited in this book follow the conference numbering. Where a figure appears only in a deck it is marked as vendor-presented and unaudited; where a vendor withheld a figure that exists elsewhere, the outside source is named. \url{https://hotchips.org/}
\bibitem{hc2026-nv-riscv} Hot Chips 2026 deck 09, Frans Sijstermans (NVIDIA), ``RISC-V CPUs in NVIDIA Servers: CUDA's Third Host ISA'': NVIDIA states the requirements for a RISC-V CPU to host CUDA alongside x86 and Arm---RVA23 compliance plus the RISC-V Server SoC and Server Platform specifications, Boot and Runtime Services (ratified August 2025) and ACPI; CUDA-specific asks reduce to coherent PCIe I/O and full-speed PCIe peer-to-peer with standard root-complex ordering. A second tier for NVLink Fusion requires NVIDIA's NVLink-C2C soft IP behind C2C hubs speaking the CHI coherence protocol, plus DOCA and NCCL support. SiFive appears on the NVLink Fusion partner slide alongside Arm, Fujitsu and Intel. The requirement document itself is stated to be available to partners only. Reported by Chester Lam, Chips and Cheese (24 Aug.\ 2026). \url{https://chipsandcheese.com/p/hot-chips-2026-cuda-targets-risc}
\bibitem{hc2026-riscv-profiles} Hot Chips 2026 deck 07, Krste Asanovi\'c (RISC-V International, SiFive, UC Berkeley), ``Update on RISC-V Standards and Adoption: Profiles and Platforms'': RVA23 ratified October 2024 with vector and hypervisor mandatory; RVA23.1 near ratification; Server Platform specification ratified May 2026. \emph{No ratified matrix extension exists}: one fast-track dot-product proposal (Zvbdota) and three competing task groups---IME (matrix-in-vector), VME (vector-matrix, adding an accumulator file) and AME (separate matrix state)---with ratification targets spanning February 2026 to April 2027. New ``optimization guidance'' options concede that a functional-only ISA string does not tell a vendor how to ship a competitive core. \url{https://riscv-europe.org/summit/2026/media/proceedings/2026-06-10-RISC-V-Summit-Europe-09h00-ASANOVI\%C4\%86-slides.pdf}
\bibitem{hc2026-canonical} Hot Chips 2026 deck 08, Gordan Marku\v{s} (Canonical), ``Enterprise Open Source on RISC-V and Convergence on RVA23'': Ubuntu's riscv64 baseline was RVA20 from 20.04 LTS through 25.04 and is RVA23 from 25.10; 26.04 LTS carries RVA23 as the long-term baseline, with RVA20 silicon retaining a supported path on 24.04 LTS for up to ten years under Ubuntu Pro. The readiness number: roughly 95\% of archive packages build on riscv64 but only \textbf{55.1\%} of tested packages pass, against 92.3\% on amd64 and 92.5\% on arm64, on roughly one seventh the test coverage---attributed to emulation pending server-grade RVA23 hardware. RISE board farm runs on Alibaba and SpacemiT hardware. \url{https://canonical.com/blog/canonical-and-ubuntu-risc-v-a-2025-retro-and-looking-forward-to-2026}
\bibitem{sifive-bigsky} SiFive, ``BigSky SF-2U870'' (24 Aug.\ 2026): the first rackable enterprise RISC-V server---32 P870-D cores at 2.0\,GHz, 256\,GB DDR5-5600, 64 lanes of PCIe Gen5, 2U, fully RVA23 compliant, shipping with Ubuntu 26.04 LTS and RHEL 10, and able to act as a CUDA head node for NVIDIA GPU clusters. Available in limited quantities; no price, no SPEC result and no named customer tape-out disclosed. Futurum's assessment is that BigSky is ``primarily a software porting vehicle rather than revenue-generating hardware,'' valuable because it moves the multi-year software phase of every SiFive-IP-based custom SoC onto hardware a customer can rack today. \url{https://www.hpcwire.com/off-the-wire/sifive-brings-risc-v-to-rackable-servers-with-bigsky-platform/}
\bibitem{riscv-consolidation} Qualcomm's acquisition of Ventana Micro Systems (announced Dec.\ 2025; terms undisclosed; Veyron V1 never reached market and V2 slipped from 2025 to 2026), and Reuters reporting (16 Jun.\ 2026) that Qualcomm was in talks to acquire Tenstorrent for up to \$10B---unconfirmed as of the currency date. Of three Western RISC-V server contenders, one shipped a server, one was absorbed without shipping its first product, and one is a rumoured acquisition target. \url{https://www.theregister.com/2025/12/10/qualcomm_riscv_arm_ventana/}
\bibitem{hc2026-memory-wall} Hot Chips 2026 deck 01, Jim Handy (Objective Analysis), ``Memory: Feeding AI's Voracious Hunger for Data'': lowest spot DRAM price up $7.1\times$ and NAND up $6.9\times$ from the mid-2025 trough; combined DRAM and NAND revenue of the five major suppliers from a $\sim$\$15B quarterly trough in 2Q23 to $\sim$\$180B in 2Q26, almost entirely price-driven rather than bit-driven; hyperscale property and equipment spending from $\sim$\$30B per quarter in 1Q21 to $\sim$\$170B in 2Q26; DRAM wafer starts flat for over a decade and new capacity taking more than two years to install. The deck's summary slide is a photograph of banknotes: the 2026 memory wall is priced, not clocked. Conflict disclosed on slide 7: the data underpin a paid report.
\bibitem{hc2026-hbm4-spread} Hot Chips 2026 decks 02 (Micron), 03 (Samsung) and 04 (SK hynix): the same JEDEC HBM4 generation is quoted by SK hynix at 8\,Gb/s and 2.0\,TB/s per stack, by Micron at 11\,Gb/s, and by Samsung at 3.0\,TB/s---a $1.5\times$ spread in per-stack bandwidth within one nominal generation, so a bill of materials specifying ``HBM4'' does not specify a part. Micron: an eight-stack HBM4 package carries 12{,}344\,mm$^2$ of memory silicon, more than $8\times$ the die area of a typical GPU. Samsung deck 28: memory rose from 52\% to 63\% of AI-chip component spend between 1Q24 and 4Q25.
\bibitem{hc2026-chbm} Hot Chips 2026 deck 03, Sangwook Han (Samsung), ``HBM Base Die: How HBM Will Evolve Using Advanced Logic Processes'': from HBM4 the base die is fabricated on a 4\,nm foundry logic process, which turns the memory vendor into a chip designer. The three-phase roadmap is reclaim xPU area by absorbing the memory controller into the base die; fill the free base-die area with repair, telemetry, test and processing elements; then delete the interposer entirely by bonding the cell-array stack directly onto the xPU (``zHBM''). Custom HBM is explicitly \emph{not} JEDEC-standardised: a customer-specific logic die sharing standard cell-array stacks. Strategically this inverts the value chain---the memory vendor comes to own an interface that belonged to the accelerator designer.
\bibitem{hc2026-scaleup} Hot Chips 2026 decks 23 (AMD, Steve Scott, ``System Architecture of the MI400 Series''), 36 (Microsoft, Maia 200), 37 (SambaNova SN50) and 33 (Meta MTIA 400): AMD builds a 72-GPU load/store domain out of UALink-over-Ethernet on merchant 512-port 200G switch ASICs, deleting the scale-up NIC from the tray; Microsoft runs one Ethernet protocol from chip to a 6{,}144-accelerator two-tier Clos, with slide 16 titled ``Why Unified?'' and listing a standard switch ecosystem as the reason; SambaNova uses ten integrated 800G Ethernet ports per socket and two merchant switches for a 64-socket domain; Meta uses RoCE on dedicated I/O chiplets. Google's optically switched torus is proprietary but is not sold; NVIDIA's NVLink~6 (130 TFLOPS of in-network compute) is the only proprietary scale-up fabric a third party can buy. Of eight parts with an off-package scale-up fabric, six build it from Ethernet-class silicon.
\bibitem{hc2026-uec} Hot Chips 2026 deck 30, Hemal Shah (Broadcom), ``Thor Ultra'': every mechanism the Ultra Ethernet Consortium argued for shipped in production silicon this year---packet spraying, out-of-order placement with in-order delivery, selective retransmission, receiver-driven credit congestion control, packet trimming, programmable transport---but as \emph{RoCEv2 extensions}. Broadcom brands its transport eRoCE and describes it as ``MRC++,'' preserving the verbs ecosystem so xCCL, NCCL, RCCL and MPI run unmodified: ``the ecosystem cost is zero.'' Neither Broadcom's nor NVIDIA's deck presents a UET transport or a libfabric stack, and NVIDIA's networking deck does not mention the UEC at all. The consortium standardised the mechanism set and lost the wire format.
\bibitem{hc2026-jalapeno} Hot Chips 2026 deck 39, Richard Ho, Ravi Narayanaswami and Chris Leary (OpenAI), ``You Can Just Build Things \ldots\ Chips'': Jalape\~no, OpenAI's inference accelerator built with Broadcom---13.4 PFLOP/s mxfp4, 216\,GiB HBM4 at 15.4\,TB/s, 700\,W, a 128-chip scale-up domain at 600\,GB/s and a 2{,}048-chip pod over a half-flattened two-level Broadcom Tomahawk~6 Clos. The organising metric is output tokens per second per all-in utility megawatt at matched interactivity; per-chip throughput and chip count are declared explicit non-goals. Architectural argument against a specialised fleet: ``dark silicon is cheaper than idle accelerators.'' More than half the core was written in the XLS hardware language with an internal model in the loop. Vendor-supplied comparisons against GB200/GB300 and MI355X are unaudited and exclude speculative decoding. \url{https://openai.com/index/openai-broadcom-jalapeno-inference-chip/}
\bibitem{openai-broadcom} OpenAI and Broadcom, ``Strategic collaboration to deploy 10 gigawatts of OpenAI-designed accelerators'' (13 Oct.\ 2025), signed as a term sheet with no financial terms disclosed; deployment from 2H~2026 through 2029. Jalape\~no unveiled 24 June 2026 as the first product. Greg Brockman: ``By building our own chip, we can embed what we've learned from creating frontier models directly into the hardware.'' Foundry, process node, unit cost and unit volume are not disclosed in any primary source. \url{https://investors.broadcom.com/news-releases/news-release-details/openai-and-broadcom-announce-strategic-collaboration-deploy-10}
\bibitem{hc2026-tpu8} Hot Chips 2026 deck 38, Norman P.\ Jouppi and Sridhar Lakshmanamurthy (Google), ``Eighth Generation TPU Family: Two Chips for the Agentic Era'': for the first time Google shipped a training part (8t: 9{,}600 chips, optically switched 3D torus, 6 HBM3E stacks) and an inference part (8i: 1{,}152 chips, seven-hop Boardfly hierarchy, 8 HBM3E stacks, 384\,MB SRAM, collectives moved into the ICI I/O die) \emph{concurrently} rather than alternately, on the stated ground that pre-training, distillation, sampling, MoE and agentic serving now need different interconnect topologies. One XLA/JAX/PyTorch stack targets both. Claimed $>2\times$ perf/TCO and $2\times$ perf/W generation over generation, without a stated workload. No per-chip FLOPS, power, node or die size published. \url{https://www.servethehome.com/google-tpu-8i-for-inference-and-tpu-8t-for-training-announced/}
\bibitem{tpu-external} Sundar Pichai, Alphabet Q1 2026 earnings call (29 Apr.\ 2026): ``As TPU demand grows from AI labs, capital markets firms, and high-performance computing applications, we'll begin to deliver TPUs to a select group of customers in their own data centers in a hardware configuration to expand our addressable market opportunity''---a small share of TPU hardware revenue in 2026, the majority in 2027. Alphabet raised 2026 capital expenditure guidance to \$180--190B in the same call. Google and Broadcom expanded their Anthropic partnership on 6 April 2026 to multiple gigawatts of next-generation TPU capacity from 2027. \url{https://www.datacenterdynamics.com/en/news/google-to-sell-tpus-to-a-select-group-of-customers-for-their-data-centers/}
\bibitem{maia200} Microsoft, ``Maia 200: the AI accelerator built for inference'' (26 Jan.\ 2026): TSMC 3\,nm, 140B+ transistors, 216\,GB HBM3E at 7\,TB/s, 272\,MB on-die SRAM, $>$10 PFLOPS FP4, 750\,W, clusters to 6{,}144 accelerators over one Ethernet protocol; deployed in US Central. The headline claim is economic, not architectural: ``30\% better performance per dollar than the latest generation hardware in our fleet today.'' NVIDIA is not mentioned. \url{https://blogs.microsoft.com/blog/2026/01/26/maia-200-the-ai-accelerator-built-for-inference/}
\bibitem{mtia400} Meta, MTIA roadmap (announced Mar.\ 2026, Broadcom co-designed; Hot Chips 2026 deck 33): four generations on roughly a six-month cadence---MTIA 300 (800\,W, in production), MTIA 400 (1{,}200\,W, 288\,GB HBM3E at 9.2\,TB/s, 6 PF FP8/MX8, 12 PF MX4, in lab test), MTIA 450 (early 2027) and MTIA 500 (later 2027). Meta has deployed ``hundreds of thousands'' of earlier MTIA parts for inference. Meta diverges from the OCP microscaling standard deliberately: 16 elements per shared exponent rather than 32, plus an MX8S $16\times16$ FP8 block format created for training. \url{https://www.tomshardware.com/tech-industry/semiconductors/meta-reveals-four-new-mtia-chips-built-for-ai-inference}
\bibitem{trainium3} Amazon Web Services, Trainium3 general availability (2 Dec.\ 2025): first 3\,nm Trainium, 2.517 PFLOPS FP8 and 144\,GB HBM3E at 4.9\,TB/s per chip, UltraServers to 144 chips, scaling to a million-chip cluster; AWS states more than one million Trainium processors deployed. Matt Garman: ``Trainium already represents a multi-billion-dollar business.'' \url{https://www.hpcwire.com/2025/12/02/aws-brings-the-trainium3-chip-to-market-with-new-ec2-ultraservers/}
\bibitem{aws-nvda} AWS and NVIDIA (26 Aug.\ 2026): two million additional NVIDIA GPUs (Blackwell Ultra, Rubin, Rubin Ultra) across AWS in 2027--2028, plus 100{,}000 GPUs for US government AI factories---and, materially, \emph{Trainium4 will adopt NVIDIA NVLink Fusion and NVIDIA custom HBM}, placing AWS's own accelerator inside NVIDIA's rack-scale interconnect and memory stack. Garman: ``Customers want the freedom to choose the best tools.'' This is the strongest single piece of evidence that NVIDIA is absorbing the custom-silicon wave rather than being displaced by it. \url{https://press.aboutamazon.com/aws/2026/8/aws-and-nvidia-to-deliver-2-million-additional-gpus-and-next-generation-infrastructure-for-agentic-and-physical-ai}
\bibitem{broadcom-q2fy26} Broadcom, second quarter fiscal 2026 results (3 Jun.\ 2026): total revenue \$22.187B, up 48\%; AI semiconductor revenue \$10.8B, up 143\% year on year; third-quarter guidance of more than \$16.0B in AI semiconductors, up about 200\%. Hock Tan attributes growth to ``increasing demand for custom AI accelerators and AI networking.'' Annualising the guided run rate ($>$\$64B a year) against NVIDIA's data-centre run rate ($\sim$\$356B a year on the July-quarter actual) places merchant custom-ASIC revenue at roughly 11--15\% of the combined total, the range set by whether actuals or guidance are annualised on each side [M]. Note that NVIDIA's \$108.0B third-quarter figure is \emph{total-company} guidance; NVIDIA publishes no segment guidance. \url{https://investors.broadcom.com/news-releases/news-release-details/broadcom-inc-announces-second-quarter-fiscal-year-2026-financial}
\bibitem{nvda-500b} NVIDIA, ``NVIDIA partners with Apollo, BlackRock, Blackstone, Brookfield, Goldman Sachs and KKR to establish AI compute infrastructure financing platforms to mobilize over \$500 billion of third-party capital'' (10 Aug.\ 2026). The agreements are memoranda of understanding, ``subject to execution of the final agreements.'' Jensen Huang: ``We began by building chips; today, we are helping create a new class of productive, investable infrastructure: AI factories'' and ``NVIDIA compute is \ldots\ fungible and transferable across customers and operators.'' David Solomon: ``We're excited for the new opportunity to create a market for credit backed by NVIDIA compute.'' The widely reported 25\% residual-value share appears in Huang's own remarks and press reporting, \emph{not} in the release or in any filing. \url{https://nvidianews.nvidia.com/news/nvidia-partners-with-apollo-blackrock-blackstone-brookfield-goldman-sachs-and-kkr-to-establish-ai-compute-infrastructure-financing-platforms-to-mobilize-over-500-billion-of-third-party-capital}
\bibitem{sec-dcs} US Securities and Exchange Commission, Office of Structured Finance interpretive position (29 Jul.\ 2026): securities in data-centre securitizations are not ``asset-backed securities'' under Exchange Act \S3(a)(79), because data centres are ``tangible, physical assets that endure beyond the tenor of the securities'' rather than self-liquidating financial assets---placing them outside Regulation AB and outside the 5\% credit-risk-retention requirement of Regulation RR. The ruling preceded NVIDIA's financing-platform announcement by twelve days. Alston \& Bird (3 Aug.\ 2026); Seward \& Kissel (10 Aug.\ 2026). \url{https://www.alston.com/en/insights/publications/2026/08/data-center-securitizations-sec-regulations}
\bibitem{nvda-ohio} NVIDIA Corporation, Form 8-K filed 17 August 2026 (Items 1.01, 2.03, 7.01): multiple residual value guaranties with SB~Energy Corp.\ on data-centre leases at the Portsmouth Site, Pike County, Ohio, whose tenant is an OpenAI Group PBC affiliate. ``NVIDIA's aggregate payment obligation is cumulatively capped at \$105 billion for its initial commitment under the Agreements.'' Approximately 4.25\,GW of IT load initially with an option on a further 3.8\,GW; twenty-year maximum term; \emph{terminates early if OpenAI achieves a satisfactory credit rating}; OpenAI must reimburse and indemnify NVIDIA for all amounts paid. First phases effective in fiscal 2029. NVIDIA also invests \$1.5B in SB~Energy and becomes exclusive AI compute infrastructure provider for the site. Earlier press reports of a \$250B commitment are superseded by the filed figure. \url{https://www.sec.gov/Archives/edgar/data/1045810/000104581026000069/nvda-20260817.htm}
\bibitem{nvda-q2fy27} NVIDIA, second quarter fiscal 2027 results (quarter ended 26 Jul.\ 2026; announced 26 Aug.\ 2026): revenue \$96.2B, up 106\% year on year; data-centre revenue \$89.0B, up 117\%; GAAP and non-GAAP gross margin 75.0\%; GAAP net income \$59.688B; third-quarter guidance \$108.0B $\pm$2\% with gross margin 74.0\%, excluding China data-centre compute revenue. GAAP diluted EPS (\$2.46) \emph{exceeds} non-GAAP (\$2.22) because GAAP other income of \$7.773B is almost entirely gains on equity securities---that is, roughly an eighth of GAAP net income for the quarter is mark-ups on stakes in companies that buy NVIDIA's chips. Shares closed up 8.7\% on 27 August, the largest single-day gain since April 2025. \url{https://www.sec.gov/Archives/edgar/data/1045810/000104581026000073/q2fy27pr.htm}
\bibitem{nvda-exposure} NVIDIA, CFO Commentary filed as Exhibit 99.2 to the Q2 FY2027 results (26 Aug.\ 2026), and Form 10-Q for the period ended 26 July 2026: maximum gross guarantee exposure \$108.5B, of which \$105B is the SB~Energy residual value guarantees and \$3.5B land, power and shell guarantees for AI cloud partners' lease obligations in the event of their default; supply and capacity purchase commitments \$279B (from \$119B the prior quarter, memory-driven); AI cloud revenue-share commitments \$36B; future equity investment commitments \$25B; data-centre leases held for third-party reassignment \$20B; non-marketable equity securities \$47.898B and publicly held equity securities \$47.740B. \url{https://www.sec.gov/Archives/edgar/data/1045810/000104581026000073/q2fy27cfocommentary.htm}
\bibitem{nvda-backstop} NVIDIA's revenue-sharing and credit-support model for AI clouds, introduced July 2026: typically six-year take-or-pay arrangements in which the operator retains rental revenue up to a backstop floor and NVIDIA takes a share of the excess; counterparties reported to include Sharon AI ($\sim$40{,}000 GB300) and Firmus. Earlier and separate: CoreWeave's 8-K of 15 September 2025 obliging NVIDIA to purchase residual unsold capacity through 13 April 2032, capped at \$6.3B; a \$1.5B arrangement with Lambda; a 9.3\% passive stake in Nebius (Schedule 13G, 20 Jul.\ 2026). NVIDIA's own framing: the model gives it ``a recurring, usage-linked earnings stream.'' SemiAnalysis (6 Jul.\ 2026). \url{https://newsletter.semianalysis.com/p/nvidia-gpu-debt-backstop-unleashes}
\bibitem{nvda-pause} Reuters and \emph{The Wall Street Journal} (27 Aug.\ 2026), one day after earnings: NVIDIA paused parts of the July revenue-sharing and credit-support programme, reportedly after employees flagged potential antitrust scrutiny and after friction over the extent to which NVIDIA could dictate which customers its partners rented to. NVIDIA's on-record statement: ``The new business model we introduced in July that opens up compute access to the fast-growing AI ecosystem is still in place and continues to evolve due to high demand.'' \url{https://whbl.com/2026/08/27/nvidia-pauses-revenue-sharing-deals-with-ai-cloud-companies-wsj-reports/}
\bibitem{nvda-defence} NVIDIA Q2 FY2027 earnings call, 26 August 2026, prepared remarks: ``We recognize the scale of this support, and we know some will call this circular financing. We see it differently\ldots\ Independent capital still underwrites every deal on its own merits. \emph{We're not making loans.} In this model, we get paid twice---once on the hardware sale, and again through the share of rental revenue.'' Huang, separately: ``We've invested nearly \$50 billion in the Frontier AI Labs\ldots\ the only regret that I have is that I didn't invest more and sooner''; and, on the Ohio structure, ``Is this circular financing? No, and OpenAI is paying the lease.'' \url{https://seekingalpha.com/article/4940563-nvidia-corporation-nvda-q2-2027-earnings-call-transcript}
\bibitem{nvda-open-models} Huang, same call: ``The world will need both closed models and open models. And both closed models and open models are skyrocketing in use\ldots\ \emph{Nearly all open models run on NVIDIA} \ldots\ our position in open models is very good, because the CUDA ecosystem is literally everywhere.'' Compare Nader Shaikh, Director of Developer Technology, NVIDIA, at the Six Five Summit (27 Aug.\ 2026): ``We build Nemotron because we know that this is going to be a multi-model world, and we see where there's the bottleneck''---arguing that durable value accrues to proprietary data rather than to model weights, which is why NVIDIA releases weights \emph{and} training data. \url{https://www.investing.com/news/transcripts/nvidia-at-the-six-five-summit-open-models-reshape-ai-strategy-93CH-4880247}
\bibitem{nemotron} NVIDIA Nemotron 3 family (15 Dec.\ 2025): Nano 30B-A3B, Super $\sim$100B-A10B and Ultra $\sim$500B-A50B, released with \emph{weights, training datasets and reinforcement-learning libraries}. Huang: ``Open innovation is the foundation of AI progress. With Nemotron, we're transforming advanced AI into an open platform.'' Nemotron 3.5 Lightning (11 Aug.\ 2026), 30B MoE with 3B active, shipped under the OpenMDW-1.1 licence with training data and recipes, on Hugging Face, ModelScope, OpenRouter, SageMaker JumpStart and Google Cloud. The Nemotron Pretraining Collection exceeds ten trillion tokens; NVIDIA also publishes open robotics, protein-structure and driving datasets. Nemotron 3 Ultra was the sixth-largest model by weekly tokens on OpenRouter at the currency date. \url{https://nvidianews.nvidia.com/news/nvidia-debuts-nemotron-3-family-of-open-models}
\bibitem{cuda-tile} NVIDIA open-sourced the CUDA Tile IR and cuTile compiler stack under Apache 2.0 with LLVM exceptions (25 Dec.\ 2025)---a Tile MLIR dialect, Python bindings, bytecode representation and conformance suite---describing it as ``the largest and most comprehensive update to the CUDA platform since it was invented two decades ago.'' TensorRT-LLM is Apache 2.0. At Hot Chips 2026 NVIDIA stated it was extending CUDA to cover the Groq-derived LP30 LPU as a supported target; outlets characterise the depth of that integration differently, from ``an LPU extension to CUDA'' to ``a fully-supported target.'' \url{https://github.com/NVIDIA/cuda-tile}
\bibitem{hc2026-lpu} Hot Chips 2026 deck 34, Igor Arsovski and Santosh Raghavan (NVIDIA), ``Think Fast: LPU Accelerator for Heterogeneous Compute'': the Groq-derived LP30---no DRAM, 0.5\,GB of on-die SRAM per chip aggregating to 128\,GB per 256-chip rack at $\sim$40\,PB/s and 315 PFLOPS FP8, a cycle-exact VLIW ISA on a single global clock with no memory barriers and no kernel boundaries, every chip its own router at 350\,ns SRAM to SRAM. NVIDIA acquired Groq's assets and key staff for a reported \$20B, completed late December 2025. The commercially significant slide is the disaggregation set: prefill on Rubin and decode on LP30, attention on Rubin and feed-forward on LP30, Rubin verifying while LP30 drafts---claimed $\sim3\times$ to $\sim5\times$ interactivity over Rubin alone. NVIDIA is publicly designing a system in which its own GPU is not the right chip for part of the inference pipeline. Groq 3 LPX entered full production 24 August 2026. \url{https://www.servethehome.com/nvidias-groq-3-lpu-accelerators-for-heterogeneous-ai-compute-at-hot-chips-2026/}
\bibitem{nvda-morgan-stanley} Morgan Stanley initiation of credit coverage on NVIDIA (24 Aug.\ 2026), rating Neutral: total credit exposure of roughly \$200B by the end of 2028, of which about \$170B off balance sheet---approximately \$38B of minimum rating adjustments, \$68B of residual-value support tied to the \$500B framework and \$64B of revenue-sharing and credit support for non-hyperscalers---against projected net cash of about \$520B in fiscal 2029. The note coins ``balance-sheet-as-a-service.'' Component breakdown reported secondhand; the original note was not obtained.
\bibitem{nvda-critique} Dave Vellante, ``Nvidia's Jensen Huang just made the AI buildout too big to fail,'' SiliconANGLE (15 Aug.\ 2026, updated 21 Aug.): ``Independent capital can fund more factories. But independent capital is not independent demand''---the requirement for credit enhancement is itself evidence that markets do not yet accept compute as enduring collateral. Three simultaneous warning signals nominated: GPU rental prices falling, utilisation weakening, financing terms tightening. Concludes the buildout is not too big to fail but ``too interconnected to fail quietly.'' For the opposing structural reading see NVIDIA's own rebuttal of vendor-financing comparisons (Nov.\ 2025): ``unlike Lucent, NVIDIA does not rely on vendor financing arrangements to grow revenue.'' \url{https://siliconangle.com/2026/08/15/nvidias-jensen-huang-just-make-ai-buildout-big-fail/}
\bibitem{amd-q2-2026} AMD, second quarter 2026 results (4 Aug.\ 2026): revenue \$11.5B, up 50\%; data-centre revenue \$6.7B, up 107\% year on year and 58\% of the total; GAAP gross margin 54\%; third-quarter guidance approximately \$13B. Form 10-Q for the period ended 27 June 2026 discloses \$4.5B of leases not yet commenced and \textbf{\$4.1B in data-centre lease guarantees}, plus data-centre leases with \$9.5B of aggregate future payments beginning in 2027--2028. \url{https://ir.amd.com/news-events/press-releases/detail/1295/amd-reports-second-quarter-2026-financial-results}
\bibitem{amd-warrants} AMD's warrant instruments: a warrant for up to 160 million shares at \$0.01 issued to OpenAI (6 Oct.\ 2025), vesting in tranches tied to 6\,GW of purchases \emph{and} escalating AMD share-price thresholds; and an equivalent 160-million-share warrant issued to Meta under an amended master purchase agreement (Form 8-K filed 24 Feb.\ 2026), vesting from 1\,GW to 6\,GW with price thresholds escalating to \$600 per share. Separately, up to \$5B of equity and 2\,GW of MI450 for Anthropic (22 Jul.\ 2026). As of the August 2026 filing, no warrant shares had vested or been issued. The customer is being paid, in the vendor's own equity, to buy the vendor's product. \url{https://www.sec.gov/Archives/edgar/data/2488/000000248826000045/amd-20260223.htm}
\bibitem{amd-open} AMD's openness posture: ROCm~7 positioned explicitly against CUDA; Instella and Instella-MoE released as \emph{fully} open 3B models---weights, data and training code. UALink 1.0 released 2025 with 2.0 due in 2026 adding in-network computing, chaired by AMD's Kurtis Bowman; the promoter group is AMD, Broadcom, Cisco, Google, HPE, Intel, Meta and Microsoft, and NVIDIA is not a member. \url{https://rocm.blogs.amd.com/artificial-intelligence/introducing-instella-3B/README.html}
\bibitem{terafab} SpaceX ``Terafab'' (confirmed by SpaceX and reported by Reuters, 6 Aug.\ 2026): \$16.8B initial investment in a vertically integrated logic and memory manufacturing, packaging and test site in Grimes County, Texas, with a minimum of 3{,}000 jobs; output targeted at Tesla Optimus, Cybercab and SpaceX space-based data centres. Musk announced the project in March 2026 with a stated ambition of up to $50\times$ the annual AI-chip output of major suppliers; a May 2026 regulatory filing showed a proposed \$55B initial and up to \$119B total across phases. Intel entered a pact with Tesla, SpaceX and xAI (Apr.\ 2026) for advanced logic and packaging, which one analyst characterises as ``likely more an Intel fab expansion with Tesla, SpaceX, and xAI as anchor customers.'' Process node and timelines are not disclosed in the August confirmation, and the figures have moved three times. SpaceX absorbed xAI in an all-stock merger at a combined \$1.25T valuation (2 Feb.\ 2026). \url{https://techcrunch.com/2026/08/06/tesla-and-spacex-will-invest-16-8b-to-start-building-terafab-chip-factory-in-texas/}
\bibitem{tesla-ai5} Tesla's dual-foundry arrangement for AI5 and AI6, confirmed by Musk (4 Nov.\ 2025): variants of the same design at TSMC and Samsung. Samsung's Taylor, Texas 2\,nm fab reached AI5 tape-out (reported 13 Jul.\ 2026), with volume production targeted for 2H~2027. Dojo was shut down in August 2025---Musk called it ``an evolutionary dead end''---and restarted as Dojo3 in January 2026 with the stated goal of space-based AI compute. \url{https://www.tomshardware.com/tech-industry/semiconductors/musk-confirms-tesla-ai5-and-ai6-will-be-made-at-samsung-and-tsmc}
\bibitem{glm53} Zhipu AI / Z.ai, GLM-5.3 (announced 13--14 Aug.\ 2026; API from 18 Aug.; open weights released 28 Aug.\ 2026): $\sim$744B total parameters with $\sim$40B active, 78 layers, 256 routed experts with 8 active per token, 1M context. The licence is neither Apache 2.0 nor MIT but a custom ``glm-5.3'' licence whose reported terms require a security review for model-as-a-service operators above \$10B of trailing revenue. The model card documents Ascend NPU deployment paths through vLLM-Ascend, xLLM and SGLang. Company-reported benchmarks lead on vulnerability discovery (CyberGym 84.5 against GPT-5.6 Sol's 83.6) and trail roughly two to one on exploit development (ExploitGym 105/130 against 216/293); Z.ai's claim of 2{,}436 vulnerabilities found across 269 projects is unaudited. \url{https://huggingface.co/zai-org/GLM-5.3}
\bibitem{glm53-flash} Zhipu AI / Z.ai, GLM-5.3-Flash (released 26 Aug.\ 2026; revealed on 27 August to have been the stealth model ``Ox Alpha''): 320B total parameters with 18B active, hybrid sparse and linear attention, the first natively multimodal model of the GLM-5 series, trained on a 30T-token multimodal corpus, 1M context---and released under \textbf{MIT}. Claimed $\sim3\times$ compute and $\sim4.4\times$ KV-cache reduction against full GLM-5.3. List price \$0.15/\$0.50 per million tokens input/output, roughly $7.5\times$ below GLM-5.3. Hugging Face records 346{,}516 downloads in the first month with 62 quantized derivatives; the model is listed on the Dell Enterprise Hub. \url{https://huggingface.co/zai-org/GLM-5.3-Flash}
\bibitem{glm53-aa} Artificial Analysis intelligence index, launch week (v4.1.1, late August 2026; superseded by the v4.3 re-basing of 7 September, on which GLM-5.3 scores 45 and Flash 42): GLM-5.3 at 60, second of 111 evaluated models; \textbf{GLM-5.3-Flash at 57, fourth of 111}, at \$0.09 cost per task against GLM-5.3's \$1.40. Frontier closed models: Claude Opus 5 (max) 63; Claude Fable 5 62; Claude Opus 5 (high), GPT-5.6 Sol (max) and Grok 4.6 (high) 61. The measured gap between the best open-weight Chinese model and the best closed model is three index points; Flash at 57 ties GPT-5.6 Sol (high) and exceeds Gemini 3.7 Flash. Caveat: the Flash score circulated in Z.ai's own materials before appearing on the public board. \url{https://artificialanalysis.ai/models/glm-5-3-flash}
\bibitem{glm53-local} Unsloth quantization measurements for GLM-5.3-Flash: UD-IQ1\_S 93.1\,GB retaining 70.9\% top-1 accuracy; UD-Q2\_K\_XL 108.7\,GB at 78.3\%; \textbf{UD-IQ3\_XXS 120.4\,GB at 81.6\%}; UD-Q4\_K\_XL 199.7\,GB at 92.2\%; BF16 641.6\,GB. Minimum combined RAM and VRAM: 100\,GB at one bit, 128--150\,GB at three bits, 162--210\,GB at four bits. Unsloth states three-bit ``works on 128GB devices like a Mac or NVIDIA DGX Spark.'' Measured decode on four DGX Spark units at native FP8: 28--34 tok/s on short prompts, $\sim$43 tok/s at 50K prompts. By contrast full GLM-5.3 at FP8 requires eight H200- or H20-class accelerators. A frontier-adjacent model has arrived on a single workstation-class machine---not on a laptop. \url{https://unsloth.ai/docs/models/glm-5.3-flash}
\bibitem{glm53-chips} Z.ai's domestic-silicon claim for GLM-5.3-Flash and what it does and does not say. Z.ai's own announcement states the model was ``previously previewed as Ox Alpha, \emph{running} entirely on Chinese AI chips''; founder Jie Tang adds ``Powered by pure Chinese chips.'' Reporting puts the preview cluster at roughly 100{,}000 domestic accelerators serving 62 trillion tokens over the preview week. Three qualifications are load-bearing. First, the claim is about \emph{serving}, not training: Z.ai has made no training-hardware statement for GLM-5.3 or Flash, and implicator.ai's summary is exact---``served is not trained.'' Second, \emph{no chipmaker is named}; Counterpoint's Ivan Lam speculates Ascend plus other domestic vendors, and CNBC reported it could not verify the claim independently. Third, this is \emph{not} a first: Zhipu claimed Ascend-only training for GLM-Image in January 2026 [P, unaudited]; the June 2026 claim for the GLM-5 family is press-derived with no located primary [A/C] (see the entry on that claim). What is new is production-scale inference on a domestic cluster carrying a globally top-ranked workload. Reports of ``100 trillion tokens per day'' are irreconcilable with router data and should not be used. \url{https://www.scmp.com/tech/big-tech/article/3365433/zhipu-ai-shares-jump-viral-ox-alpha-model-revealed-glm-53-flash-chinese-chips}
\bibitem{zhipu-ascend-prior} Prior Zhipu domestic-training claims, for calibration: GLM-Image trained end to end on Huawei Atlas 800T A2 (Kunpeng 920 plus Ascend 910), reported January 2026, with \emph{The Register} noting that no server counts or timings were disclosed; and the GLM-5 family reported in June 2026 as trained on roughly 100{,}000 Ascend 910B with MindSpore and no NVIDIA at any stage. Neither claim has been independently audited, and for the GLM-5 claim no primary Zhipu statement was located. Against them stands the reporting that DeepSeek's R1 training did not complete reliably on Ascend and reverted to NVIDIA---which is why the training-versus-serving distinction carries weight rather than being a quibble. \url{https://www.theregister.com/2026/01/15/zhipu_glm_image_huawei_hardware/}
\bibitem{openai-luna} OpenAI, ``Advancing the price-performance frontier with GPT-5.6'' (30 Jul.\ 2026): GPT-5.6 Luna---the family's cheapest tier, corresponding to the earlier ``nano'' class---cut 80\%, from \$1.00/\$6.00 to \textbf{\$0.20/\$1.20} per million tokens input/output; Terra cut 20\% to \$2/\$12; \emph{Sol held unchanged} at \$5/\$30. OpenAI attributes the cut to efficiency at every layer and states that GPT-5.6 Sol autonomously optimised production kernels, ``reducing serving costs by 20\% and improving token-generation efficiency by 15\%.'' On 21 August OpenAI cut Sol to \$4/\$20---explicitly ``for the next 3 months,'' a promotion rather than a repricing. No OpenAI executive has attributed any 2026 cut to Chinese competition on the record. \url{https://openai.com/index/advancing-the-price-performance-frontier-with-gpt-5-6/}
\bibitem{price-inversion} The consequence of the 30 July cut, and the fact most easily missed: at \$0.20/\$1.20 GPT-5.6 Luna is \emph{cheaper than most Chinese open-weight models on list price}---below GLM-5.3 (\$1.40/\$4.40), Kimi K3 (\$3/\$15), Qwen3.8-Max (\$2/\$6) and, on input, below DeepSeek V4-Pro. On a combined million-in, million-out basis: Luna \$1.40, GLM-5.3 \$5.80, Kimi K3 \$18, Claude Opus 5 \$30. Measured elasticity, 27 July to 14 August: Luna usage rose to $13.8\times$ its pre-cut average and Terra to $5.6\times$, against $1.1\times$ for the undiscounted Sol, with roughly 75\% of the gain taken from competitors rather than cannibalised. The incumbent did not merely match the commodity price; at the bottom of its ladder it went underneath it. \url{https://openrouter.ai/rankings}
\bibitem{deepseek-raise} Caixin Global (14 Aug.\ 2026): DeepSeek launched V4-Pro and \emph{raised} API prices by as much as 1{,}100\% effective 13 August 2026, introducing peak and off-peak tiering---``a break from the prolonged price war among China's AI developers.'' No rationale was stated. Note the sequencing: the deepest Western cut (30 July) preceded this increase, so neither move can be read as a direct answer to the other, and the simple story in which cheap Chinese models drive Western cuts does not survive the price leader reversing direction mid-cascade. \url{https://www.caixinglobal.com/2026-08-14/deepseek-launches-v4-pro-and-raises-api-prices-by-as-much-as-1100-102473919.html}
\bibitem{openrouter-aug26} OpenRouter public rankings, week to 29 August 2026, by tokens routed: Ox~Alpha (GLM-5.3-Flash, then unattributed) 20.7T, up 216\%; DeepSeek V4 Flash 12.3T; Xiaomi MiMo-V2.5 10.0T; \textbf{GPT-5.6 Luna 6.75T, up 31\%}; Tencent Hy3 6.62T; \textbf{NVIDIA Nemotron 3 Ultra (free) 5.62T, up 27\%}; GLM-5.3 Flash 4.62T; \textbf{Gemini 3.7 Flash 4.05T, up 156\%}. The two fastest-growing Western entries are the two that were discounted, and the sixth-largest model on the board is a chip vendor's open-weight release. Router volumes are a small and promotion-distorted slice of total inference and should be read as a directional signal. \url{https://openrouter.ai/rankings}
\bibitem{openai-econ} Reported unit economics, all via \emph{The Information} and not confirmed by the companies: OpenAI first-quarter 2026 revenue \$5.7B against \$3.7B of cash burn, with adjusted gross margin around 33\% against a 46\% target, inference cost rising from $\sim$\$8.4B in 2025 to a projected $\sim$\$14.1B in 2026, and a revenue run rate of \$40B by August 2026; Anthropic at a \$65B run rate at the end of July 2026 with \emph{API gross margin reported above 80\%}. Sources conflict on the direction of OpenAI's margin. The relevant caution for this book's argument: no credible published estimate of loss per token exists, and no analyst asserts that these prices are below cost---``inference loses money'' is a model-mix and pricing claim, not a universal one. \url{https://www.axios.com/2026/08/17/anthropic-revenue-run-rate-ipo-openai}
\bibitem{hc2026-ai-design} Hot Chips 2026 decks 38 (Google) and 39 (OpenAI)---the only two of forty decks claiming AI inside the silicon design flow, both from companies that design accelerators for their own frontier models. Google: an ``AI Designer'' taking specification to RTL on O(100) TPUs for about a week and an ``AI Optimizer'' doing RTL and physical-design optimisation on O(100) CPUs for about a week, yielding a 6\% power and 5.8\% area reduction on the TPU~8t matrix unit (6\% more FLOPs in the same budget) and 5.3\% area on 8i. OpenAI: measured block-level wins against an ``optimized human baseline'' and the qualitative claim that ``major changes landed up until the day of RTL freeze.'' Neither reports full-chip PPA, bug escapes or schedule saving against a baseline, and \emph{nobody claimed AI-driven verification}, which is where most of a chip programme's cost sits. Separately, OpenAI reports an internal agent driving a DeepSeek MLA kernel from 0.31\% to 88.94\% of roofline in about forty hours---bring-up software, historically the thing that decides whether a novel accelerator is usable, produced by an agent in days.
\bibitem{hc2026-china-bench} The benchmark surface of Hot Chips 2026: nine decks report inference economics on Chinese open-weight models. d-Matrix reports GLM~5.2 and Kimi~K3 serving rates; OXMIQ builds its cost model on Kimi-K2 1T; NVIDIA's Rubin deck measures agentic throughput on DeepSeek-v4-PRO and its Spectrum-X deck uses DeepSeek~V3 training step time as the headline multi-tenancy result; Intel plots every published frontier model from Llama~2 70B to Kimi~K2 1T to establish its capacity-versus-bandwidth thesis; SambaNova quotes MiniMax M2.7; OpenAI reports DeepSeek~R1 on GB300. US silicon vendors now validate and publish their inference economics against Chinese weights, because those are the frontier-scale weights that can be obtained and published against. No Chinese chip is named anywhere in the conference, and no Chinese vendor presented.

% ---------------- Buying the open-model supply chain; AI for Science in China ----------------
\bibitem{nvda-poolside} NVIDIA and Poolside, reported by the \emph{Wall Street Journal} from deal documents and announced in Poolside's shareholder letter (20--24 Aug.\ 2026): approximately \$7B in total---\textbf{\$6B for a non-exclusive licence to Poolside's ``Model Factory''} training platform and \$1B of equity at a \$12B pre-money valuation---together with offers to \textbf{109 Poolside engineers and researchers to join NVIDIA and work on Nemotron}, NVIDIA's open-weight family. Poolside remains independent, its founders do not join, and the licence is non-exclusive. The stated purpose is to build a leading US open-weight model against DeepSeek and Moonshot. Earlier and separate: a reported \$500M--\$1B NVIDIA investment inside a \$2B round (Oct.\ 2025) alongside ``Project Horizon,'' a 2\,GW West Texas site with CoreWeave. Sources differ on whether the announcement date is 20 or 24 August. \url{https://www.forbes.com/sites/jonmarkman/2026/08/24/nvidia-pays-poolside-6b-to-license-its-model-factory-and-109-workers/}
\bibitem{nvda-hf} NVIDIA and Hugging Face---\emph{reported, not concluded} [A]: \emph{The Information} reported on 26 August 2026 that NVIDIA was closing on an acquisition of Hugging Face for about \$12.9B, and as of the currency date neither company had confirmed it, with \emph{Business Insider} reporting that talks had not produced a signed agreement. \emph{As of the currency date neither company had confirmed it}, and \emph{Business Insider} reported that talks had not produced a signed agreement and ``could still atomize''; TechCrunch notes that NVIDIA's silence is itself unusual for a company that moves quickly to correct reports it considers inaccurate. Context: NVIDIA participated in Hugging Face's 2023 Series B at a \$4.5B valuation and had a standing model-hosting and DGX Cloud partnership; Hugging Face reportedly \emph{declined} a \$500M NVIDIA investment at a \$7B valuation in 2025, on the ground that it did not want a dominant investor able to sway its decisions. Hugging Face revenue is around \$150M a year. Unlike NVIDIA's licence-plus-acquihire structures, an equity acquisition would require Hart--Scott--Rodino review. Treat as reported, not concluded [A]. \url{https://techcrunch.com/2026/08/26/nvidia-closes-in-on-hugging-face-acquisition/}
\bibitem{nvda-structured} The pattern of licence-plus-acquihire transactions that leave the target nominally independent and, on the reported structure, outside Hart--Scott--Rodino merger review: \textbf{Groq}, $\sim$\$20B (25 Dec.\ 2025)---a non-exclusive inference-technology licence with roughly 90\% of staff including chief executive Jonathan Ross and president Sunny Madra moving to NVIDIA, Groq continuing under a new chief executive; \textbf{Enfabrica}, $\sim$\$900M (2026); \textbf{Poolside}, $\sim$\$7B (Aug.\ 2026)~\cite{nvda-poolside}. Roughly \$27B in nine months. Senators Warren and Blumenthal wrote to NVIDIA on 23 March 2026: ``by licensing its technology and hiring its most important employees, NVIDIA has effectively acquired Groq in all but name,'' and ``we are concerned that this takeover could stifle competition, further entrenching NVIDIA's dominance in the AI chip industry.'' An earlier letter of 4 February 2026 from Warren, Wyden and Blumenthal asked the FTC and DOJ to investigate reverse-acquihire deals at Meta, Google and NVIDIA. No enforcement action had been taken by the currency date. \url{https://www.warren.senate.gov/newsroom/press-releases/warren-blumenthal-question-whether-nvidias-20-billion-groq-deal-is-attempt-to-avoid-antitrust-laws/}
\bibitem{nvda-segments} NVIDIA's Q2 FY2027 CFO Commentary (26 Aug.\ 2026) discloses a change in reporting structure that most coverage missed: \textbf{``Gaming'' has ceased to be a reported market platform}, with GeForce revenue folded into a new \textbf{Edge Computing} line (\$7{,}198M, +27\% year on year), and \textbf{Data Center is now split into Hyperscale (\$48{,}710M, +102\%) and AI Clouds, Industrial \& Enterprise (\$40{,}313M, +138\%)}---management stating that the latter ``will represent roughly half of our data center business'' going forward. Networking was described on the call as a record quarter growing 18\% sequentially, but \emph{no networking dollar figure was published}. Shipments of Hopper products to China were under 1\% of data-centre revenue and \emph{no} China data-centre compute revenue is assumed in the outlook. Approximately \$26.0B was returned to shareholders in the quarter with \$99.0B remaining authorised. On the call management guided to roughly \textbf{70\% revenue growth in fiscal 2028}, against a sell-side consensus near 44\%---a forward statement that does not appear in the press release. \url{https://s201.q4cdn.com/141608511/files/doc_financials/2027/Q227/Q2FY27-CFO-Commentary.pdf}
\bibitem{nvda-huang-q2} Jensen Huang, Q2 FY2027 press release (26 Aug.\ 2026): ``AI has reached its inflection point. It's doing useful work. Its tokens are productive and profitable. Now, compute is revenue. And demand is accelerating. This time last year, one lab alone was driving the buildout; today, we have a golden age of new AI labs and startups, multiple frontier labs scaling in parallel, \emph{a thriving open-model ecosystem} and physical AI coming online.'' On the call, on OpenAI: ``OpenAI's existing and planned commitments represent approximately \textbf{12 gigawatts} of NVIDIA compute'' through 2030. On custom accelerators: ``many of these XPUs are inference specific chips for one cloud or one service. NVIDIA is an entire AI factory platform that spans the entire AI lifecycle that you can use in any cloud,'' and, of the frontier laboratories designing their own silicon, ``I have 100\% confidence that our technology will continue to be extraordinary for them\ldots\ They're going to use it.'' \url{https://nvidianews.nvidia.com/news/nvidia-announces-financial-results-for-second-quarter-fiscal-2027}
\bibitem{jalapeno-results} OpenAI, ``Jalape\~no: first results'' (25 Aug.\ 2026): vendor-run comparisons claiming \textbf{1.5--1.9$\times$ more AI work per watt at peak throughput and 1.7--3.6$\times$ lower end-to-end latency} against GB200/GB300 and MI355X, measured on GPT-OSS 120B, DeepSeek R1 670B and Kimi K2.5. Two qualifications are load-bearing. First, deployment: OpenAI states it will ``begin deploying Jalape\~no within OpenAI's compute infrastructure by the end of the year''---\emph{the part was not deployed at the currency date}. Second, and decisively for how the challenge should be read, the same post says: \textbf{``We will continue to widely deploy accelerators from NVIDIA and other partners for both training and inference workloads.''} SemiAnalysis notes the comparison sets HBM4 parts against older HBM3E systems and calls it ``somewhat incomplete and unfair.'' Neither OpenAI's nor Broadcom's announcement names NVIDIA; the anti-NVIDIA framing is analyst- and press-supplied. \url{https://openai.com/index/jalapeno-first-results/}
\bibitem{jalapeno-analysts} Analyst assessments of Jalape\~no: Alexander Harrowell (Omdia), ``this is the biggest competitive threat to NVIDIA,'' noting that roughly half of AI infrastructure spending comes from hyperscalers running or able to run custom-silicon programmes; Adrien Sanchez (Yole), that a hyperscaler-designed chip can now match or beat Blackwell-class GPUs on inference efficiency, threatening NVIDIA's inference margins specifically. Against them, Daniel Newman (Futurum) calls the ``NVIDIA killer'' framing ``exhausting'' and credits Broadcom rather than OpenAI for the engineering. Jensen Huang's own response, on CNBC, 26--27 August 2026: ``Lots of projects get started. Lots of projects get canceled\ldots\ I don't have to take anything personally\ldots\ We have the supply chain and the technology scale to be their largest supplier.'' \url{https://www.ibtimes.com/custom-ai-chips-are-coming-nvidia-openais-jalapeno-could-biggest-warning-yet-3806846}
\bibitem{openai-exec} OpenAI executive departures, and the direction of the talent flow. \textbf{Chris Malone}, OpenAI's Head of Data Centers, departed in the week of 17 August 2026 (reported 25 August) after sixteen months, in what TechCrunch describes as a continuing stream of high-profile exits; OpenAI said it had recently reorganised its infrastructure organisation. Malone had never worked at NVIDIA, and the ``head of data centers'' title was OpenAI's alone: he spent close to five years at Meta as a \emph{distinguished engineer} and nearly twelve years at Google working on data-centre infrastructure and IT hardware research and development, per Data Center Dynamics' account of his career. OpenAI's infrastructure now reports to \textbf{Sachin Katti}, recruited in November 2025 from \textbf{Intel}, where he was Chief Technology and AI Officer, ``to work on building out the compute infrastructure for AGI.'' NVIDIA's data-centre leadership is intact---Ian Buck remains VP of Hyperscale and HPC---and the senior technical flow in 2026 ran \emph{toward} NVIDIA, not away from it: Groq's leadership and roughly 90\% of its staff in December 2025, and 109 Poolside researchers in August 2026~\cite{nvda-structured,nvda-poolside}. \url{https://techcrunch.com/2026/08/25/openai-loses-a-top-data-center-exec-as-stream-of-high-profile-departures-continues/}
\bibitem{hunyuan-hy4} Tencent Hunyuan, \emph{Hy4 preview} (28 Aug.\ 2026): open-sourced, \textbf{770B total / 49B active} mixture-of-experts with a context window above one million tokens, distributed on Hugging Face, Tencent Cloud TokenHub and OpenRouter at \$0.834/\$2.501 per million tokens. Tencent's own evaluation across 163 experts and 203 engineering tasks scores it 2.99 of 4 against GLM-5.3 at 2.92 and Kimi K3 at 2.94---vendor-reported and unaudited. It is the largest open-weight release of the last week of August 2026, and three of the four open-weight releases in that window were Chinese. \url{https://technode.com/2026/08/28/tencent-open-sources-hy4-preview-with-770b-parameters-and-a-1m-token-context/}
\bibitem{deepseek-round} DeepSeek's first external funding round, reported by \emph{The Information} and widely relayed (16--17 June 2026): approximately \$7.4B at a valuation above \$50B, with the state-backed National AI Industry Investment Fund (the ``Big Fund'') as the sole holder of direct equity and voting rights, Tencent and CATL in a non-voting limited partnership under a five-year lock-up, and Liang Wenfeng reported to have invested roughly RMB~20B personally. A late-August report of a \$74B valuation ahead of a 2027 onshore listing conflicts with this and is not used [A]. \url{https://www.forbes.com/sites/anishasircar/2026/06/17/deepseek-just-raised-74-billion-heres-the-catch/}
% ---- China's AI-for-science strategy ----
\bibitem{ai-plus-science} \emph{Opinions of the State Council on Deepening the Implementation of the ``AI+'' Initiative} (Guofa [2025] No.~11), signed 21 August 2025 and released 26 August 2025. The document sets out six key actions, and \textbf{``AI+ science and technology'' is the first of them}---ahead of industry, consumption, welfare, governance and global cooperation. Its three science sub-clauses direct the acceleration of ``AI-driven new research paradigms'' and of ``from 0 to 1'' discovery; the construction of \emph{scientific large models} (kexue da moxing) and AI-upgraded research platforms; and---the analytically important one---the integrated development of ``AI-driven technology R\&D, engineering realisation and product landing,'' which chains discovery explicitly to engineering and to product. The document's numeric targets are \emph{diffusion} targets (over 70\% penetration of new-generation intelligent terminals and agents by 2027, over 90\% by 2030); \textbf{it sets no numeric target for science}. \url{https://www.gov.cn/zhengce/content/202508/content_7037861.htm}
\bibitem{most-ai4s} Ministry of Science and Technology and the National Natural Science Foundation of China, launch of the ``AI for Science'' special deployment (the AI-for-science special deployment), 27 March 2023. It names four basic disciplines---mathematics, physics, chemistry, astronomy---and four applied priorities---drug development, gene research, agricultural breeding, and new materials; commits to domain-specific AI4S platforms, to the national next-generation AI public computing platform, to heterogeneous integration of high-performance and intelligent computing centres, and to scientific data sharing. It is a ministerial deployment of work rather than a numbered regulation, and Western sources that describe it as a ``plan'' overstate its form. The traceable competitive funding instrument that followed is the National Key R\&D Programme's ``Disruptive Technology Innovation'' special programme, whose 2025 AI-for-science call (13 May 2025) names five directions: frontier science, life and health, new materials, engineering design, and other AI-driven research-system technologies. \url{https://www.gov.cn/xinwen/2023-03/27/content_5748495.htm}
\bibitem{fyp15-ai4s} The Party's \emph{Recommendations} for the 15th Five-Year Plan (adopted at the Fourth Plenum, released 28 October 2025) contain, in section 14, the operative instruction \textbf{``use artificial intelligence to lead a transformation of the scientific research paradigm''}---use artificial intelligence to lead a transformation of the scientific research paradigm---and an official interpretation built around that phrase carries commentary from CAS academician E~Weinan. A precision point that cuts against the enthusiastic reading: \emph{the phrase does not appear in the Outline (gangyao) adopted by the National People's Congress in March 2026}. In the Outline, AI appears under compute, data and diffusion (Chapter~12 on compute, algorithm and data supply, which directs the study and construction of ``ultra-large-scale intelligent computing clusters'' and government purchase and leasing of compute; Chapter~13 on digital empowerment), while basic research appears in Chapter~8 \emph{without} AI. Anyone writing that the 15th Five-Year Plan makes AI-for-science a central pillar of the state plan is reading the Party recommendations rather than the plan. \url{https://www.gov.cn/zhengce/202511/content_7049042.htm}
\bibitem{ec-ai4s-biblio} M.~Bianchini, V.~Di Girolamo, J.~Ravet and D.~Arranz, \emph{Artificial Intelligence in Science---Promises or Perils for Creativity?}, European Commission, 7 April 2025 (ISBN 978-92-68-25300-7): a bibliometric study of roughly three million papers across eighty fields, 2000--2022, finding that \textbf{China overtook both the United States and the European Union in AI-assisted scientific output from 2016}, and produced more than 25,000 AI-assisted papers in 2022 against roughly 15,000 for the EU and 12,000 for the US---with advantages reported on novelty and citation impact as well as volume. The finding is the European Commission's own, which is why it is used here rather than a Chinese source. \url{https://research-and-innovation.ec.europa.eu/knowledge-publications-tools-and-data/publications/all-publications/artificial-intelligence-science-promises-or-perils-creativity_en}
\bibitem{nsfc-ai4s} R.~Arcesati, MERICS (4 March 2026): roughly half of the project categories in the National Natural Science Foundation of China's current call foreground AI for science, with individual project budgets of EUR 240{,}000--600{,}000 and topics including generative AI for fusion reactor design and biological component design for extreme-environment biomanufacturing [A]. \url{https://merics.org/en/comment/science-funding-call-provides-insight-chinas-ai-priorities}
\bibitem{intern-s1-2025} Shanghai AI Laboratory, \emph{Intern-S1} (released 26 July 2025; arXiv:2508.15763): 241B total parameters---a 235B Qwen3 mixture-of-experts language backbone with a 6B InternViT vision encoder---pre-trained on 5T tokens of which \textbf{2.5T are scientific-domain}, released under \textbf{Apache~2.0} on Hugging Face, ModelScope and GitHub. Self-reported ChemBench 83.4, MatBench 75.0, MSEarthMCQ 65.7, MicroVQA 63.9; comparative claims against Western frontier models are vendor-reported and no independent replication was located. The successor platform, Shusheng Duanyan (Intern Duanyan), announced 17 July 2026, is a five-layer scientific discovery stack with more than two hundred cross-disciplinary agents across six domains and a dry-lab/wet-lab closed loop---\emph{and no open-source release was stated for it}. \url{https://huggingface.co/internlm/Intern-S1}
\bibitem{protenix} ByteDance AI4Science, \emph{Protenix} (Apache~2.0; v1.0.0 released 5 February 2026 at 368M parameters, v2 on 8 April 2026 at 464M): described as the first fully open-source model to outperform AlphaFold3 across diverse benchmark sets at matched training-data cutoff, model scale and inference budget, evaluated with the accompanying PXMeter toolkit over more than six thousand complexes against AlphaFold3, Boltz-1 and Chai-1. Code, weights, training and inference pipelines and a web server are all released. The significance for this book is not the benchmark but the licence: the open-weight ratchet of Chapter~\ref{ch:model} has reached structural biology, and it reached it from a Chinese laboratory. \url{https://github.com/bytedance/Protenix}
\bibitem{deepmodeling} The DeepModeling community and DP Technology (Shenshi Keji, founded 2018; chief scientist Zhang Linfeng): DeePMD-kit, the deep-potential molecular dynamics package (\emph{J.\ Chem.\ Phys.} 159, 054801, 2023); \textbf{DPA-2}, ``a large atomic model as a multi-task learner'' (\emph{npj Computational Materials} 10, 293, 19 December 2024), covering 73 elements across 18 pre-training datasets spanning alloys, semiconductors, battery materials, drug molecules, ferroelectrics, solid-state electrolytes and catalysis, with density-functional labels from VASP, Gaussian \emph{and} the domestic ABACUS code; the ABACUS electronic-structure package; and the \textbf{OpenLAM} initiative (launched 1 January 2024, ``Conquer the Periodic Table!''). The community reports more than ten thousand members across more than twenty countries. DP Technology raised a Series~C above RMB~800M (about US\$114M) in December 2025 in which three of six named investors are Beijing municipal state funds, and reports more than three million users across over a thousand universities and research organisations plus corporate clients including PetroChina, CATL and BYD [A]. Code, weights and datasets are open; on the evidence located this is among the most genuinely open scientific-AI ecosystems anywhere, Western or Chinese. \url{https://www.nature.com/articles/s41524-024-01493-2}
\bibitem{cas-scienceone} Chinese Academy of Sciences, \emph{ScienceOne} , unveiled at the World AI Conference, 26--28 July 2025, built by twelve CAS institutes including Automation, the Computer Network Information Centre, the National Science Library, the Academy of Mathematics and Systems Science and the Institute of High Energy Physics. It is built on \emph{Chinese open-source foundation models} and integrates AlphaFold and MatterGen as tools; it powers a digital-cell platform, particle-collider simulation and telescope coordination. Performance claims are self-reported; parameters, training compute and open-source status were not disclosed. \url{https://english.cas.cn/newsroom/cas_media/202507/t20250728_1048587.shtml}
\bibitem{ai4s-institutions} The institutional layer of Chinese AI for science: the \textbf{AI for Science Institute, Beijing} (Beijing Kexue Zhineng Yanjiuyuan, founded September 2021 under CAS academician E~Weinan), which runs the DeepModeling community and has shipped the DPA potential family, the ABACUS density-functional code, the DeepFlame combustion platform and the AIS-Square model-sharing platform; the \textbf{Shanghai AI Laboratory}, the most productive shipper of scientific models; the \textbf{Shanghai Academy of AI for Science} (2023, with Fudan), which produces the FuXi weather and climate family; \textbf{BAAI}'s 2026 ``Wujie'' series (Brain$\mu$ for neuroscience, OpenComplex~2.5 for biomolecules, Physis for physical world modelling); \textbf{Pengcheng Laboratory}; and \textbf{Zhejiang Lab}, whose Three-Body Computing Constellation launched its first twelve compute-carrying satellites in May 2025 against a stated target of roughly a thousand by around 2030. \url{https://www.ncsti.gov.cn/kjdt/ztbd/szn/snhc/202501/t20250103_191555.html}
\bibitem{pengcheng-cloudbrain} Pengcheng Cloudbrain (Pengcheng Yunnao), the first fully domestic large-scale intelligent computing system, built by Pengcheng Laboratory with Huawei, Peking University and Tsinghua, awarded the \textbf{National Science and Technology Progress Award, First Prize} in the 2025 cycle (announced 8 July 2026). It is reported to support training for more than half of China's domestic foundation models and to have guided the construction of twenty-five intelligent computing centres. The award is the clearest available signal of how the Chinese state ranks \emph{domestic scientific computing infrastructure} against the models that run on it: the prize went to the machine. \url{https://www.stdaily.com/web/gdxw/2026-07/08/content_544307.html}
\bibitem{fuxi-weather} The FuXi (Fuxi) end-to-end weather and climate system (Fudan University and the Shanghai Academy of AI for Science, with the National Climate Center on the Fengshun sub-seasonal model), selected among China's 2025 Top Ten Meteorological Science and Technology Advances (announced 21 March 2026): China's first purely data-driven end-to-end forecasting system, ingesting satellite remote sensing directly, operationally deployed across Chinese meteorological departments and integrated into an early-warning platform serving forty-three countries. Developer-reported claims---medium-range ensembles beating ECMWF on 98.1\% of variables by CRPS, sixty-day forecasts in three minutes---are not independently replicated, and the open-weight status of the current system is not stated. By contrast, Huawei's \emph{Pangu-Weather} has been checked by third parties in the peer-reviewed literature (\emph{Remote Sensing} 17(2):191, 2025; \emph{Meteorological Applications}, 2025; \emph{Weather}, 2025, on both Pangu-Weather and FengWu), which makes weather the one Chinese AI-for-science domain with genuine external validation. \url{https://news.fudan.edu.cn/2026/0323/c1247a148495/page.htm}
\bibitem{ai4s-industrial} Evidence that the Chinese AI-for-science programme is aimed at industry rather than at publication. Of the ten joint results Shanghai AI Laboratory announced at the 2026 World AI Conference, three are aerospace design with a state aerospace research institute and a commercial rocket firm, electroplating-materials formulation with Xiamen University, and an analogue chip-design agent with Shanghai Jiao Tong University claiming five- to tenfold sizing-optimisation efficiency and a halving of layout-planning cycles. Three of the ten carry foreign co-authors (Columbia, Cambridge, University College London), so the programme is not hermetically sealed. The ``AI+'' initiative's science pillar makes the industrial intent textual: its second sub-clause chains AI-driven research to ``engineering realisation and product landing''~\cite{ai-plus-science}. \url{https://www.shlab.org.cn/news/5444277}

% ---------------- The toll road, the export offer, the telecom precedent, and September 2026 ----------------
\bibitem{moonshot-cloud-talks} Reuters exclusive (Liam Mo and Fanny Potkin, 26 Aug.\ 2026), three people familiar with the talks: Moonshot AI is negotiating \emph{revenue-sharing} agreements with Microsoft, Amazon and Google under which the three US clouds would host Kimi K3, Moonshot seeking \textbf{up to 30\% of revenue from K3-related services}, terms said to align with those offered to large customers of the open-weight model; discussions ``at an early stage'' with ``no certainty they will result in agreements''; revenue split, data access and token auditing unresolved; Moonshot has already signed similar agreements with smaller cloud platforms. Microsoft, Google and AWS declined comment; Moonshot did not respond. It would be the first such pact between a Chinese laboratory and a major US cloud. The word in every account is \emph{revenue}, not profit, and no source attributes to the talks any motive beyond monetization. \url{https://finance.yahoo.com/technology/ai/articles/exclusive-chinas-moonshot-talks-microsoft-075340033.html}
\bibitem{kimi-k3-license-maas} Moonshot AI, Kimi K3 licence (Hugging Face, \texttt{moonshotai/Kimi-K3/LICENSE}): ``If the Licensee or any of its affiliates operates a Model as a Service business, and the aggregate revenue of the Licensee and its affiliates exceeds 20 million US dollars\ldots\ over any consecutive 12 months, the Licensee must enter into a separate agreement with Moonshot AI before using the Software or its derivative works for any commercial purpose,'' plus a display requirement above 100M monthly active users or \$20M monthly revenue. The clause is new relative to K2. It is the legal lever behind the cloud negotiations: a hyperscaler cannot serve K3 as it serves DeepSeek or Qwen without contracting with the laboratory, which is why K3, unlike K2.5, is not on Bedrock. Chinasoft International disclosed a ``token revenue-sharing and joint innovation agreement'' with Moonshot on 20 July 2026, percentages undisclosed. \url{https://huggingface.co/moonshotai/Kimi-K3/blob/main/LICENSE}
\bibitem{qwen-revshare} Reuters (7 Aug.\ 2026, two people familiar): Alibaba will require large commercial users of the open-weight Qwen3.8-Max to share revenue, the rate under negotiation and the design explicitly modelled on Moonshot's; the Qwen3.8-Max licence's US\$50M model-as-a-service threshold is the instrument. \url{https://finance.yahoo.com/technology/ai/articles/alibaba-plans-revenue-sharing-next-131417995.html}
\bibitem{us-clouds-chinese-weights} What the US hyperscalers already serve, from their own documentation: DeepSeek-R1 on Amazon Bedrock Marketplace from 30 January 2025 and fully managed from 10 March 2025; six open-weight models added fully managed on 10 February 2026---DeepSeek V3.2, MiniMax M2.1, GLM~4.7, GLM~4.7 Flash, Kimi K2.5, Qwen3 Coder Next; Google Vertex model-as-a-service for Kimi K2 Thinking (13 Nov.\ 2025), MiniMax M2, DeepSeek-V3.2, GLM~4.7 and GLM~5 (10 Feb.\ 2026); DeepSeek in Microsoft Foundry's ``sold by Azure'' category. AWS's third-party model terms state that DeepSeek and Qwen serverless models ``are sold by AWS'' and that customer data is not shared with model providers. No cloud discloses any payment to a Chinese laboratory for serving permissively licensed weights, and Bedrock lists DeepSeek V3.2 at roughly two to five times the laboratory's own price. Kimi K3 is reachable in Foundry only through Fireworks AI as a partner model, at \$3.30/\$16.50. \url{https://aws.amazon.com/legal/bedrock/third-party-models/}
\bibitem{webster-regulability} G.~Webster and Y.~Xu, ``Kimi's tightrope: openness and revenue,'' \emph{Transpacifica} (27 Jul.\ 2026)---``with revenue comes regulability'': downloading open weights left it ambiguous whether any transaction with a Chinese laboratory occurred, but ``if a US company needs a contract with Moonshot to provide the inference tokens\ldots\ the picture looks different''; under a supply-chain rule ``it's harder to argue there's no transacting going on,'' and if Moonshot were sanctioned ``the contract to provide inference is clearly an instance of doing business.'' The only serious analysis located runs against the reading that hosting on US clouds makes a US restriction harder; no source was found that supports it. Separately: Axios (20 Jul.\ 2026) reported a draft executive order that would require US companies hosting Chinese models to guarantee security and accept liability, not issued at the currency date. \url{https://herecomes.transpacifica.net/p/kimis-tightrope-openness-and-revenue}
\bibitem{ndaa-1532} FY2026 National Defense Authorization Act \S1532 (enacted 18 Dec.\ 2025): the Department of Defense must exclude ``covered AI''---defined by reference to DeepSeek and High-Flyer, entities on the Consolidated Screening List or the 1260H list, and companies domiciled in covered nations with unmitigated foreign control---and contractors may not use it in DoD work, with case-by-case waivers. The only enacted federal restriction on Chinese models at the currency date; it does not distinguish self-hosted from cloud-served. The \emph{No Adversarial AI Act} (S.2177/H.R.4142, introduced 25 Jun.\ 2025) remained in committee. \url{https://www.congress.gov/bill/119th-congress/senate-bill/2177/text}
\bibitem{mofcom-model-curbs} Reuters (7 Jul.\ 2026, three sources): China's Ministry of Commerce met Alibaba, ByteDance and Z.ai over the preceding month to discuss options including barring public release of the most advanced models, domestic-only access, curbs on overseas access to unreleased systems, criminalizing model leakage and screening foreign capital; a panel of legal scholars proposed a tiered filing, security-review and domestic-only scheme; ``officials have decided nothing yet,'' any curbs might apply only to future models, and no timeline was given. Bloomberg (26 May 2026) reported passport-surrender and travel restrictions extended from DeepSeek to Moonshot, StepFun, ByteDance and others. Note that Zhipu's two-week hold on the GLM-5.3 weights in August was, on every account located, the company's own safety review and not a government one. \url{https://time.com/article/2026/07/07/china-ai-models-alibaba-bytedance/}
\bibitem{cisa-distillation} NSA, FBI and CISA, joint cybersecurity advisory AA26-251a (9 Sep.\ 2026), naming DeepSeek, Moonshot, Alibaba, MiniMax, StepFun and Z.AI as conducting distillation campaigns against US frontier models, with traffic routed through cloud providers and third-party aggregators. The advisory concerns distillation, not hosting. \url{https://www.cisa.gov/news-events/cybersecurity-advisories/aa26-251a}
\bibitem{deepseek-ascend-order} Bloomberg (4 Sep.\ 2026, anonymous sources), relayed by TechNode and Mobile World Live: DeepSeek plans to order at least \textbf{160{,}000 Huawei Ascend 950DT} accelerators for a gigawatt-scale inference site at Ulanqab, Inner Mongolia, with delivery paced by Huawei's and the domestic HBM supply chain's capacity and expected to take up to a year; training is reported to remain on NVIDIA hardware. The implied cost of roughly \$2.6B is a derived estimate. The order matters to this book less for what DeepSeek gets than for what it shows about supply: a single domestic customer absorbs a year of Ascend output. \url{https://technode.com/2026/09/07/deepseek-huawei-ascend-950dt-chips-inner-mongolia/}
\bibitem{remote-access-rule} \emph{The Information} (28 Aug.\ 2026), relayed by Tom's Hardware: the Commerce Department is drafting a rule to block Chinese firms from renting NVIDIA compute in third countries, reportedly triggered by Moonshot training Kimi K3 on servers in Thailand; the draft may reach industry in September and its legal authority is contested. Not published at the currency date [A]. \url{https://www.tomshardware.com/tech-industry/policy/new-us-export-controls-reportedly-target-chinese-access-to-remote-ai-servers-trump-admins-cut-down-ai-diffusion-rule-could-be-shared-with-industry-as-soon-as-september}
\bibitem{moonshot-ipo} Reuters and HK01 (3 Sep.\ 2026): Moonshot AI has made a confidential A1 filing for a Hong Kong listing, reportedly to raise about \$3B at a pre-money valuation near \$50B; the company declined comment [A]. Reported annualized revenue rose from \$100M in March to \$300M in June 2026 (\emph{The Information}, unverified), with no geographic split disclosed. \url{https://technode.com/2026/09/03/moonshot-ai-reportedly-submits-confidential-hong-kong-ipo-filing/}
\bibitem{aa-v43} Artificial Analysis, Intelligence Index v4.2 (4 Sep.\ 2026) and \textbf{v4.3 (7 Sep.\ 2026)}: two re-basings in one week (Terminal-Bench v4.0 replacing v2.1, AutomationBench-AA replacing $\tau^3$-Banking, GPQA Diamond dropped, private-set weight 45\%). On v4.3: Claude Fable 5.1 (max) 53, GPT-6 Astra (max) 53, Claude Opus 5 (max) 51, Muse Spark 1.3 48, GPT-5.6 Sol 47, GLM-5.3 (max) 45, \textbf{GLM-5.3-Flash 42}, Qwen3.8-2.4T-A95B 40, GPT-5.6 Luna 38, DeepSeek V4 Pro 36, DeepSeek V4 Flash 35. The launch-week score of 57 for GLM-5.3-Flash (and 60 for GLM-5.3, 63 for Opus~5) was on v4.1.1. Any index number in this book carries its version; cross-version comparisons are meaningless. \url{https://artificialanalysis.ai/articles/artificial-analysis-intelligence-index-v4-3}
\bibitem{glm53-flash-agentic} Artificial Analysis independent sub-scores for GLM-5.3-Flash on v4.3: Terminal-Bench v4.0 \textbf{33\%} against 12\% for GPT-5.6 Luna and 12\% for DeepSeek V4 Flash (GLM-5.3: 42\%); AutomationBench-AA 60\% (Luna 50\%); but GDP.pdf 15\% (Luna 24\%), CritPt 15\% (Luna 21\%), output 61--71 tokens/s against Luna's 120, and 34.5\,s to the first answer token; cost per task \$0.25 against Luna \$0.18 and GLM-5.3 \$2.01. Vendor-reported figures on the model card (Terminal-Bench 2.1 84.3, DeepSWE v1.1 63.4) are not independently replicated, and no SWE-bench Verified, SWE-bench Pro, LiveCodeBench or Aider result exists from any source. The model is, on independent evidence, the cheapest point on the price--agentic-coding frontier and not a general replacement for the flagship tier. \url{https://artificialanalysis.ai/models/comparisons/glm-5-3-flash-vs-gpt-5-6-luna}
\bibitem{openrouter-sep26} OpenRouter rankings, data through 8 September 2026, by tokens: Tencent Hy4 preview 18.7T; GPT-5.6 Luna 14.5T; DeepSeek V4 Flash 12.4T; \textbf{GLM-5.3-Flash 12.2T}; DeepSeek V4 Flash 0423 5.0T; MiniMax M3 (free) 4.3T; Hy3 3.7T; Nemotron~3 Ultra (free) 3.7T; MiMo-V2.5 3.4T; GLM-5.3 3.4T. Of the ten, the active-parameter counts where disclosed are 49B, 13B, 18B and below---the volume tier is a sparse-mixture tier. Tom's Hardware (4 Sep.\ 2026), citing an OpenRouter engineer: platform token volume up more than 25$\times$ in a year and doubled in the last month; Claude Fable~5 in ``the low billions'' of monthly output tokens against roughly 12T each for Luna and GLM-5.3-Flash. Router volumes remain a small, promotion-distorted slice of global inference. \url{https://openrouter.ai/rankings}
\bibitem{vercel-aug26} Vercel AI Gateway: open-weight models at 36\% of July 2026 tokens on 8.6\% of spend, a single-day reading of about 62\% of tokens on 22 August 2026, self-hosting break-even ``past 50 million tokens per day,'' 81\% of tokens on models absent six months earlier; the top four frontier laboratories still take 95\% of spend, Anthropic 61--65\% of spend at roughly 4.4$\times$ the average token price. The barbell---volume in the open sparse tier, spend in the closed frontier---is the single most important fact about the model market at the currency date. \url{https://vercel.com/i/open-weight-models}
\bibitem{slm-belcak} P.~Belcak et al.\ (NVIDIA Research), ``Small Language Models are the Future of Agentic AI,'' arXiv:2506.02153 (2 Jun.\ 2025): the argument that most agentic sub-tasks are repetitive, narrow and non-conversational, and are therefore better served by small models with a large model held in reserve---together with NVIDIA's own subsequent agentic flagship, Nemotron~3 Super at 120B, which is not small. Against it, the ``Hierarchy of Agentic Capabilities'' study (arXiv:2601.09032, Jan.\ 2026): on 150 realistic workplace tasks ``weaker models struggle with fundamental tool use and planning,'' and even the best models fail about 40\%. Morgan Stanley (3 Aug.\ 2026) on the rent consequence: open weights ``may really kill\ldots\ the idea that the foundation model layer can sustain high rents''---not compute demand. \url{https://arxiv.org/abs/2506.02153}
\bibitem{glm5-ascend-unsourced} On the widely repeated claim that the GLM-5 family was trained on roughly 100{,}000 Huawei Ascend 910B with no NVIDIA hardware at any stage (June 2026): no primary Zhipu document stating it was located. It circulates through TechTimes (21 Jun.), Tom's Hardware (30 Jun.) and Decrypt (18 Jun.) without a cited source; the GLM-5 repository carries no hardware statement; and Huawei's own Ascend announcement of 12 February 2026 describes ``0-day'' \emph{inference} support and post-release training guides, not a training claim. The one Zhipu domestic-training claim with a primary is GLM-Image (January 2026), with no accelerator count. This book therefore grades the GLM-5 training claim as press-derived [A/C] and does not stack later claims on it. \url{https://www.hiascend.com/activities/dynamic-news/648}
\bibitem{waico} World Artificial Intelligence Cooperation Organization, founding agreement signed in Shanghai on 16 July 2026 (Wang Yi for China; the UN Secretary-General attended): an intergovernmental organization headquartered in Shanghai with a Council of members and a Secretariat; civilian domain only; funded by member contributions scaled to Universal Postal Union classes; enters into force on three ratifications; one-year withdrawal notice. Twenty-nine signatories at founding, reported to include Russia, Brazil, South Africa, Indonesia, Kazakhstan, Pakistan, Malaysia, Serbia, Kenya and most of ASEAN; no official roster was retrievable, and India, the Gulf states other than Oman, Turkey and Egypt are absent. Xi Jinping's WAIC keynote of 17 July pledged 5{,}000 training places for developing-country personnel over five years, regional ``application cooperation centres'' and a meteorological warning system for thirty countries---and \emph{no} compute, financing, datacentre or open-weights commitment. An academic mapping (arXiv:2606.23860) scores WAICO lowest on formalization of fifteen comparable bodies: ``a name, a host city, and a statement of purpose.'' The earlier \emph{Global AI Governance Action Plan} (26 Jul.\ 2025) and \emph{AI+ International Cooperation Initiative} (23 Sep.\ 2025) likewise contain no infrastructure or financing clause. \url{https://english.www.gov.cn/news/202607/17/content_WS6a59a226c6d00ca5f9a0c432.html}
\bibitem{huawei-egypt} Bloomberg (26 Aug.\ 2026, anonymous sources), relayed by The Next Web, Al Jazeera (4 Sep.) and Al-Monitor: Huawei has bid to build Egyptian government AI datacentres---about 1{,}408 Ascend 950-series accelerators for a training cloud and roughly 600 for two inference clusters, a twelve-month build, reportedly with iFlytek---which would be the first confirmed Ascend export anywhere; the US State Department contacted NVIDIA, AMD and Microsoft to assemble a counter-offer and warned that Ascend use may breach US rules, Egypt having needed a US licence for advanced AI chips since 2023. No price, no financing terms, nothing signed; Xi Jinping's state visit to Cairo (31 Aug.--2 Sep.) produced no announced AI agreement. \url{https://www.aljazeera.com/news/2026/9/4/as-the-china-us-ai-race-reaches-cairo-egypt-faces-a-strategic-decision}
\bibitem{huawei-malaysia} Bloomberg via \emph{The Edge} (7 Sep.\ 2026) and SCMP (8 Sep.): Telekom Malaysia is evaluating Huawei Ascend 910C as the backbone of a RM2B (about US\$494M) sovereign-AI project to keep government, military and intelligence data in-country; quantity unclear; US officials warned in 2025 that 910C use could violate US export rules; Malaysia accepted export-control commitments in its 2025 reciprocal trade agreement with the United States; unsigned. In May 2025 a deputy minister announced and within hours retracted a 3{,}000-chip Ascend deployment. \url{https://www.scmp.com/week-asia/politics/article/3366792/will-malaysias-potential-huawei-ai-deal-risk-clash-over-us-trade-pact}
\bibitem{brazil-ai-supercomputer} Brazil, 21 August 2026: an R\$2.3B (about US\$444M) AI supercomputing programme of which R\$1.3B (about US\$251M) goes to Huawei and iFlytek for a Rio de Janeiro machine on Ascend and a Portuguese-language model, from July 2027, with a further R\$1B tender for a Rio Grande do Norte machine expected to go to NVIDIA; financed by Brazil's own FNDCT; 3{,}000 professionals to be trained over five years. ``The strategy is not to depend on a single company, technology or country.'' The only contracted sovereign-AI deployment on Chinese silicon outside China located at the currency date---and it is hedged, half-NVIDIA and Brazil-financed. \url{https://www.aljazeera.com/economy/2026/8/21/brazil-launches-ai-supercomputer-push-while-balancing-us-and-chinese-tech}
\bibitem{kazakh-huawei} Kazakhstan--China package of 16 July 2026 (Tokayev in Shanghai, ``\$15B, seventy-plus agreements''): a strategic partnership between the Ministry of AI and Digital Development and Huawei; a Samruk-Kazyna agreement to acquire Huawei equipment; and a Kazakhtelecom--Hengtong agreement on ``basic principles'' for a 1\,GW ``Data Center Valley'' at Ekibastuz, first facility targeted for June 2027; project-level values undisclosed; a mix of binding agreements and memoranda. Pakistan's Punjab government agreed six datacentres and a thousand ``AI master trainers'' with Huawei on 28 August 2026, with no money, chips or contract described. \url{https://www.intellinews.com/kazakhstan-china-sign-15bn-worth-of-deals-as-tokayev-visits-shanghai-455262/}
\bibitem{us-ai-exports} Executive Order 14320, ``Promoting the Export of the American AI Technology Stack'' (23 Jul.\ 2025), and the Commerce Department's American AI Exports Program: full-stack packages (hardware, data, models, security, applications) with federal financing referrals---direct loans, guarantees, equity, political-risk insurance---and an Export--Import Bank ``ExportAI'' framework under consideration from May 2026; 78 consortium proposals received by 30 June 2026 and \emph{no packages designated} at the currency date. Neither the Chinese nor the American export offer had been delivered anywhere by 9 September 2026; the contest in Cairo is between two offers, neither of which yet exists in the field. \url{https://www.trade.gov/press-release/commerce-statement-application-conclusion-next-steps-american-ai-exports-program}
\bibitem{ascend-supply-ceiling} On the supply side of any Ascend export: Beijing's reported requirement that state-funded compute run at least 80\% domestic silicon under a roughly RMB~2T national compute programme (Bloomberg, 9 Jun.\ 2026, relayed) [C]; a SemiAnalysis-derived estimate that 2026 domestic HBM supports only some 250{,}000--300{,}000 Ascend 910C-equivalents [A/C]; DeepSeek's 160{,}000-unit order alone reported to take up to a year~\cite{deepseek-ascend-order}; and BIS guidance of 13 May 2025 that using Ascend ``anywhere in the world'' risks violating US controls, extended on 1 June 2026 to state that licence requirements follow PRC-headquartered parents abroad. Huawei was reported in July 2025 to be pitching low thousands of 910B to the UAE, Saudi Arabia and Thailand with no finalized deals; the Gulf went to US suppliers. No PRC policy authorizing or restricting Ascend export was located. \url{https://www.aljazeera.com/economy/2026/6/1/us-says-ban-on-ai-chip-shipments-applies-to-chinese-firms-outside-china}
\bibitem{qwen-derivatives} Hugging Face, ``The state of open models, summer 2026'': 151{,}448 Qwen derivative models, about 2.6 times Meta's; roughly 39.6 million monthly GGUF downloads for Qwen against 7.5 million for Llama. The diffusion mechanism that actually operates in the non-aligned world is this one---weights that cost the recipient nothing and lock the recipient into nothing---and it is a different instrument from the hardware-plus-service-plus-debt bundle that built Huawei's telecom position. \url{https://huggingface.co/blog/state-of-open-models-summer-2026}
% ---- the telecom precedent ----
\bibitem{huawei-ar-2025} Huawei Investment \& Holding, \emph{Annual Report 2025}: revenue RMB~880.9B (about US\$126B), net profit RMB~68.0B, ICT infrastructure RMB~375.0B (+2.6\%), R\&D RMB~192.3B (21.8\% of revenue); ``more than 170 countries and regions'' of operation. Earlier: 2006 US\$8.5B; 2009 US\$21.8B; 2020 peak RMB~891B before the Honor divestiture and sanctions; 2021 RMB~636B; 2024 RMB~862.1B. \url{https://www.huawei.com/en/annual-report/2025}
\bibitem{delloro-huawei} Dell'Oro Group telecom-equipment share series as reported: Huawei 27.7\% (2018), 28.1\% (1H19), 29\% (2Q19, Nokia 15.7\%, Ericsson 13.1\%), \textbf{31\% (1H20)}; and \textbf{31\% again in 1H25} (Nokia 13\%, Ericsson 12\%, ZTE 10\%)---\emph{21\% outside China and 40\% outside North America}. Dell'Oro (17 Feb.\ 2026): the RAN market stabilized in 2025 after losing about a fifth of its value in 2022--24, with Huawei and Nokia gaining and vendor concentration at a ten-year high. Omdia (Aug.\ 2025): Huawei the leading RAN vendor in three of five regions (Asia-Pacific, Middle East and Africa, Latin America), third in Europe, absent from North America; 2024 sales only 4\% below 2020. \url{https://www.delloro.com/news/ran-market-stabilized-in-2025/}
\bibitem{ahrens-huawei} N.~Ahrens, \emph{China's Competitiveness: Myth, Reality, and Lessons for the United States and Japan---Case Study: Huawei}, CSIS (Feb.\ 2013): China Development Bank credit facilities for Huawei's customers of US\$10B (2004) raised to \textbf{US\$30B (December 2009)}, with Export--Import Bank support of about US\$600M, and roughly US\$10B drawn by 2013; pricing about 30\% below Western competitors in developing markets and about 12\% below Russian rivals; the rural-first entry (more than two hundred staff across Heilongjiang counties against three or four for Ericsson) and the sequence Russia, Thailand, Brazil, South Africa, Ethiopia (2003), Nigeria (US\$200M CDMA, 2005) before Europe; parity with Ericsson by 2011. \url{https://csis-website-prod.s3.amazonaws.com/s3fs-public/legacy_files/files/publication/130215_competitiveness_Huawei_casestudy_Web.pdf}
\bibitem{wsj-huawei-state} C.-W.~Yap, ``State support helped fuel Huawei's global rise,'' \emph{Wall Street Journal} (25 Dec.\ 2019): up to US\$75B of state support---about US\$46B in loans, credit lines and other financing from state lenders, about US\$25B in tax savings 2008--18, plus grants and discounted land---and, per analysts and customers, prices undercutting rivals by some 30\%. Huawei disputed the characterization. Counter-data: a 2017 European study found Huawei's prices only slightly lower than rivals' in Europe, and a 2019 study found them about 25\% \emph{higher} in China. The 2011 EU inquiry put the US\$30B CDB line at roughly ten times Huawei's operating cash flow. \url{https://www.lightreading.com/business-management/no-one-wants-to-talk-about-huawei-s-state-subsidies}
\bibitem{huawei-5g-seps} LexisNexis IPlytics, \emph{5G SEP Report} (Jan.\ 2025): Huawei first in declared 5G patent families and first in 3GPP technical contributions over 2015--24 (Ericsson second, Nokia third, Qualcomm fourth, ZTE fifth); Chinese-headquartered firms above 40\% of declared 5G families. The standards position is the endgame move of the telecom campaign, executed inside a Western-founded body rather than around it. \url{https://www.lexisnexisip.com/wp-content/uploads/2025/01/2025-LexisNexis-5G-SEP-Report.pdf}
\bibitem{huawei-bans} The pushback and what it removed: Australia (Aug.\ 2018); US Entity List (16 May 2019); United Kingdom (14 Jul.\ 2020: no new equipment after 2020, removal by end-2027, up to \pounds2B and up to three years' delay); Sweden (Oct.\ 2020); EU 5G toolbox (Jan.\ 2020) and the Commission's proposed mandatory phase-out of high-risk suppliers (20 Jan.\ 2026). US ``rip and replace'': about US\$5B appropriated, 53 of 126 projects (42\%) complete by June 2026 against a missed May 2026 deadline. Strand Consult: Chinese vendors' share of European 5G sites about 32\% in 2Q22 \emph{and} 4Q24---flat---with Germany at 59\%; the United Kingdom fell from 41\% to about 20\%; of 23 states that legislated, seven imposed restrictions. The bans bit in the Five Eyes, the Nordics and Japan, and the May 2020 US foundry rule---a semiconductor chokepoint, not a network-security finding---is what triggered the British reversal. \url{https://www.lightreading.com/5g/huawei-has-hardly-been-weakened-in-european-5g-data-shows}
\bibitem{dou-huawei} E.~Dou, \emph{House of Huawei: The Secret History of China's Most Powerful Company} (Portfolio, Jan.\ 2025), and K.~Chan's review, ``Huawei the Hydra,'' \emph{Phenomenal World} (1 May 2025), for the industrial-policy reading; and, on Nortel, S.~Cooper, \emph{Global News} (Aug.\ 2020): the 2004 breach traced to Shanghai (1{,}488 documents including optical and 3G/4G research), Huawei's roughly 40\% underbid on the 2008 Telus/BCE contract, and the caution that a 2014 University of Ottawa study found no evidence that the intrusion caused Nortel's 2009 bankruptcy. \url{https://phenomenalworld.org/reviews/huawei-the-hydra/}

\end{thebibliography}